\documentclass{article}
\usepackage[utf8]{inputenc}
\usepackage[spanish]{babel}
\title{Documento Procesado}
\author{Procesador Automático de PDF}
\date{}

\begin{document}
\maketitle

\section*{THE LIBRARY OF THE UNIVERSITY OF CALIFORNIA LOS ANGELES}
\section*{Gift of Mrs. Lawrence C. Lockley}
\section*{MUSH LIBRARY}
\section*{Digitized for Microsoft Corporation by the Internet Archive in 2007. From University of California Libraries.}
May be used for non-commercial, personal, research, or educational purposes, or any fair use. May not be indexed in a commercial service.

\section*{ph d by Microsoft ®}
4

VIOLIIST TONE-PECULIARITIES BY- FREDERICK CASTLE, M. D.

\section*{FIRST EDITION.}
\section*{LOWELL, INDIANA. Alfred H, Miller}
\section*{Copyrighted 1906, By Frederick Castle, M. D.}
\section*{H. H. RAGON\&SON, Printers, Lowell, Indiana. 1906. IV}
\section*{DR. FREDERICK CASTLE.}
oe \} SS Figs

\section*{Univ Calif - Digitized by Microgoft ® ree}
Music Librurj

It is gratifying to acknowledge Major Gilbert Thompson, Washington, D. C., and Mr. Frank Spalding, Township Principal, Griffith, Indiana, as persons giving valuable aid in making this book presentable. FREDERICK CASTLE.

EXPLANATION OF CHART: Because different samples of sounding-board wood must receive different treatment concerning graduation, therefore values for thicknesses are not given; but in lieu of figures, shading is employ- ed as the means of indicating relative quantitive val- ues. The single experiment noted in the text, and, having maximum evenness of tone-power in view, operated to augment altissimo tones in a more marked degree than tones of lower pitch. Because power in altissimo tones is desirable and difficult to secure, therefore, this method for graduation is given record, and, with the hope that future stu- dents of the violin may continue the experiment of shortening the length of sounding-board activity to augment tones of higher pitch. Trial will un- doubtedly determine a better ratio than 2-3 for shortening fiber-activity beneath the lighter strings.

\section*{VI}
CHART, indicating by shading, the relative quantitive values of violin sounding-board graduation for maximum

\section*{evenness of tone-power.}
\section*{ae es = a eet: Digitized by Mi}
\section*{CONTENTS.}
CHART 6 INDEX 8 INTRODUCTION 12 LECTURE I -Violin Tone-Character 17 LECTURE II- Accident No. 1 26 LECTURE III-The Fraudulent Violin Builder 42 LECTURE IV- Violin Sounding-Board Wood 57 LECTURE V-Disintegration of Violin Interior Surfaces 78 LECTURE VI -My Wornout Violin, or Varnish Phenomenon No. 1 92 LECTURE VII- Varnish Phenomenon No. 2- . 113 LECTURE VHI-Uniformity of Violin Tone- Values 130 LECTURE IX- Violin Tone-Modifiers 150' LECTURE X-Violin Tone-Modifiers, con- tinued 160 LECTURE XI -Violin Tone-Modifiers, con- tinued 171 LECTURE XII -Violin Tone-Modifiers, con- tinued 180 LECTURE XIII -Violin Tone-Modifiers, con- cluded 190 LECTURE XIV- Maximum Evenness of Violin Tone-Power 205 LECTURE XV -Maximum Evenness of Violin Tone-Power, concluded 222 LECTURE XVI-Maximum Violin Tone-Power 235 LECTURE XVII-Philosophy Involved in Vio- lin Interior-Surface Conditions 256 LECTURE \textbackslash\{\}VlII-Our Pass-Over 274 APPENDIX 289

\section*{vn}
\section*{INI3KX.}
ACCIDENT NO. I, 26. of, 93. AUDIBLE SOUND, GUSHER, limited to twelve inches the, 285. of the violin sounding- HARMONICS, board, 53. HARMONICOVER- value of, 96. ANGLES, incidence and reflection, TONES, 294. 304. INTRODUCTION XII, BOOKS, method of arriving at unreliable evidence in, conclusions, xiv, eight 35. propositions, xv. BEST SOLO VIOLIN, INANITY, not best orchestra vio- continuing the demand lin, 37. that soloists only appear BACK PLATE, with a Stradivarius, or functions of, 47; as a Guarnerius, 39. tone-producing agent, INVITATION, 107. to my sheep-herder vio- CARE OF VARNISH, lin, 101. example, 93. LONGEVITY CREMONA VARNISH, of the violin, 98; exam- 125. ple, 98. COMBINED TONAL EF- LOSS OF TONE-POWER, FECT, causes of, 299. half dozen violins of METEORICCONDI- uniform tone-pitch, 136. conditions of test, 136. TIONS, 23. COLD TONE, 296. Musical sound, explicit DISSONANTOVER- law, 30; affecting dis- tance traveled by violin TONES, tone, 280. cause of, 30; 295. MAXIMUM EVENNESS FRAUD, 43. GENERAL PURPOSE of violin tone-power, 205; fundamental tones, Violin, 38. 208; reasons leading to GUMS, a new method for sound- hard, inelastic, eifect ing-board graduation,

\section*{vin}
INDEX. 214; demonstration of PLAYERS, sounding-board areas Varying opinions of, 36. augmenting tone of each PHILOSOPHY string, 216; ratio for Involved in violin in- lengths of fiber-activi- terior-surface condit- ty beneath each string, ions, 256; list of princi- 218; concert pitch 200 pals modifying violin years ago, 223; demon- tone, 259; number of stration that errors in tone-qualities absolute- sounding-board gradua- ly at command of the tion cause uneven tone- violin builder, 261; in- power, 225. tensity of violin-tone MAXIMUM TONE-POW- the product of four fac- ER, 235; tors, 262; properties of "big" tone, 236; violin air affecting intensity tone separated by two of violin tone, 264. irreconcilable factors, PENETRATION 237; violin aestheticism of 274. run mad, 238; list of PASS-OVER, oil, factors producing max- 276. imum violin tone-power; POST, physical appearances of problem of, 301. sounding-board wood POST-SETTING, 303. producing maximum RICH TONE, tone-power combined causes of, 96; descrip- with "rich" tone, 246; tion of wood yielding normal and transverse rich tone, 97; illustra- vibration in the sound- tion, 296. ing-board, 249; woody SCIENCE, tone, 255; dispersion of never made a violin, 32. force, the rebounding SCIENTIFIC STATE- ball, 271. MENTS, NOISE, received with caution 32 definition of, 25. SOUNDING-BOARD NODES, 294. WOOD, 57; 0, THOU POST, 276. mite contributed by sci- PITCH, ence, 58; value in scru- Musical sound, 30. tinizing wood of used

\section*{IX}
\section*{INDEX.}
violins, 61; unequal back, 144; thickness of skrinkage, example, 62; sounding-board, 144; example, 63; faults of arching, 147; high a"rch, pine, 66; Michigan pine, example, 148; laws gov- 67; white cedar, 69; col- erning lines, of sound- or changes in pine, 70, wave travel, 151; un- independent action of solved problem in arch- contiguous fibers, 71; ing, 152, the bar, 152; pathetic, 78; preserving wolf caused by mal-po- interior surfaces, 79; sition of bar, 153; wolf example of interior-sur- I caused by graduation, face disintegration, 83. 154; position of bar di- SYMPATHETIC A C- minishing power of the TION, 108; D, 158; the post, 1QO; the law of, 109. bridge, 164; the finger- SWEET OLD VIOLIN, 91 -board, 165; the strings, THE KING, 181; re- enforce block, TWO 20. 175; the trembling vio- MISTAKES, 22. lin, 181; interior sur- TONE-DISTANCE TEST faces, 185; the exits, out of doors, 23; the 187; depth of ribs, 199; listening ear in the bet- the mute, 302;bowhair, ter position, 24. 202. TRADE PROMOTERS UNIFORMITY industry of, 39. of violin tone- values, TWO-DOLLAR MAN, 41. 131. THOU, VIOLIN, 45. VERDICT, TONE-PITCH, 45; Label, Varnish and -"'. eight principals' gov- Price vs Sweet Tone, erning, 46; application 305. rulel, 48; rule 2, 48; VIOLIN TONE-CHAR- rule 3, 49; rule 4, 49; ACTERS, ruleS, 50; rule 6, 50; place of usefulness, 19. rule 7, 51; rule 8, 51. VENEERING WORKS, TONE-MODIFIERS, list, 141; varnish, 142; 53. bending the sounding- VARNISH PHENOME- board, 143; bending the NON NO. I, 104.

INDEX. VARNISH PHENOME- VENTRAL SEGMENTS, NON NO. II, 113. 295. WOLF, VIBRATION, caused by the bar, ex- normal and transverse, ample, 153; caused by 291; velocity compared, sounding-board gradua- 299. tion, example, 154. ERRATA. Page 19, line 8, for "tenora," read tenoro. Page 44, line 2 verse, for "Music," read Music's. Page 48, line 12, for "give," read gives. Page 57, line 24, for "govern's," read governs. Page 96, line 14, for "a basso,'' read a bassa. Page 173, line 1, for "diminish," read increase. Page 216, line 20, for "purfing," read purfling. Page 217, lines 1, 13, for "purfing," read purfling. Page 297, line 3, for "MOVEMENT," read MOVEMENTS.

\section*{XI}
\section*{VIOLIN TONE-PECULIARITIES.}
INTRODUCTION. These lectures, addressed to an imaginary audi- ence of violin students, were originally written for and partly published in the Western Musician, Dixon, Illinois, and, for entertainment of the many readers of this musical journal. Two of the lect- ures now appear in print for the first time. As a familiar style was employed, therefore abstruse technical terms were avoided as far as possible without interfering with clearness and precision. The experiments, results, and conclusions, as thus recorded are not romances of the imagination, as might be inferred at first, but, are conclusions gained by practical experiments, and also by acci- dents occurring in my experience. Thus, when violin patients came to my hospital, I was happy, and, because of my enthusiatic devo- tion to problems in tone-diagnosis, I worked upon them, and over them, until pronouncing them cured, or incurable. Some of those violin-patients were as some human patients, blest with inherently good constitutions to begin with, and, were capable of receiving enhanced tone values from careful adjust- ment of tone-modifying factors, while some of them were so inherently bad from the day they were named "violin," (misnamed,) that only noisy tone was their inheritance; yet, noisy tone made "inter- esting cases" of the latter class because of offering incontestable reasons for inferior tone reasons conclusively demonstrating truth in the statement, "without superior material, without superior vio- lin."

\section*{xn}
\section*{VIOLIN TONE-PECULIARITIES.}
Throughout my period for active work, the fol- lowing questions were always in view: "How does the violin operate to produce musical sound?" 'What ' agents, connected with the violin, oper- ate to modify tone?" "What are the causes for inferior violin tone?" "What are the causes for superior violin tone?" Some of these questions I have solved to my own satisfaction, but, no claim is advanced that such solutions will be acceptable to other students of violin tone-phenomena; nor is the claim advanced that all such tone-problems have received solution. Some of my conclusions are at variance with con- clusions of noted scientific investigators, but, for my own conclusions, infallibility is not claimed. To err is human. To follow error is also human. Thus, I followed a scientific conclusion concerning pro- duction and modification of violin tone requiring experiences of twenty-five years to dispel the delu- sion. Upon this ground, the violin student is warned of danger in following abstract theory under the guise of science. It is my belief that theories, even when based upon oft-repeated prac- tical demonstrations upon various violins, should be presented only as conclusions of an individual attempting solution of a problem wherein capric- ious action of wood has ever been, is now, and for- ever may remain an unknown quantity; and, I pre- sent the thought that such unknown quantity is the reason why science meets defeat in attempting to build a violin to order.

\section*{xm}
\section*{VIOLIN TONE-PECULIARITIES.}
The following problems remain for elucidation: ' 'Inherent, capricious spring-action of wood. ' ' ' 'Varying degrees of sound-wave concentration at the exits as governed by varying degrees of plate-arching." 'The ' phenomenon of raising tone-pitch by en- larging area of exits." The opinion is presented that solutions for the first two of these problems will place violin tone- quality at command of the will. Notwithstanding doubts of solving problems involved in capricious action of wood, yet, the value in such solution re- mains a powerful incentive to continued effort. The desire for violins possessing "rich" tone com- bined with marked intensity of tone is a stimulus surpassing the stimulus of fine gold; and, whoever discovers a method producing such violins at com- mand of desire will become a king in his own right. My method for arriving at conclusions concern- ing the potency in each modifier of violin tone is to run down the causes for noisy tone, sweet tone, powerful tone, hollow tone, thin tone, tone "all inside," tone "all outside," volume of tone, intens- ity of tone, tone pitch, unmusical double-stop tones, powerful open tones with feeble altissimo tones, resultant tones, or harmonics a bassa, consonant overtones, dissonant overtones, the "rich" tone, the "cold" tone, sympathetic tone, evenness of tone-power, and tone character . based upon tone character of the human voice. In this work, the conclusions herein presented follow experiments directed upon both old and new

\section*{xiv}
\section*{VIOLIN TONE-PECULIARITIES.}
violins, and the number of such violins runs into hundreds. From deductions thus obtained, it is my desire to give prominence to the following propositions: 1. Tone-peculiarities, existing in. a given violin, may not exist in any other violin. 2. Writing up tone-peculiarities existing in a given violin as infallible necessities for all violins is misleading. 3. Finding two violins possessing precisely simi- lar tone- values is equally difficult with finding two voices possessing precisely similar tone- values. 4. No violin maker whatever, is, or has been able to give marked tone- value to each and every violin. 5. That the mountain herder of sheep may pro- duce a violin possessing tone-values equal to the best. 6. That th3 skillful mechanic, guided by unerring musical instinct, produces a vastly great- er number of superior violins than the mechanic minus such instinct. 7. That all .violin makers may meet occasional defeat. 8. That, barring accident, the superior violin is a product of superior mechanical skill combined with superior musical sense, all being directed up- on superior material. Than the method herein presented for determin- ing the potency and operation of each factor enter- ing into the production and modification of violin tone, there seems no other method offering equal

\section*{xv}
\section*{VIOLIN TONE-PECULIARITIES.}
value to conclusions. To the weighing of such factors, I have given a lifetime; not in abstract theorizing, but in sitting at the bench while repeating demonstration after demonstration, year after year, decade after de- cade, from youth to old age, determined upon iso- lating, weighing, and knowing the operation of each and all factors underlying violin tone-phenom- ena or die in the attempt. At sixty-three years of age, death came near, and three years later re- mained near, leaving but my right arm sufficiently useful to guide the pen. It is now certain that I shall not reach the goal of my ambition. Under such difficulties, writing is laborious, be- sides, the matter herein is made up wholly from memory, no notes having been made with the view of publication. At the present moment, necessary conservation of strength confines me to a limited daily period for work; hence abandonment of in- tended re-writing of preliminary publication, of which needed corrections are made, and from which some paragraphs are omitted, and to which lectures xvi, and xvii, are added. This publication is presented as my legacy to both the violin student and the American violin maker. That the follow- ing record is reduced to writing and published is matter wholly due to encouragement offered by a modern violin maker; therefore, whatever of en- tertainment, or whatever of other value may be found upon these pages is something not attribut- able alone to courage of, FREDERICK CASTLE. Lowell, Indiana, March 20, 1906.

\section*{xvi}
\section*{VIOLIN TONE-PECULIARITIES.}
LECTURE I. GENTLEMEN OF THE VIOLIN STUDENT CLUB: At this, our first session, I take the opportuity to of- fer you congratulations for the following interest- ing facts. First: The causes for noisy violin tone are discovered. Second: A successful way to pre- serve interior surfaces of the violin from disinte- gration by heat and moisture has been perfected. Third: Areas of the violin sounding-board, re- sponsible for production and augmentation of tone are now located and defined. Fourth: A method for sounding-board graduation, securing maximum evenness of tone-power has received demonstration. Sixth: Principles, governing violin tone-intensity, are brought out into the light. Seventh: Princi- ples governing violin tone-power are reduced to words. Eighth: Quality of sounding-board wood, permitting ' 'rich ' ' violin tone, is described. Ninth : The power of accident to dispel darkness and de- lusion is placed on record. Tenth: Some scientific 'how conclusions concerning ' the violin operates to produce musical sound" have experienced a "jar." Eleventh: The fallacy in the claim that "the best of Cremonas are necessary vehicles for inter- pretation of Haydn, Mozart and Beethoven scores" is made manifest. Twelfth An attempt has been made to right the wrongs heaped : upon the modern violin maker by the "old- violin-trade-promoter. " During our course of study, you may be present- ed with some ideas pertaining to violin tone hereto- f ore" not given expression. Indeed, I promise you but little stale hash in our menu. As it is not

17

\section*{VIOLIN TONE-PECULIARITIES.}
my intention to rob you of the pleasure afforded by anticipation, therefore you will be given but small doses at one time. This plan is adopted to avoid injuring your digestive power, and, to secure your regular attendance. Tone taste varies varies up through every de- gree of musical culture. A tone suiting one may in no wise suit another. No player can play his, or her, best upon an instrument whose tone to him, or her, is disagreeable. Fortunately, the violin affords a variety in tone-quality so infinitely great that every violin player on earth may own one having tone-quality to his taste. One might think it possible to make violins hav- ing a single standard of tone-quality, but the fact is, inherent peculiarity of action in wood stands in the way. There is something approaching an unvarying tone-standard for the whole range of wind instru- ments and instruments of percussion. But unvarying tone-quality suddenly halts in the presence of the violin family. We can imagine the unbounded surprise of one who never heard other than wind musical instruments upon introduction to this violin family. As he picks up violin after violin, viola after viola, cello after cello, he finds no two possessing an identical tone-character. Each violin, each viola, each cello, has a tone-qual- ity peculiar to itself. These peculiarities are so marked that he soon becomes able to name each violin with whose tones he is familiar, although blindfolded, or in a distant apartment, naming

18

\section*{.VIOLIN TONE-PECULIARITIES.}
them as unerringly as he can name different sing- ers with whose tones he is acquainted. He becomes curious to know the reason, or rea- sons, for the infinite tone-peculiarity of that wond- erful musical instrument called 'Violin." By observation, he finds its G and D strings sometimes possess the bass tone-character; at an- other time these strings possess a baritone-tenor* character; in some instances A and E strings pos- sess mezzo-soprano character; in other instances A and.E strings possess soprano tone-character alone. Because of these tone peculiarities he divides violins into four classes, thus: 1. Basso-mezzo-soprano. 2. Basso-soprano. 3. Baritone-mezzo-soprano. . 4. Baritone-soprano. By experiment, he finds that these four classes of tone-character can be given to violins at will; and that they are dependent upon various degrees of .soundingrboard thicknesses, together with such tone-modifiers as size and position of exits, air capacity of the violin, etc. From observation he finds a field of usefulness peculiar to two of these classes. Thus: The vio- lin of basso-mezzo-soprano tone-character is the more agreeable solo instrument, while the baritone soprano tone-character is decidedly the more effec- tive for orchestra uses; that the latter fact is due to high tone-pitch, therefore its tone- waves ride on the topmost wave of harmony parts. Surprising as are these peculiarities, he yet finds

19

\section*{VIOLIN TONE-PECULIARITIES.}
another fact in violin tone still more surprising; that is, some violins possess human tone-quality in a degree far outranking all other musical applian- ces. Thus a new world of expression is opened to him. Here is a musical device capable of talking; tak- ing part in dialogue a la "Arkansas Traveler" an instrument capable of arresting the song of wild birds, causing them, with outstretched necks and wonder-lighted eyes, to look about for that other strange singer pouring forth those enchanting trills; an instrument capable of breaking out into joyous laughter a la the laughter score in Pagan- ini's "Carnival;" an instrument capable of uttering prayer devout in Mozart's "Song Without Words;" an instrument thrilling the air with those trouble- forgetting tones in Schuman's "Traumerei;"anon, awaking human tenderness with sympathetic tones of "Sweet Home;" anon, making human eyes weep with those matchless, touching, despairing, farewell tones, as brave, loving, unfaltering Norma, by a stern Druid father condemned to death by fire, singing in "duetto e scena ultima" as she ascends that blazing funeral pyre God! the agony of it! and the blinding tears.! In all the wide world there's no musical instru- ment approaching the violin. Our violin investi- gator is now become a violin devotee. The en- trancing tones of this human-attuned wonder com- pel him to bow down and worship it as ' 'The King. " Thou, Violin! Thou that smil'st on beggar and king!

20

\section*{VIOLIN TONE-PECULIARITIES.}
Thou thing that laughs weeps prays sings! Thou thing of beauty! Thou joy forever! Nor kings, nor queens, nor potentates whatever Reign with thy absolutism. Thou, Violin! Do you condemn this man for such idolatrous worship? I have given more than fifty years to the search for causes of violin tone-peculiarities; finding some of them, so I think. I claim no superior knowledge of physics, nor superior penetration. I only claim merit for tenacity. The physiognomist might say of me, ' 'you have a square jaw." The phrenologist might say, "you have a remarkable developement in the region of never-let-go". Both might conclude by saying, "you have nothing else worthy of remark." Only , tenacity can hold any man to fifty years work upon any single problem. I worked forty years in trying to make all violins sweet in tone; trying, in other words, to find the cause for "noise" in violin tone. I had concluded that sweet violin tone is an accident, when a real accident occurred revealing the cause of noise in less than ten min- utes. Irony? Much. I 'confess that a solution by accident is better than no'solution at all. When a man, even by help of accident, lives to demonstrate a principle of benefit to humanity, he may then depart, knowing

21

\section*{VIOLIN TONE-PECULIARITIES.}
the world is better for his having lived therein. Is it not a fact that when a man can lift humanity above ear-rending, soul-harrowing, suicide-impelling vio- lin noise, he has sufficient basis for any reasonable claim on earth or in the heavens?. Let millions having "nerves" attest. Enough! Some violins are sweet in tone; some are not. Truthfully, a few violins are thus sweet; many are not. To the listening ear, sweetness is the chief ele- ment of value in violin tone. Strangely, there are a few violin players who place no value upon sweetness of tone, saying, "I'll take care of the sweetness of tone if I can only get tone-power." Never was mistake greater. The best violin players, from Ole Bull down to Mr. "Saw-yer-head-off " could not, nor ever can, conceal "noisy" violin , tone. Admitting a differ- ence in favor of skillful bowing, yet, skillful as one may be with the bow, ninety per cent of an audi- ence will say, "That fellow can't play the violin." By sweet tone I mean tone unaccompanied by sound-waves pitched at inharmonious keys. Again, it is a mistake to suppose the "noisy" tone to travel an equal distance with sweet tone. On this point, the following test for carrying- power gives the loud tone devotee a good opportun- ity for disillusion, and also to part with wealth. It is a test that I have repeatedly made, and made with unvarying results.

22

\section*{VIOLIN TONE-PECULIARITIES.}
As you know, a loud toned, noisy violin, played in a small, bare room makes one long for the quiet of shady woods. From a number of violins, tested in such small room, select the noisiest and the sweetest and take them to an open field, level, and affording at least 1400 lineal feet of unobstructed distance. Select a day when the winds are at rest; a cloudless day is best, because anything approach- ing the nimbus cloud greatly assists the propaga- tion of sound. Select an hour from 10 a. m. to 4 p. m., because in those hours of any given day sound is propagated with greater difficulty. To make the record of value, take along a thermometer, a ba- rometer, and a hygrometer, and record the read- ings of these instruments at hour of test. Thus the carrying power, of a violin so tested, can be guar- anteed to repeat its performance at any time hav- ing similar meteoric conditions. Meteoric condi- tions, as you know, greatly modify distances at which sounds travels. In our summer months, these conditions often cause violin tone to be dis- appointing. Thus, any violin may be the object of unwarranted tone-criticism. In making this tone distance-test, two persons, at least, must assist. One will play a melody upon the G string; the other retires across the field to a distance at which the melody is faintly distinguished. This distance, being measured, is placed to the credit of that string. Thus a record is made for each string of each violin. This test establishes:

23

\section*{VIOLIN TONE-PECULIARITIES.}
(1) Evenness of tone power. (2) Intensity of tone. (Carrying power.) (3) Purity of tone. (Sweetness of tone.) In my experience, the sweeter tone invariably travels the greater distance. I have known the sweeter toned violin to win by 250 feet. In even- ness of tone, I tested one reputable old violin whose distance record\textasciicircum\{\} varies from G string 1000, to E string, 1480 feet. Only two violins have I thus test- ed, securing an equal distance-record for each

string. I assure you this tone distance-test may cause profound admonishment to the participants. My- self, after long experience, dare not risk anything on the result. This test affords ample proof that the listening ear is in the better position for judg- ing tone. At this point I present "noise." It is a familiar thing, truly. It is a thing which ought not be found in the violin. Not always in text books do we find definition of "noise." Noise seems to be a painful subject. Proximity accentuates its painfullness. The existence of "noise" is as the density of population. To noise may be charged the existence of "nerves." This fact may be proven by rolling back a few centuries when there were fewer people on earth, fewer things moving about, and vastly fewer fid-

24

\section*{VIOLIN TONE-PECULIARITIES.}
dlers, and there we find no record of "nerves." Science would have us believe that where there are no ears, there is no sound, no "noise." Do you believe this story? If you do not care to express disbelief you may, at least, call it a paradox. But, to think of a place where there's no noise! Blessed place! That must be a place where angels do not fear to tread. Let it not be invaded by the "noisy" violin, for the most distracting noise, noise "le plus terrible," can come from the violin. What is noise? Should you not find a definition to suit you, read the following: Noise is an aggre- gation of sound-waves pitched at inharmonious keys. The definition itself makes one shiver. I'm proud of it; the shiver, I mean. The shiver is proof that my definition is correct. What is the cause of "noise" in violin tone? This question is filled to the brim with absorbing interest to the whole violin world. The accident, affording a solution to this mo- mentous question, is so provokingly accidental as to rob me of all honor for the solution. I had worked a lifetime for this solution; I had read books, and books; had bought a few books my- self; had borrowed more; I had learned therein some facts requiring twenty-five years to unlearn; I had given up the possibility of a solution, believ- ing sweetness of violin tone to be an accident, when the following accident occurred, thus:

25

\section*{VIOLIN TONE-PECULIARITIES.}
\section*{LECTURE II.}
GENTLEMEN: I was working upon a used violin, a violin of "noisy" tone-faults, a violin belonging to a "business" player. Re-graduation was com- pleted, a new bar ad justed, interior surf acing- work finished and all ready for assembling when the owner sent in a "hurry up" call. To hurry up meant working at night, and to please, my friend, I resolved to do assembling work that same night. Accordingly I returned to my workroom, lighting a 20-candlepower lamp placed before a reflector and seated myself immediately in front of the lamp. Reaching to my right, I picked up the finished sounding-board. As it came between my sight and the lamp, I noticed some dark spots upon the inner surface. Thinking that in some way I had soiled the surface, I laid the sounding-board down, in- tending to remove these spots with a cloth. But, as the sounding-board lay upon the bench, there was no soiled spots to be seen. They had disappeared. The inner surface was clear and bright. But, cer- tainly I saw dark spots. Again I hold the sound- ing-board up to the lamp. There they are, plain enough; several of them; some large, some small. They must be caused by spots on the varnished surface. Critically I examine the varnish. No spots appear thereon. Possibly these spots may be due to opacities in the wood, but I find no opaci- ties; the wood being clear and bright. Again I hold the sounding-board to the lamp. One, two, three, six cloudy areas are there, and located on the up- per tone-producing area of the sounding-board.

26

\section*{VIOLIN TONE-PECULIARITIES.}
Motionless, I gaze at those cloudy areas, thinking, thinking hard for an explanation. Slowly the ex- planation came, slowly of course. I am not brilliant; never posed as a prodigy; am slow, but am a ' 'stay- er," also a smoker. As you know, smoking is a prerequisite of the philosopher. We cannot disas- sociate the pipe and the German, neither can the world match him in solving knotty problems. In imitation of the philosopher, I light my pipe. At once, through clouds of smoke, I see the cause for these cloudy areas. Thus does "smoke" aid phi- losophy. The great Hahnemann established the principle that "like cures like;" or, as Hahnemann states in choice Latin, ' 'Similia similibus curantur. ' ' I believe in Hahnemann as far as "smoke" goes; yet, that great thinker did pierce the thick skin of the "regular" doctors with the sharp-pointed fact that our doses often surpass generosity. And now myself, one of the ' 'regulars, ' ' am going to dem- onstrate that the "infinitesimal" can cause "noise" inside the violin too. I believe in accident. In due course of time I shall present to you two acci- dental occurrences pointing directly at solutions of violin tone-problems hitherto pronounced un- solvable. Therefore when meeting unsolvable questions, I do pray for an accident. Only for accidently carrying this sounding-board with its convex surface towards the lamp, I yet might be searching for the cause of "noise" in violin tone; moreover, might never have found the cause. In the following narration you can plainly see the unquestionable evidence afforded

27

\section*{VIOLIN TONE-PECULIARITIES.}
by accident, for, in the light of theories of musical sound, this evidence stands forth in clear-cut outline. To make my process of reasoning clear, I call up a few familiar facts in natural philosophy: Thus, light rays, falling upon the convex surface of a con- cavo-convex body, (as the violin sounding-board, ) after passing through such body are bent from a direct line at the moment of leaving the concave surface, and thence travel in converging lines to a focus; theiJ\textasciicircum\{\}ore, objects upon the concave surface become magnified; but, reversing the direction of travel, light rays leaving the convex surface, are dispersed; hence, objects upon the convex surface appear diminished in size; therefore, had I carried this sounding-board with its convex surface towards the lamp, I could not have seen those cloudy areas. [The clear varnish upon the C( nvex surface of this sounding-board greatly'* facilitates transmission of light rays; and, in applying this test for perfect graduation work upon sounding- board wood "in the white" it is necessary first to cover the convex surface with some transparent me- dium, as oil, or clear varnish, or better, a mixture of these two available substances.] In doing graduation work upon this sounding- board, I supposed myself to be doing fine work; using calipers in a careful manner; yet, I had fail- ed of doing perfect work, as will now be shown. Had this sounding-board been replaced without dis- covery of those cloudy areas, there would yet have been mixed with its tone, sound-waves pitched at

28

\section*{VIOLIN TONE-PECULIARITIES.}
inharmonious keys or "noise." As before, hold- ing the sounding-board to the lamp, with a pencil I circumscribe one of those areas, and therein, with a scraper remove a little wood. Now, in this area, light comes through more freely. I continue thus with scraper until this circumscribed area becomes equally luminous with adjoining areas. "Qoctor, how much wood did you thus remove?" I cannot answer this question with precision; but say perhaps the 100 part of an inch in some places; in other places of deeper cloudiness, a greater amount. You must understand the form of gradu- ation applied to this sounding-board to be 9-64 inch at position of bridge; thence, in either up ward o downward direction, down to 3-32 inch at the f edge as near as practical to work. Thus, absolute precision in diminution of sounding-board thickness- es becomes a matter of difficulty; in fact, this form of sounding-board graduation is not surpassed in diffi- culty by any other form known. You must also understand that where thickness equals 9-64, but little or no light passes through; also, that great- est luminosity, on this plate, occurs at the thickness of 3-32 inch; that between these two points, lumin- osity diminishes, or increases, with a regularity exactly proportionate to diminution, or increase of sounding-board thickness; therefore, irregularities in thickness produce cloudy areas by interfering with the passage of light rays. I recall rules gov- erning tone-pitch. I call up the fact that irregu- larities of 100th inch in the organ-reed, in organ- pipe, or on a tuning 1 fork, are things inadmissable.

29

\section*{VIOLIN TONE-PECULIARITIES.}
Is such irregularity also inadmissable on the violin sounding-board? Have I, at last, found the cause for "noise" in violin tone? My nerves are tingling. Forty years have I worked upon used violins trying to find a formula for eradicating "noise" therefrom; worked with the usual amount of "noisy" failure. Quicker than I can tell it, the intelligence was flashed to my sensorium that those slight irregularities of sounding-board thicknesses could produce a screaming, ear-piercing, soul-har- rowing, immeasurably high-pitched tone; one such tone for each irregularity in the tone-producing area of the violin sounding-board. You have heard how "prospectors" may "pros- pect for a lifetime without striking a "lead;" how, when they do strike a "lead" they immediately celebrate the event by going crazy. Gentlemen: The tone of this violin is not accom- panied by sound-waves pitched at inharmonious keys; and you now know the cause of "noise" in vio- lin tone. In connection with the word "violin," the word "noise" vividly recalls untold suffering; therefore, I shall press this subject upon you no further than is necessary for complete understand- ing of those high-pitched, dissonant overtones pro- ceeding from the imperfect violin sounding-board. Science states that musical sound is comprised, in tone-pitch, from 27 to 4000 vibrations per second; and, as violin "noise" is exasperatingly high-pitch- ed, it therefore follows that violin ''noise" pos- sesses a pitch exceeding 4000 vibrations per second. Science also states that sound may be pitched so

30

\section*{VIOLIN TONE-PECULIARITIES.}
high as to make its accurate measurement impos- sible. Thus those ear-piercing violin tones, (not legitimate, harmonic overtones,) become legitimate objects for "guessing" at their pitch; and when guesswork becomes legitimate, 'tis the "guesser's" own fault if his "guess" is too low; there- fore, I "guess" the vibrations of these suicide-sug- gesting "noises" to range anywhere from 20, 000 to 20,000,000 vibrations per second. In making this "guess," I admit that my figures may be large in proportion to the largeness of my ear; anyway, the opportunity to set these little, pestiferous annoy- ances out in their true light, becomes greatly sooth- ing to my long suffering nerves. 'Tis said revenge is sweet. I believe it. Until after the occurrence of this accident reveal- ing the cause of violin "noise," I never once sup- posed the violin sounding-board to be thus sensi- tive; yet enough related examples were near at hand to create precisely such supposition. The rusted piano wire, even only rusted in local areas, cannot, nor ever can, produce tone unmixed with "noise;" and the reason is plainly due to irregularities in wire-diameter caused by irregular oxidation. The law governing production of musi- cal sound seems explicit, and seems not to permit of even infinitessimal deviations in tone-producing agents. Again, an organ reed slightly out of "voice" is easily remedied. If its pitch be slightly too low, re- moval of an immeasurable amount of metal, from

31

\section*{VIOLIN TONE-PECULIARITIES.}
its free end, perceptibly raises its pitch; if too high, removal of only enough metal from its base to brighten the surface, lowers its pitch. I now believe a carefully graduated violin sound- ing-board to be equally as sensitive as either piano- wire, organ-reed, organ-pipe, or tuning-fork. "Science never made a violin." Gentlemen, please keep your seats. As a veteran violin-student, I am forced by ex- perience to acquiesce to the truthfullness in this startling statement, no matter how distasteful ac- quiescence may be. I am not among those who ignore science; but I am among those having learn- ed to accept scientific statements with caution. Science in chemistry, geology, mineralogy, astrono- my, botany, even science in violinology, so far as science can go, possesses fascination for me. Once I believed every scientific statement from every source whatever. I have lived to a' disillusion. Let me warn you of the danger in implicit belief of statements originating from human sources. There is but One Source for infallibility. In my day, chemistry was taught as an exact science, equally as exact as figures in arithmetic. Therefore, when I read in text books on chemistry that water is a combination of two elementary, gaseous bodies, hydrogen and oxygen, that water equals one equivalent each of these two bodies, that therefore the symbol for water is H 0, I be- lieved that statement to be infallibly correct. To- day the symbol for water is changed; so greatly changed that, without assurance, I could not know

32

\section*{VIOLIN TONE-PECULIARITIES.}
how H represents water. a Upon 1 a certain occasion at the University of Michigan, an incident occurred making an impres- sion on my mind to last until death. Twas at the hour for lecture on chemistry, wherein we were favored by the presence of an old gentleman who appeared deeply absorbed in. matter under discus- sion. Among other topics presented at this hour was distilled water, a gallon bottle of which . stood upon th6 lecturer's table. Instantly the lecture closed this attentive old gentlemen, with unsus- pected agility, climbed over railing to operating floor, pulled the cork from bottle of distilled water, applied his nose, shook his head,; shook the bottle, critically examined the "bead," then slowly in-- quired, "Is there T- any thing about this dis- tilled water intoxicating?"' Today I am the old man. You are the alert young fellows. In a subdued manner I ask, ' Is your new H= O. stronger than old HO?" Then, hydrogen was an element, therefore indiv- isible. Today hydrogen is separated into argon and boron. Pin not your faith upon human infal- libility; it is not in existence. Science is a Gollossus, yet, even today, science can't make a violin. Approaching the violin for its secrets, sci- ence, proud giant, receives a slap; on the face as a reminder that in the world there's one thing "none of his business. " In other directions, decade after decade, science advances with prodigious 'bounds; but in no scientific work can L find a solution for

33

\section*{VIOLIN TONE-PECULIARITIES.}
violin tone. English and German philosophers both dismiss violin-tone with a Latin evasion, "sui generis." French philosophers dismiss violin-tone by giving it a name belonging to all peculiar musical tone, "ft'w\&re." Science never made a violin. "Doctor, what did make the violin?" Experiment, alone, made the violin. Science came along later, as is often the case, with a post facto explanation of a few facts con- cerning violin-tone, but never an explanation of all facts concerning violin-tone. [Notwithstanding the statement of apologetic writers, I can find no corrobation for the statement that Stradivarius was a scientist. That he was not a scientist seems clearly established by his earlier work; that he was a tireless experimentalist, is proven by his history. That he is entitled to cred- it is proven by a few of his violins, I ask you, members of The Violin-Student Club, to state spec- ifically what invention for, or modification of the violin is due to Stradivarius. Silence? 'Tis no wonder.] In the violin bibliographic list are nearly 200 books. From these books only can we procure ev- idence. After wading a distance through this mass of evidence colored by writers who apologize for the Strad "chunks" by stating that those ab- normalities must have been built "on contract," colored by writers possessing exhuberant imagina-

34

\section*{VIOLIN TONE-PECULIARITIES.}
tion, who forgot to say they have no practical knowledge of either violin-building, violin-wood, or violin- varnish; colored by writers who only know the ' 'yellow" Cremona; again, by writers who only know the "red" Cremona; again, colored by trade promoters, one awakens to the fact that now he knows less of the Stradivarius violins than after having read but a single book. Neither. Gaurnerius nor Stradivarius appear to have improved violin-outline, nor detail of construc- tion, nor violin color- work, nor violin varnish. The first half of Stradivarius' life work was gone be- fore he reduced his table-arching to that given to the violin by Maggini 100 years before. Specific- ally, I can only determine that both Gaurnerius and Stradivarius reduced table-thickness to a degree wherein lighter strings could overcome sounding- board inertia and rigidity; and that, up to their day, their selection of sounding-board wood, for solo- violin use, was not equaled. Than these two facts, I can find nothing farther in favor of those two reputed violin-builders. If added facts in their favor can be shown me, I will honor them. That Italian music-builders, demanding three oc- taves of musical tone from each violin-string, gave the hint for reductions in sounding-board thick- nesses is extremely plausible. That Gaurnerius and Stradivarius were first in supplying such de- mand is in evidence. I think I may safely consider that both reduction in table-thicknesses and the selection of sounding- board wood of softer fiber, were then looked upon as

35

\section*{VIOLIN TONE-PECULIARITIES.}
innovations; and that, by other violin-builders, both of these men were looked upon as "cranks." But both men, as solo-violin-builders, were winners, al- though each adopted a different plan for accom- plishing 1 sounding-board reductions of thicknesses. Ole Bull demanded four octaves of musical tone from each violin-string. He found satisfaction in a ' 'Joseph. ' ' Paganni, it is claimed, delighted in light violin- strings. Paganni was the first great soloist to exploit the Gaurnerius violin. Honeyman, in an indefinite way, states that the Gaurnerius sound- ing-board is "thinnest throughout the central portion. ' ' It is a matter of interest to note the varying opinions of violin players concerning the varying tone in the violins of the great violin-builders. Ole Bull expresses surprise at seeing his contem- porary, DeBeriot, appear in public with a Maggini. In speaking of his Strad, Ole Bull substantially says, "I never play it in public because of its nasal twang," (See Ole Bull Memoirs, by Sarah C., his wife. ) I call your attention to notable omission in statements concerning the Stradivarius violins that the speaker or writer generally omits to men- tion in which of Strad 's "periods" his particular violin was made. I can only account for such omis- sion upon the hypothesis that the speaker, or writ- er, considers violins from only one of Strad 's four "periods" to possess tone- value worthy of remark. Gentlemen, you make a mistake if you interpret me as not joining in the praise universally accord-

36

\section*{VIOLIN TONE-PECULIARITIES.}
ed these two famous violin builders; but my praise is bestowed only for a definite reason that is, the great violins of these two great makers ; are only great as solo-violins. "Pis honor enough, and richly deserved for their bold innovation and correct judgement of sounding- board wood, My point is this: The unlimited praise bestowed upon these solo-violins has driven an army of vio- lin-makers to life-long effort at reproduction of solo-violins, until today, there is scarcely a good orchestra violin on the market. "Doctor, do you mean to say that the best solo- violin is not also a best orchestra violin?"

Yes, sir. "How is that?" The sounding-board of the best solo-violin is too light in wood to withstand the terrific force of sound-waves from full orchestra instruments. " "Doctor, please explain. Thus: The violin sounding-board to yield three octaves of agreeable tone upon each string, must be reduced in thickness to the point of weakness in presence of harmony- waves from full-orchestra; and because of such weakness, its tones are smoth- ered by overpowering harmony-instruments. "But, Doctor, we understood that first- violin sound-waves ride on top of harmony- waves. ' ' Correct, sir, so long as first-violin-waves have sufficient propelling strength to keep them on top. The swimmer, so long as his strength is sufficient, may keep himself on the crest of a huge wave, but

37

\section*{VIOLIN TONE-PECULIARITIES.}
when strength fails his next position is at the bot- tom of the trough; next, at the bottom of the. sea. To say, ' 'the violin reached perfection 200 years ago," means much and means little. To say the solo- violin reached perfection 200 years ago is near- er correct, and is not so misleading. We should insist upon definite specifications in this matter, because of different tone-characters being demand- ed from the violin. We only need a moment of attention to understand the absurdity of so much misdirected effort to produce solo-violins while the demand for orchestra violins is in the ascendant. Only a few solo-violins are in actual employment at any given date, while the general purpose-violin and the orchestra violin in large numbers, are in hourly employment. "Doctor, please explain your meaning by gener- al-purpose violin and full-orchestra violin." With pleasure, sir. Thus: The general-purpose violin is one whose tone-pitch, upon G and D strings, is distinctly of the bass character, while the tone-pitch of A and E strings is distinctly of soprano character: the full-orchestra violin is one whose G and D string tone-pitch is of baritone character, while the pitch of a A and E strings is distinctly of soprano char- acter. "Can these different tone-characters be given to violins at will?

Yes, sir. "Please tell us how." By combining, in the work of violin-building, cer-

38

\section*{VIOLIN TONE-PECULIARITIES.}
tain rules to be found in Lecture III, governing pitch of tone-producing agents. By reason of repeated, practical demonstrations, I am able to assure correctness of these rules. Before leaving the question of ' Violin perfection 200 years ago," let me direct your attention to two of its evil consequences. First: It acts as a discouragement to further experiment. Second: It acts to place a fictitious value upon a few violins of that date. Industrious "trade promoters" have lost no op- portunity for befogging the public in the matter of old, solo- violin values; thus the public erroneously demands that every soloist shall appear only with either a ' 'Joseph* ' or a ' 'Strad' ' under his arm. But, today there are unmistakable signs that the public is emerging from the fog; a few observing persons, having the courage of conviction, pronounce "pass- ing benediction" upon the "old violin." My own veneration, reverance, worship, call it what you will, for the "old violin" was cruelly shattered years ago by discovering that the violin may be- come "too old", also, by discovering the fact that a modern violin can be made to appear ' 'remark- ably preserved." But a few years ago the public was given an op- portunity for observing the inanity of continuing the demand that the soloist use only old violins upon all occasions. 'Twas at the Chicago Auditorium. My friend, Mr. Beebe, himself a violin-maker of repute, upon first occasion, seated himself quite

39

\section*{VIOLIN TONE-PECULIARITIES.'}
near the stage; upon a second occasion, seated him- self farther away from the stage; and, upon a third occasion, in the most distant seat; all for the pur- pose of judging tone from a Strad in the hands of a star soloist wise enough to permit only piano ac* companiment for his old violin. Mr. Beebe rejoices in scheduling a musically-trained ear, and I, of my own knowledge, have assurance that said ear, in the presence of violin tone, stands conspicuously forward. 'Therefore, when Mr. Beebe states that, while in a distant Auditorium seat, certain efforts of the great violinist were disappointing because of weakness in the tone of that old violin, I accept such statement as coming from a reliable authority. This incident shows the fictitious value placed upon old violins. That Strad is valued at several thousands a Beebe at a few hundreds. I value most that violin having the most satisfying tone. Gentlemen: Let me advise that you build violins with a view to some particular place of usefullness. Thus, when building for solo use, reduce sounding- board thickness to a degree yielding three octaves of musical sound from each string. When building a general-purpose violin, reduce dimensions of bar and sounding-board beneath >G and D strings to give those strings the basso tone-character, while leaving sufficient wood beneath A and E strings to give thein the soprano tone-pitch. In building for orchestra us\textasciicircum\{\}e, then leave sufficient wood beneath all strings to secure the baritone-soprano tone- character. ' < ' After fifty years of practical work given to violin

40

VIOLIN TONE-PECULIARITIES. tone-peculiarities, such are my conclusions in the matter of violin building. I call your attention to the fact that, in skillfully made violins, a wide range of tone- value exists; and to the fact that such wide range of tone- value is due to inherent differences in quality of sounding- board wood. Neither you, nor I, nor any person whatsoever, can exactly pre-determine the tone- quality of any given sample of sounding-board wood. In this fact lies the reason why science cannot build a violin. I call attention to the fact that 'tis an easy mat- ter to build the "noisy" violin sounding-board. For such work the equipment is contained in the following bill: 1 Three-quarter-inch-gouge. 1 Jack-knife. 1 Twenty-five-cent workman. For the latter "tool" we must resort to importa- tion from Germany. Even then it is a failure on this side of the Atlantic. When once that twenty- five-cent "tool" puts foot on American soil, "he" becomes a two-dollar man. Peace to thee, America!

41

\section*{VIOLIN TONE-PECULIARITIES.}
LECTURE III. GENTLEMEN: You understand the aim of our society includes the study of both musical and un- musical sound produced by the violin. It is not agreeable to admit other than musical tone as a pos- sibility for the violin. Instinctively, we are wont to consider "sweetness" and "violin" to be synon- ymous terms. Even to speak of the "noisy" vio- lin, is rasping to our sensory nerves, while enforced listening to such violins is positive torture. Under the heading, "Cruelty to Humanity," the sale and use of "noisy" violins should receive statutory pro- hibition. It is a lamentable fact that the most mu- sic-destroying sound can come from close imitations of the violin. The outward appearance of these fraudulent imitations is the trap set for the inno- cent purchaser. Because of appearances, the pur- chaser thinks himself in possession of a violin. Such deluded victims are today found in every direction. In the credulous ear of every one of these victims was sung that siren song, ' 'The violin always im- proves with age and use." In trying to apply this old song to the imitations of the violin, trouble be- gins. The victim's family is made to suffer from an attack of "nerves," and, .his immediate vicinity is wholly avoided by Music. His industrious"bow- arm is spurred to increased activity by hope of tone- improvement. Steadfastly and affectionately, he addresses that imitation as "old shell." The years roll by with no improvement in tone. Frequently gray-heads, with doubting mein, have brought these imitations to me. (Doubtfully,) "Doctor, can

42

\section*{VIOLIN TONE-PECULIARITIES.}
you tell me what ails my old shell?" Rapidly I run the bent "sound" across the inner surface of the sounding-board, the rattle therefrom suggesting the tinner's cart on corduroy road. "Yes, sir, I can tell you." "What is it?" "It's fraud." As I remove the sounding-board, permitting his doubt-expressing eyes to look upon the scene with- in his "old shell," the frozen smile on his wrinkled face is something truly pathetic. The graduation, (begging pardon for the use of that word, ) of this imitation sounding-board is entirely accomplished with gouge and jack-knife. The bar is a "chunk" left by haste of gouge, and thoughtfully smoothed by jack-knife upon the side exposed to your inspec- tion. Because you cannot see them, the two upper corners are forgotten. The linings, of various di- mensions, are both too long, and too short. Out- wardly, this fraud appears like the violin. Although the varnish is of cheap quality, the coloring is quite artistic. Here is a lifetime devoted to music, yet, wasted, wasted by this damnable imposition upon an inno- cent, credulous, unwary purchaser. Moreover, in this great, enlightened, (beg pardon once more, ) liberty-giving United States of North America, this same premeditated, unmitigated fraud is sold by thousands. Such gigantic imposition upon music- loving people is sufficient cause for murderous sen- timents within a heart of stone. These frauds are the product of combined cupidity, stupidity, and heart-

43

\section*{VIOLIN TONE-PECULIARITIES.}
lessness. Cupidity, in these frauds, is shown by charging more than kindling-wood prices, notwith- standing the fact that quotations are placed in the "pfennige" column. In these frauds, stupidity is shown in the fact that fifteen minutes more work upon the sounding-board might place quotations in the "marks" column. In these frauds, heartless- ness is shown by the premeditated murder of Music. The criminality of these fraud-builders ex- ceeds the criminality of one who robs little child- ren, because they combine murder with robbery. Poor Music! Standing aghast at sight of her fav- orite haunt no longer habitable! While professing friendship, you, Mr. Fraud-builder are the cause of her anguish. To you, and to such as you, and of you, I say ' 'Hell knoweth not a greater villain. ' '

\section*{Thou Violin! Thou Who art sweet Music s only King! Must thy witching form conceal Cupidity's rankling deadly sting? Men! Arise! Stand up now In ranks of those who wish to fight To kill this robbing deadly thing, And give to Music hers by right.}
"Doctor, how may the innocent purchaser know Cupidity's hand?" My friend, the case is hopeless when the purchas- er is so innocent as to be thereby unable to read, 'Made ' in Germany. ' ' After a lifetime given to the production of violin-

44

\section*{VIOLIN TONE-PECULIARITIES.}
tone, I am compelled to acknowledge the disbelief in the possibility of satisfactory solution for all problems involved therein. I am also compelled to state that my experience, and conclusions, in the matter of how the violin operates to produce musi- cal sound, do not add corrobation to the theories of Savart, Honeyman, et al. Although unable to solve all problems in violin tone, yet, I do not re- gard my efforts thereat to be entirely fruitless. I give you my conclusions and principles, so far as they are satisfying to myself, in terms as clear in meaning and as concise in form as I am able to command. I shall not indulge in abstract theories. My conclusions are all worded after actual, practi- cal, repeated tests upon the used violin, except the principles of sounding-board graduation for maxi- mum evenness of tone. As will be shown later, this principle, for reasons of physical disability, re- ceived but a single demonstration. Whatever of value may exist in my conclusions is hereby be- queathed to you. I now call up the matter of violin tone-pitch. [This feature of violin tone is worthy of careful consideration. Other things being satisfactory, tone-pitch should govern the place of usefullness for each violin. Thus: The violin of low tone-pitch throughout should not be placed at the head of full orchestra, and, for the reason that its tones have not sufficient intensity, (carrying power) to satisfy the distant listening ear. The place for the violin of low tone-pitch throughout, is that of the solo instrument. Again, the violin of high tone-pitch

45

\section*{VIOLIN TONE-PECULIARITIES.}
throughout, should not be employed as a solo in- strument, and, because its high-pitched tones, alone, are disagreeable to the listening ear. I have known the reputation of a creditable soloist to be ruined by his employment of a good, high-pitched, orchestra violin as a solo instrument.] I have formulated eight rules, each of which is capable of modifying violin tone-pitch. For con- venience of reference these rules are numbered consecutively from I to VIII. As frequent refer- ence will be given to these rules, they are therefore placed in a group. Rule /.Lengthening a tone-producing agent, other dimensions remaining equal, lowers tone- pitch. Rule //.Shortening a tone-producing agent, other dimensions remaining equal, raises tone-pitch. Rule III. Increasing thickness of a tone-produc- ing agent, other dimensions remaining equal, raises tone-pitch. Rule IV. Diminishing thickness of a tone-pro- ducing agent, other dimensions remaining equal, lowers tone-pitch. Rule V. Lengthening confined perpendicular air columns, lowers tone-pitch. Rule VI. Shortening confined, perpendicular air columns, raises tone-pitch. Rule VII. Enlarging sounding-board exits, raises tone-pitch. Rule VIII. Reducing sounding-board exits, low- ers tone-pitch. In practical violin building, or in violin toning,

46

\section*{VIOLIN TONE-PECULIARITIES.}
each of these rules possesses value. In my practice, application of these rules is given to the sounding- board alone. In no way do I apply any of these rules to the back plate. After long continued ex- periment, I am come to the conclusion that the chief function of the back is that of a reflecting medium. Therefore, I only demand of the back sufficient rigidity to withstand, without a tremor, the terrific charge of molecular wave-movement originating at the inner surface of the sounding- board; and, that its fiber be fine, and dense, and susceptible of high polish. Thus, when "fin- ished," the inner surface of the back presents a perfect reflecting surface. I do not say that the back cannot modify tone-pitch. On the con- trary, I state that the rigidity of the back may be reduced until the tone-pitch of the violin is lowered to a hollow, feeble, worthless degree. I have found many violins thus ruined; more thus ruined than by great reduction of sounding-board rigid- ity. Indeed, many violin builders attempt to reg- ulate tone-pitch by reducing rigidity of the back. In my experience with the long-distance test for intensity, (carrying power) of tone, this manner of regulating tone-pitch has proven itself seriously erroneous. The tone of every violin, thus regulat- ed, has fallen far behind in the distance record. Speaking for myself, regulation of tone-pitch by reduction of sounding-board rigidity, has proven itself to be the safer method. Therefore, these rules for governing tone-pitch are applied to the sounding-board.

47

\section*{VIOLIN TONE-PECULIARITIES.}
I will now consider the practical application of Rule I: "Lengthening a tone-producing agent, other dimensions remaining equal, lowers tone- pitch." In application of this "rule to the sounding- board, its effects are greater upon the tone-pitch of G and D-strings than upon tone-pitch of A and E-strings. As an aid to illustration, I will suppose that here is a sounding-board having a uniform thickness of 8-64. This thickness of sounding- board, together with bar dimensions, in length, lOi inches, thickness, 3-16, depth, I, tapering to ends, give 5 a high-pitch character to Gand D-string tone, not the highest, yet high. Desiring to lower th tone-pitch of the G and D-strings, I proceed to iengthen the tone-producing agent according to Rule I. For the purpose of limiting the tone-pitch to the G and D-strings, at one-half inch upon either side of the bar I draw pencil lines extending to as near the purfing and end blocks as is practical for reduction of sounding-board thickness. Within these lines, beginning at the position of the bridge, I gradually reduce the thickness down to 4-64 at ends of the plate; also, outside of these lines, thick- ness gradually increases to 8-64. Upon trial, the tone-pitch of the G and D-strings is found to be perceptibly lowered. Thus is the fact demonstrat- ed that, in the sounding-board of 8-64 throughout, the whole length thereof does not act with suffic- ient energy to produce audible sound; also, thus is demonstrated the correctness of Rule I. I call up Rule II: "Shortening a tone-producing agent, other dimensions remaining equal, raises

48

\section*{VIOLIN TONE-PECULIARITIES.}
tone-pitch. ' ' Evidently, the application of this rule must occur at the time of constructing the sound- ing-board; thereafter, it cannot apply. In the il- lustration for lowering tone-pitch, by lenthening the tone-producing agent, is sufficient data for raising tone-pitch by shortening the tone-producing agent. Rule III: "Increasing thickness of a tone-pro- ducing agent, other dimensions remaining equal, raises tone-pitch. ' ' This rule is applied at time of construction. I have found a few sounding-boards having a thickness of i at the position of the bridge; and two having the same great thickness through- out the length of the center join. The extremely high-pitched tone from the sounding-board thus heavy in wood, is not enjoyed even at a distance out of doors. In troubador days, such high-pitch- ed, out-of-doors tone was popular; in fact, quite necessary to satisfy the public taste. Today the street musician, taking pattern after "ye olden time," is yet successful in making the public part with its pennies by use of such high-pitched violins; and personally, I do not object, so long as he ac- cepts my penny as a hint to move on. (Honeyman again, is authority for the statement that the Nicolas Amatus sounding-board has a thickness of i inch throughout the length of the center join, thence, laterally, down to at the edges. ) Rule IV: "Diminishing thickness of a tone-pro- ducing agent, other dimensions remaining equal, lowers tone-pitch. " The high-pitched tone from the thick sounding-board is lowered by reducing

49

\section*{VIOLIN TONE-PECULIARITIES.}
thickness thereof. Lowering tone-pitch by reduc- ing sounding-board thickness requires caution, be- cause this work brings into employ two powerful factors, as will be shown after Rule V, towit: "Lengthening confined perpendicular air columns, lowers tone-pitch." Thus, in lowering tone-pitch by reducing sounding-board thickness, the same re- duction also lengthens perpendicular air columns within the body of the violin; therefore, two fact- ors are in employ; and hence the necessity for cau- tion. In a sounding-board of sonorous wood, with the thickness reduced to 8-64 throughout, a further reduction of 1-64 causes perceptible lowering of tone-pitch; and, from 7-64, the further reduction of 1-64, lowers tone-pitch two times more than the lowering effected from 8-64 to 7-64. I attribute such arithmetical increase in pitch-lowering to the simultaneous employment of two factors. Rule VI: "Shortening confined, perpendicular air columns, raises tone-pitch. ' ' For the purpose of raising tone-pitch of a violin whose table thicknesses are already at, or below the danger figures, I have successfully applied this rule by diminishing depth of ribs; thus shortening perpendicular air columns within the body. The fact that shortening perpen- dicular columns raises tone pitch is demonstrated in numerous familiar ways, as in the organ-pipe, the steam whistle, et cetera. As a practical dem- onstration of this rule, I have diminished depth of ribs by 1-16 inch, making a test at each reduction. The result affords ample proof of correctness in Rule VI.

50

\section*{VIOLIN TONE-PECULIARITIES.}
Rule VII: "Enlarging exits, raises tone-pitch." In practice, I only apply this rule after testing the tone of any given violin. Experience affords proof of correctness in this rule. As a matter of safety, the areas of the exits should be made rath- er small at the time of constructing the sounding- board. After testing the pitch of the finished vio- lin, should the pitch be found but slightly lower than desired, the fault may be remedied by a slight enlargement of exit area. In this work, I prefer to make such enlargement by removing wood from the inner edge of the exits, and for the purpose of bringing the exits nearer the focal points of sound-wave concentration. [Upon a later occasion, sound-wave concentration at the exits will receive consideration.] Rule VIII: "Reducing area of exits, lowers tone-pitch." Application of this rule should be made at time of construction, although its applica- tion to the finished violin is possible. Reduced area of exits is a powerful factor for lowering tone- pitch, and for diminishing volume of tone. For the production of the most agreeable violin tone for studio uses; I know of nothing among tone-pitch modifiers, of greater importance than the small exit. Musical sound is a fascinating subject to the vio- lin student. In that quality of sound called ' 'pitch" there is ample material for study. Some of the problems in sound are solvable; some are not solv- able. The problem of absolute pitch is solvable by mechanisms able to record the number of blows

51

\section*{VIOLIN TONE-PECULIARITIES.}
per second delivered upon air molecules. Thus, the pitch of any sound is governed by the striking agent, or, tone-producing agent. But, this fact by no means is a solution of all the problems involved in pitch. Absolute pitch does not cover all the ground. Words seem inadequate to even state all the problems in pitch of sound. Viewing the prob- lems of pitch from the standpoint of the violin, we are compelled to admit eight modifying factors, (as formulated in Rules I- VIII, ) in the production of violin sound, while but one of these factors is the striking agent! Tis hard work to formulate the problems in vio- lin tone; 'tis even harder work to solve these prob- lems after formulations. I bequeath their solution to you. By your side-long glances, and averted faces, I infer you do not estimate my bequest as a valuable asset. Perhaps such asset is notvaluable. 'Tis an easy matter to bequeath assets of which one is not "seized." But, because solution for violin- tone-problem is not now forthcoming, 'tis not certain that such solution never will be forthcom- ing. I advise continuance of effort. Little by lit- tle the solution of these problems may come to each persevering student. Myself, to continue the subject, must resort to assumption of fact, and mix it with a few positive facts, together with reasoning from analogy. Whether or not, such mixtures may be found agree- able to your palates, depends largely upon your en- thusiasm, together with my power of persuasion. Proceeding: A body may vibrate and yet not

52

\section*{VIOLIN TONE-PECULIARITIES.}
produce audible sound. To produce audible sound, the sound-producing agent must deliver blows of sufficient energy, upon contiguous air molecules, as to cause movement in every air molecule between the striking point and our tympani; otherwise we are not conscious of sound. These are facts. The following point is an assumption. Thus: I assume that the entire 14-inch length of the assembled sounding-board cannot act upon air molecules with sufficient energy to produce audible sound; but, I am compelled to be content with an approximation thereto. The difficulties in the way of precision in this matter are great. The differences in the elas- ticity of different samples of wood are infinite. Again, but slight increase in height of arch adds to sounding-board rigidity. As a mere assumption, I give the greatest length of sounding-board, act- ing to produce audible sound, as equaling 12 inches. In practice, to secure such 12-inch length, I reduce thickness from position of the bridge gradually down to 4-64 as near the ends of the sounding- board as is practical. I consider sounding-board thickness at the edges, when down to 4-64, to be at the limit of safety; and, this degree of reduction I only employ when desiring the lowest, practical, tone-pitch for G-strings. Should it happen that too great reduction of tone-pitch follows reduction of sounding-board thickness, a reduction of thickness sufficient to to cause a hollow, feeble tone, I know of no reme- dy by increasing sounding-board thickness; that is, no such remedy having value. I do know that in

53

\section*{VIOLIN TONE-PECULIARITIES.}
the city of Chicago there lives a man who claims ability to restore tone-pitch of violins by means of thin veneer of wood applied' to the inner surface of the sounding board. A violin, thus treated, was brought to me with the request that something be done for its dullness of tone. Not knowing what had been done to this violin, I first applied the bow as usual, expecting thereby to determine necessary work. But, its tone effectually took away all my conceit. In no way could I pre-determine the cause of such extraordinary dullness of sound. In haste I removed the sounding-board - Chicago! You, so careful not to keep your lights under a bushel! How have you lost distinction! Of your surgeons you have taken good care. Their reputation is world-wide. But the reputation of your fiddle-surgeon has not been equally nursed. By your permission, I will assist in correcting the oversight. First, I will describe the display of plastic surgery within this violin. Second, I will prescribe a course of treatment for the surgeon. The area covered by veneer equals I of the sound- ing-board; but, such areas are divided into two ter- ritories, one at either end of the plate. The thick- ness of the veneer equals 1-32. Its color is dark brown. In fiber it is like coarse paper. Its coher- ent strength does not equal that of good paper. It is cut in pieces of miscellaneous dimensions; some of the pieces being 2x4 inches, some down to by 2 inches. In places the veneer is doubled. In placing the fragments in position, no order of work prevails; the length of some being across the grain

54

VIOLIN TONE-PECULIARITIES. of the sounding-board; of others, with the grain, and others, oblique to the grain, while between patches are areas unaccountably overlooked. The scene arouses unbounded astonishment. Prior to this moment I prided myself on ability to determine the cause of violin tone-faults by use of the bow. Now is my pride humbled. 'Tis said one never becomes too old to learn. 'Tis a truth. Poor vio- lin! By your degradation am I learning. Before your pitiable condition became known to me, I had supposed those insatiable monsters, heat and moist- ure, were your most relentless enemies. But, you have come to teach me that I was in error, to teach me your worst enemy is man. And this man lives in Chicago! These patches of veneer are not attach- ed to the sounding-board by glue, but, by thick paste. Thanks for the little mercy! With cloth wrung from hot water the paste is easily softened, and I lift the various patches entire. Those patch- es, where veneer is doubled, fill me with surprise at the generosity of a Chicagoan. His bill might be equally large by saving the extra pieces. The whole affair is. inexplicable. How this man can live in Chicago and escape notoriety is one of the mysteries. But, the prescription, richly merited by this man, has nothing of mystery about it. It is a prescription dictated by justice. Here is a vi- olin without fault of serious degree, except that of model. The arching drops at the ends of the plates too suddenly for tone power; otherwise it is a fairly good violin. Its sounding-board is really a superior selection. After being again in playing order, I

55

VIOLIN TONE-PECULIARITIES. find this violin to be worth \$75. But, its worth, as it came from the veneering works, was less than seventy-five cents. I love Music! Sweet Music! I delight in doing her homage, even to the degree of worship. Therefore, when I see murderous hands reaching out towards Music, the lion within me arouses. Take this guilty wretch out on the lake, (not so far as to reach pure water,) tie a weight to his neck, tie a life-preserver to his feet, (otherwise that head will not go under, ) toss him overboard, hold him under until, by cessation of rising bubbles, it becomes certain that within his cranium there is an increment of gray matter.

56

\section*{VIOLIN TONE-PECULIARITIES.}
LECTURE IV. GENTLEMEN: At this hour I present one of the most absorbing questions pertaining to the violin. It is sounding-board wood. To some violin students the varnish question is particularly at- tractive. To others the color question may be par- ticularly attractive. Color study, by itself, is a definite department of art. No one ought to im- agine himself capable of mastering violin-color work in a brief space of time. Neither can mas- tery of color work be bought for any certain amount of cash. [I once had an amateur violin-making friend who enjoyed all the advantages afforded by wealth, and who lost his life in a second European tour search- ing for violin colors. Poor man! He vainly im- agined that money could buy a talent which is eith- er a gift, or is a result of long-time application. Color-sense is a thing not in commercial channels. He had not patience to practica color-work during the years required for the development of color- sense.] But, notwithstanding the great interest attach- ing to varnish and to colors, sounding-board wood retains paramount interest. It seems to be the universal opinion that quality of sounding-board wood largely governs violin tone-value. Everything in, or about, sounding-board wood has been subjected to hawk-eye scrutiny. In such scrutiny, science has given but meagre assistance. To experiment alone, are we indebted for a limited knowledge thereof. Incidentally, science offers a

57

\section*{VIOLIN TONE-PECULIARITIES.}
mite of corroborative value in the record of con- ductive power for sound-waves in air, gases, wat- er, metals, and different varieties of wood. Bot- any contributes a mite in the fact that heart- wood capillaries become smaller as a tree advances in age, therefore heart-wood fibers possess greater spring-action. [To be precise, this statement does not include fibers immediately at the heart, but fibers near to the heart.] As a matter of interest to the violin student, I quote a few points from the conductivity table in that admirable treatise on the theory of sound in its relation to music by Peitro Blaserna, Royal Uni- versity, Rome. [This work is complete in detail and faultless in diction except wherein the translator, throughout the book, with but two lonely exceptions, erro- neously uses the word "note" when meaning "tone"] Blaserna's figures are given in the met- ric system, towit:

\section*{"Velocity of sound in various bodies." 'Air 32 Fahr., meters per second, 330." 'Copper 68 Fahr., " meters " per second, 3556." 212 " 3295." 392 2954." 'Acacia wood, along the fibers, 4714. ' ' " across the rings, 1458. ' ' with the rings, 1352." Tine along the fibers, 3322. ' ' across the rings, 1405." with the rings, 794."}
58

\section*{VIOLIN TONE-PECULIARITIES.}
These figures present some interesting facts to the violin student. Thus: (1) Sound-waves along the fiber in pine, travel with more than ten times the velocity traveled in air. (2) In copper (and in all other metals) as heat increases, distance traveled by sound-waves diminishes. [The conductivity of metals, diminishing under higher temperature, is here introduced for the pur- pose of calling attention to the fact that both high and low temperatures greatly modify the conduc- tive power of all bodies. Thus, in the increased temperature of the air from midsummer heat, all sounds, musical and otherwise, can be only propelled to comparatively diminished distances. The same phenomenon occurs in the low tempera- ture of winter. Thus, the intensity, or carrying power of a violin should not be determined by a test given in either extreme of temperature.] But one variety of wood exceeds pine in the pow- er of conducting sound. That wood, acacia, or cin- namon, is not available for sounding-board uses. Therefore pine is placed at the head of the list. But as there are several species of the genus pine, we therefore must choose between them. It is un- fortunate that the particular species of pine, em- ployed in making the record, is not given. We may legitimately suppose that the sample used for the record was taken from near the heart of a mature tree, because, as stated in botany, capillary tubes of heart-wood in the mature tree, diminish in caliber with age until they cease to carry sap. Therefore heart-wood possesses greater density

59

VIOLIN TONE-PECULIARITIES. than sap-wood. Heart-wood also possesses much greater spring-action than sap-wood, and also great- er weight. But, when seasoned, sap-wood may be harder to cut than heart- wood. Mere hardness in cutting, so far as my observation extends, is a quality detrimental to highest tone-quality. Upon examining the new stump of a mature tree, the recent yearly growths beneath the bark are not yet well marked into hard fiber and connective tissue, but appear more as a homogeneous mass. Ap- proaching the heart, the divisions of grain become well defined. Long experience demonstrates be- yond all doubts that wood of well defined grain yields the better single tone, and immeasurably the better double-stop tone. Experience also dem- onstrates beyond all doubts that heart-wood possesses the greater spring-action. Thus, when forcibly bent, upon release from force, heart-wood returns to the point of rest with the greater rapid- ity. The point of value in the Indian bow lies in the rapidity with which it returns to the point of rest. Therefore heart-wood is selected for such bows. As youths, many of us learned how unsat- isfactory is the bow made from either an immature tree, or from sap-wood of the mature tree. Such dissatisfaction came from the slowness in which the bow returned to the point of rest. The action of such bow is weak, and the arrow can only be projected in a weak manner. In the sounding- board of immature wood I find precisely similar weakness; also, weakness of tone from sounding- boards of mature trees not possessing well defined

60

\section*{VIOLIN TONE-PECULIARITIES.}
grain. In my fifty years' work in re- toning used violins I have carefully scrutinized many varieties of wood. In doing such work I have opened several hundreds of violins; and am now satisfied that this plan af- fords more satisfactory evidence bearing upon tone-peculiarities than any other possible plan. I think this plan unquestionably discloses the reason why many skillful workmen fail in securing a great- er number of tone-art masterpieces. Barring mod- el and workmanship, I now believe that to quality of sounding-board wood is due every violin master- piece of tone-art. By the words "masterpiece of tone-art" I mean the limit of tone beauty. Tis not enough that the single-stop tone is beautiful. All double-stop tones must also be beautiful. In this latter demand is where the sounding-board oftenest displays its lack of highest value. If you have sufficient patience with my slow ways, you will know, in due course of time, my reasons for ascribing violin tone-value to the sounding-board. I shall not claim infallibility for such reasons. I only shall claim such reasons to be satisfying to myself. In the course of my experience with used violins I have seen changes on the interior sounding-board surfaces not generally known, as I believe. As an instance, I have neither heard of, nor read about unequal shrinkage of sounding-board fibers after the violin had left the builder's hands. Yet, I have seen cases, in violins of faultless workmanship, wherein unequal shrinkage of the sounding-board

61

\section*{VIOLIN TONE-PECULIARITIES.}
caused serious injury to tone. I call up a certain violin having a clear record for 150 years, a violin unquestionably built by a workman, a violin where- in unequal shrinkage was located beneath the E string. These inequalities in thickness ran along the fiber in the form of ridges and valleys, and were due to shrinkage of certain fibers in a greater degree than in neighboring fibers. There were three such ridges and two valleys, and their loca- tion, either, alone, might injure E-string tone. The only possible remedy lay in re-establishing uniform thickness in that particular area. I assure you that no further work, of any kind, was needed either upon the interior or exterior of this violin. The result of re-establishing uniformity of thick- ness in this area was the removal of ''noise" from E-string tone. I am satisfied that the original graduation work upon this sounding-board was done with precision. [Although not here to the point, yet 'tis fighting against my nature to omit notice of the fact that the bewitching sweetness of tone, now possessed by this old violin, is something to keep in memory. Its power of tone is but slightly diminished by age and use. So exactly precise was sounding-board thickness guaged for the guage 2 E-string that the slight amount of wood removed from beneath that string resulted in a perceptible diminution of tone- power thereof. One example of unequal shrink- age in the sounding-board, after leaving the build- er's hands, may be sufficient explanation of one reason why some well made violins, possessing

62

\section*{VIOLIN TONE-PECULIARITIES.}
superior sounding-board wood, fail of becoming tone-art masterpieces. Yet, I will offer another example of shrinkage in sounding-board wood per- haps still more striking. Thus: From a certain plank, in my possession for forty years, I fashion- ed a sounding-board. For reasons, this sounding- board was not used until two years thereafter. My surprise was unbounded at discovering additional shrinkage occurring subsequent to its graduation. This fact affords proof that quantity greatly modi- fies shrinkage of wood while seasoning. The thin shaving from plank, however dry, may yet further dry out. The thin shaving also swells the sooner in the presence of heat and moisture. These facts are of interest to the violin student. Thus, no matter how long the rived block, or plank, may have seasoned, we must yet consider such block, or plank, as new wood in a degree; and, when the sounding-board, from block, or plank, is reduced down to thicknesses considered as the limit of safety, we must not be surprised at yet further shrinkage. Neither must we be surprised at the greater rapidity of swelling in the presence of heat and moisture. Heat and moisture, as will be shown at- a later hour, are relentless enemies of the violin. They also combine to defeat intention of the violin builder.] Because vapor of water can greatly diminish both resonance and brilliance of tone, I therefore depend upon a hygrometer to note the per cent of satura- tion existing in the air in my workroom, even when retoning long used violins. As moisture in

63

\section*{VIOLIN TONE-PECULIARITIES.}
wood is drawn out by dry air, I am therefore care- ful to employ artificial heat when necessary to se- cure dryness therein. After becoming satisfied that all moisture in wood is drawn out, then I am ready to do the work of hermetical sealing describ- ed upon a later occasion. Unequal shrinkage by no means occurs in all sounding-boards. How to pre-determine unequal shrinkage I do not know. Undoubtedly the builder of this violin never even dreamed that future shrinkage of a few fibers in the sounding-board would occur to the injury of tone. I consider the workmanship displayed upon the interior of this violin to reach the limit of human skill. That its sounding-board is a good selection, otherwise than those few shrunken fibers, we may know by the beautiful tone. The following point also contains something of interest to the violin student, that is, the faulty tone in this violin never could have been improved by the use of the bow. Nothwithstand- ing the age and workmanship belonging to this violin, because of tone faults from unforseen caus- es, its value was inconsiderable. Today its value is changed. He who can buy this violin now, is al- ready rich. Wherein lies its value? 'Tis not in tone-power. Its tone is neither powerful nor weak. The question is difficult to answer. 'Tis not in age. There are violins of greater age, but of less value. 'Tis not in sweetness of tone. There are violins of equally sweet tone, but of less value. The answer de- fies enunciation. Inadequately, the question is ans- wered by saying, ' 'Its tone arouses a feeling of ten-

64

\section*{VIOLIN TONE-PECULIARITIES.}
derness." Tender feelings make us part with wealth. I know of no other logic explaining why one will part with more wealth for a certain sweet ton- ed violin than for another equally sweet in tone. Therefore the price paid may be taken as a meas- ure of the purchaser's tenderness. But, why, or how, one sweet toned violin arouses greater tender- ness of feeling than another sweet toned violin, re- mains to me an unsolved problem. The fact exists. That's all I know about it. Yet, there is something further that might be said about it. Thus: When the tone of a violin arouses tender feelings within the performer, then, and there-for, will the performer give us his utter- most musical expression. In selling such violin, the owner adds a price for his feelings. In the course of my work upon used violins, I made some acquaintance with sounding-board wood from different countries, as Scandinavia, Russia, France, Switzerland, Austrian Tyrol and Italy. Briefly, the most resonant sounding-board wood, coming under my observation, grew in Austrian Tyrol; yet, some samples from Switzerland and Italy have proven difficult to surpass in this im- portant quality. Undoubtedly, heat and moisture greatly modify evenness of yearly growth, density of fiber, and the presence, or absence of fat in the genus pine. Because altitude modifies the amount of both heat and moisture, it follows that pine, of widely varying widths of grain, and varying amount of fat, may be found in all mountainous lo- calities where pine grows. Therefore, every sam-

65

\section*{VIOLIN TONE-PECULIARITIES.}
pie of pine whatever, from any country whatever, is not of value for sounding-board purposes. For such purpose, pine presents five serious faults; and the finding of even a single tree without one of these faults, even in the most favored localities, is a matter of difficulty. Those five faults are: (1) Unevenness of grain width. (2) Crooked grain. (3) Unmarked grain. (4) Too great density of grain. (5) Fat, or pitch. From some parts of Norway comes pine of per- fectly straight, even grain; but, while sonorous in a marked degree, its density is so great as to give a somewhat disagreeable, stinging quality to tone. Were its density less, I believe no better sounding- board wood can be found. Presuming you to be familiar with foreign grown sounding-board wood, I pass on to the considerat- ion of domestic wood. Foreigners pronounce our wood to be of no value for sounding-board use. They are mistaken. Colorado, in a limited way, offers spruce of unsurpassed value for sounding- board purposes; while Michigan, also in a limited way, offers pine vastly superior to the majority of "ready-made" European sounding-boards kindly sent over to us. The kind of Michigan pine to which I refer, is peculiar in the color of grain. Thus, in European pine the different grains are se- parated by dark lines; but, the grains of this Mich- igan pine are separated by white lines. Another peculiarity in this Michigan pine is the fact that its

66

\section*{VIOLIN TONE-PECULIARITIES.}
white line is the softest part of the grain. There is another peculiarity yet more striking. It pos- sesses, as transverse markings, numerous scales of silvery brightness; and, as I know by observation, these transverse markings retain their brightness more than half a century. Except in mere power of tone, this quality of domestic pine, for sounding- board use, is unsurpassed by any pine whatsoever. The tone from this wood, for studio, parlor, or the smaller audience rooms, cannot be excelled so far as my observation extends. Even in new violins, double-stop tones from this wood possess exceeding attraction. Its value is only lacking in great pow- er. Upon this wood, varnish must be applied sparingly. It has no coarseness of tone to be smothered. For the bar and sound-post, I have used no other wood in many years. As the bar, from no other wood can I obtain equal beauty for G-string tone. My experience with white cedar for sounding- board uses is quite limited, having given it but two trials. For these trials I used cedar which had seasoned out of doors during more than thirty years under my own observation. One marked peculiar- ity of this wood is its power to resist disintegration from heat and moisture. Even after standing un- covered for so many years, its surface only show- ed but slight traces of decay, while beneath, the color was exceedingly bright. The presence of deep cracks interfered much in the work of secur- ing samples for sounding-board purposes. Indeed, only by joining four pieces could I get a perfectly

67

\section*{VIOLIN TONE-PECULIARITIES.}
sound, straight grained sample. The rapidity with which this wood grows produces great width of grain. In my opinion its width of grain, as a rule, is too great for high value for the sounding-board. As a mere supposition, its width of grain might not be objectionable for the bass sounding-board. The tone of this wood, from my two experiments, was peculiar. Volume of sound was very marked, but intensity was very feeble. Thus, at near-by dis- tances, the sound was great; to long distances its sound could not be forced. Intensity of sound is something defying explana- tion. We know that sound may possess marked volume while only able to propel itself to short dis- tances; and again, sound may possess diminished volume while able to propel itself to greater dis- tances. Undoubtedly the cause for this phenome- on lies in peculiarity of the sound-producing agent. The human voice affords abundant examples of the differences between volume and intensity of sound. Some voices can be easily propelled to distances impossible to other voices; and, the voices traveling the greater distance, may not sound more than half so loud as voices at near-by points. The same phe- nomenon is presented by the sounding-board. The quality of intensity cannot be pre-determined in any sound-producing agent. Marked intensity of tone is a valuable asset of the violin. It is an asset largely governing prices. As a mere business prop- osition to the soloist, the value of the violin should be based upon the square of the distance to which its tone can give pleasure to the listener. Thus:

68

\section*{VIOLIN TONE-PECULIARITIES.}
The tone of violin A gives satisfaction to the listen- ing ear at 400 feet. Employing 100 as a unit, the square of 4 units equals 16. The tone of violin B gives satisfaction at 800 feet; and the square of 8 units, equaling 64, therefore the actual, practical value of these two violins is as 16 to 64.] The tone from these white cedar sounding-boards was of a dis- agreeable, coarse quality. Strangely, a few violin users pronounced such coarse tone to be quite sat- isfying. Heavy coating with varnish was required to subdue coarseness of tone; but, in thus removing coarseness, loss of tone-power followed in a mark- ed degree. There was yet one peculiar feature in the tone from these sounding-boards. Thus: Powerful bow-pressure produced no greater tone than moderate bow-pressure. For these reasons I do not consider white cedar the equal in tone- value with the Michigan pine herein described. That the violin sounding-board is responsible for violin tone-value is a question no longer in dispute. Therefore, to the violin student, wood for sounding- board purpose becomes a matter of chief import- ance. Seemingly, a life-time given to re-toning used violins might enable one to pre-determine with accurracy the tone-quality of many grades of pine and spruce. Although such experience af- fords grounds for close guesswork, yet I assure you of my disbelief in the possibility of precise pre-de- termination of the tone-qualities of any given sample of sounding-board wood. In my experi- ence, the surprises awaiting test by the bow have been infinite. I will describe a grade of pine pos-

69

\section*{VIOLIN TONE-PECULIARITIES.}
sessing the nearest approximation to uniformity of tone-quality. But, it is necessary to bear in mind the fact that my description is based upon physical qualities and appearances of pine long seasoned after being fashioned into sounding-board form. This necessity arises from the fact that rapidity in color- changes is greatly modified by quantity of wood. Thus, original brightness of color is longer retained in the log than in the rived block; also longer re- tained in the block than in the sounding-board. Therefore color cannot be dependable evidence in determining the time elapsed since any given sam- ple of wood was cut from the stump. And further, different samples of sounding board pine take on widely differing depths of color with advancing age. I have one sounding board, of respectable age, showing but a mere shadow of color-change beneath the surface. In others the color-change is marked, having passed to brown-red through and through. The cause for color-changes coming on during the time of seasoning is a mystery baffling explanation. But, by long observation, I know that different tone-qualities follow different depths of coloring in the sounding-board. Thus, from the sounding-board longest retaining original bright- ness, the tone is of a certain stinging quality de- manding the utmost of careful bowing; whereas, the tone, from the brown-red, is quite agreeable whatever the bowing. The greatest difference is shown in double-stop tones. From the bright wood, the most beautiful effect from double-stop tones are impossible, while from the brown-red wood,

70

\section*{VIOLIN TONE-PECULIARITIES.}
double-stop tones possess the highest degree of agreeableness. Than beautiful double-stop tones, nothing from the violin gives greater pleasure to the listener. From the violin-soloist, double-stop tones are imperatively demanded. Therefore in selecting a violin for solo use, this knowledge be- comes valuable. Because the brown-red sounding- board easily is split, we may suppose that the con- nective tissue between the denser fibers, is of feeble power. We may also suppose that the absence of power in connective tissue permits greater independence of fiber action, and that to independ- ent fiber action, agreeable double-stop tones are due. The bright wood is not easily split. Its con- nective tissue is much the stronger. Hence its fibers, or denser part of grain, has less independ- ence of action. As an assistance in accounting for the brown-red color in pine, long seasoned, I offer the following: Back in the fifties of the last century, a certain barn and shed, several hundred feet in length, were enclosed with Michigan pine. In those by-gone days only large trees in pine forests were felled for lumber. Then the old fashioned "sash" saw was in the heights of its glory. Those large logs were cut "through and through, " and each board, or plank, was "edged" with a smaller saw. Thus the lumber was of various widths, and, without assort- ing, as in the present day, was sold to the consum- er. Much of the lumber covering this particular barn would be graded today as "first clear. " Many of the boards were 30 inches in width. "Dressed"

71

VIOLIN TONE-PECULIARITIES. lumber was then unknown. So was paint al- most. The covering of this barn was neither "dressed" nor ever painted. Forty-five years later I examined these unprotected boards. A majority of the wider boards were split at the cen- ter. I could easily determine every board coming from the center of the log by the evidence of disin- tegration. Thus: Upon the surface of all those ' 'center cut' ' boards, the soft connective fiber has decayed, turned to dust, and decayed to varying depths on different boards, or rather on boards from different logs. The denser part of the grain remained as ridges. [Some old, worn violins present similar ridges at the ends of sounding-board; and some do not. The difference is plainly attributable to varying degrees of toughness of connective tissue. In my observation, those old sounding-boards, showing greater loss of connective tissue, invariably pos- sessed greater double-stop tone- value, and greater mellowness of tone generally. Also, such sounding- boards are more easily split.] On all those weather-worn boards coming from that part of the log near the bark, the ridges of denser fiber are much less prominent. Indeed, some of these boards show no ridges whatever. One striking 1 peculiarity of these ridges is in their line. This line, on the great majority of those old boards, follows a zigzag course. Only upon a very few of them is the ridge-line straight. I am now come to the color question. As I dress the weather surface of all "center-cut" boards, there appears a great variety of colors. But, only upon the widest

72

\section*{VIOLIN TONE-PECULIARITIES.}
boards, do I find the brown-red. Unquestionably those 30-inch boards came from the mature trees. I now return to those boards showing straight ridge-lines. Their color runs from pale to butter color, and their general color-effect is bright. In a very few of them I find silvery bright, transverse markings, or scales. The evidence afforded by these old boards points towards age in the standing tree as the source of brown-red color developed by seasoning. In wood- craft there is much of value to the violin student. Therein may we learn that the connective tissue in young, or immature trees, is very much tougher than in the mature tree; also, that immature wood, while permitting greater degrees of bending, yet, when released, remains partially bent; whereas, mature wood, permitting less bending, but, when released, quickly returns to its original point of rest. This feature of quickly returning to the point of rest is of great importance to tone-pro- ducing agents. Experience clearly demonstrates the fact that immature wood, for sounding-board use, is inferior to mature wood. In my violin re- toning experience I have found sounding-boards made of two samples of wood taken from trees of widely varying age. Such sounding-boards have in- variably given proof that mature wood possesses the greater tone-value. Thus, in such sounding- boards, when the left half is of mature wood, while the right half is of immature wood, the G and D strings will yield good musical tone, while A and E strings will yield but a "dead," or lifeless tone;

73

\section*{VIOLIN TONE-PECULIARITIES.}
and vice versa, the G and D tone will be lifeless. I am how come to a fact concerning violin sound- board action which seems to have not attracted much attention. I allude to independent action of contiguous fib- ers. From my point of view, other facts being correct, independent action of contiguous fibers, in the sound -producing area of the sounding-board, makes violin tone- value. This view is based upon observation of the fact that agreeable quality, can- not be produced from wood having rigid connect- ive tissue. The latter fact I have observed in many cases. The following considerations explain necessity for independent fiber action in the sound- ing board: The open tone A requires for its pro- duction that certain fibers of the sounding-board strike 450 blows per second upon contained air; open tone E requires 675 such blows. For the sim- ultaneous production of A and E tones, contiguous sounding-board fibers evidently must vibrate at different rates per second. The difference for these tones is the difference between 450 and 675. Therefore, those fibers producing E, move 225 more times per second than the fibers producing A. That the fibers producing the tones A and E are contiguous will be shown upon a later occasion. In this case the evidence seems conclusive. Again I call your attention to pine possessing the brown- red color as being a grade offering the minimum of connective tissue rigidity. As further evidence bearing upon the question of independent action of

74

\section*{VIOLIN TONE-PECULIARITIES.}
contiguous fibers, I present this violin which has been in use 40 years within my own knowledge. Two times have I opened this violin in the effort to make its double-stop tones agreeable. By apply- ing a bow, you instantly perceive that my efforts are failures. I pronounce agreeable double-stop tones from this sounding-board to be an impossi- bility; and the reason for such impossibility is due to inherent rigidity of connective tissue. This vio- lin has something of a history; and, because the history of some violins is the only asset of value they possess, I present the history of this one as evidence. Upon first opening it, I found beneath an accumulation of dirt a legendary label bearing the following: Andreas Guamerius Sub titulo Santa Thtresia 1645. In appearance, this label is the embodiment of inno- cence; and, in outline and profile, this violin is a fac sjmile of the Andreas. In my earlier observa- tions of the violin, this one is remembered as giving swelling pride to Prof. . Many times I had gazed upon it in open-mouthed awe. Wanted to own it? Why, yes! But, I didn't have the wealth of Croesus. Another young fellow was more for- tunate. Fortune smiled upon this particular young fellow. Prof. found necessity for moving to other parts. But, before he could make such move, his unreasonable landlady exacted pay for his board. Under such distressing circumstances the Prof, regretfully "came down" in his price, and that young Croesus became the proud owner

75

\section*{VIOLIN TONE-PECULIARITIES.}
of an Andreas Guarnerius. Time not only rolls around swiftly, but also brings swift changes. This young man, becoming obliged to earn a living, laid his Andreas away in the garret together with his other violin of common value. Mice always prefer to live in the garret. 'Tis the warmest apartment in the house. Hitherto mice were not considered as violin "experts." But now, proof of their expert judgment of violins is established by the fact that they gnawed their way into the Guarnerius. Hence this violin came to me for re- pairs. [Although not an old violin expert by any means, yet, concerning the modern violin, in the 40 years since I first saw this one, I have learned something. Therefore I am not surprised to read upon the in- ner surface of this sounding-board, and written with pencil upon the wood, the following words: Johann Winkerline, Mittenwald Den ersten Oktober, 1853. [The suggestion of Mittenwald makes it quite certain that Andreas died some time before making this violin.] I reduced thickness of sounding-board to what I consider the limit of safety, but only succeeded in removal of a slight amount of its woody tone. After a period of use I again opened this instru- ment of torture, and gave the sounding-board an hour of massage.. [As you know, massage is quite often applied to- day to human sounding-boards for a fee. When such fee is large there is but little difference in the

\section*{7G}
VIOLIN TONE-PECULIARITIES. tone of these two varieties.] This treatment did just perceptibly diminish dis- agreeableness of double-stop tones, yet, in a small, bare room they remain suggestive of self-destruc- tion. I think my experience warrants me in positively stating that, in violins having sounding-boards of a rigidity and density equaling this sample, despair and death precede tone-improvement. ''Doctor, do you mean to state that violins do not always improve with age and use?" I mean to state that one long life-time is not suf- ficient to witness tone-improvement in some violins. "How then may we select new violins certain of tone improvement?" By ability to correctly judge sounding-board wood, and to judge the graduation thereof by the tone therefrom. "But, Doctor, all this requires experience?" Certainly, sir. But in lack of experience it is quite safe to trust the experience of some skillful, ambitious, conscientious, tone-knowing, violin- playing, violin-loving, violin-maker. You have my assurance that such violin makers are not yet all in heaven.

77

\section*{VIOLIN TONE-PECULIARITIES.}
LECTURE V. GENTLEMEN: The effect of age upon violin tone is a subject filled to the brim with interest. Among all musical devices, in the matter of tone-improve- ment by age and use, the violin pre-eminently stands alone. Strangely, tone improvement in the violin follows loss of value in wood in its use for other purposes. Because of such steady diminution in value of wood, unlimited tone-improvement for the violin becomes an impossibility. The fact of loss in wood-value affords a pathetic side to violin history. The best violins have been first in suc- cumbing to those insatiable enemies, heat and moisture. The inexorable law of disintegration has stayed not in its hand for best of the Maggini, nor best of the Amati, nor best of the Guarneri, nor best of the Stradivari. The better tone-producing sounding-board sooner yields to attacks of its en- emies. Soon will the best old violins only be known in memory. The many hours of pleasure contributed to hu- manity by those worn-out violins is something be- yond expression in figures. 'Tis well for humanity that the amount of such pleasure cannot be com- pressed into one hour. Such compressed sweetness could depopulate the world. 'Tis a peculiarity of humanity to love most that which gives the most pleasure. Of all inanimate things contributing to the sum of human pleasure, the violin stands at the head of the list and without a rival. From with- in the "dugout" on Western ranch to beneath the gilded domes of Czar, the violin carries sunshine.

78

\section*{VIOLIN TONE-PECULIARITIES.}
Deep down in recesses of human affection the vio- lin finds ample room. Must this priceless thing remain a sacrifice to the law of disintegration? Ought not humanity bestir itself to find a safe de* fense for the violin against this law? In the light of my experience in searching for means of such defense, I say; " Tis but folly to sit down and cry while doing nothing to save the vio- lin." In the light of my experience, I say to every man who, without experience condemns all efforts to preserve interior violin surfaces from in- evitable "ravages of disintegration," "You are doing an act of criminal violence to Music." To such as have made efforts to thus preserve the vio- lin, and who have abandoned such efforts because of injury to tone, I say, "Persevere." Do as much in such efforts as I have done in ten years. Then, if you are yet dissatisfied with results, possibly my method of preserving those interior surfaces from disintegration may be of interest to you. If called upon, I now consider myself amply provided with proof that interior vio- lin surfaces may be indefinitely protected from dis- integration, and without injury to tone. Should my caller be not prejudiced too much, I hope to convince him that the tone of such protected vio- lins remains equally sweet, while the tone qualities of brilliance and intensity are greatly augmented. There was once a man whose genius as an imit- ator of the Strad violin, not only filled the violin world with amazement, but also, filled this man's pockets. This man, J. B. Vuillaume, impostercon-

79

\section*{VIOLIN TONE-PECULIARITIES.}
fessed, said that interior violin surfaces cannot be protected without injury to tone. Because this man was an imposter, therefore when he said, "No," possibly he meant, "Yes." [Strange how an imposter's ''No," to this day, influences the violin world. As I view this mat- ter, 'tis but a display of inanity when I act upon proffered advice from a known swindler. As I think, the statements of men, even under oath, who have no other object in life than money, should be taken with the proverbial grain of salt. Indeed, there are cases wherein a pound is better. In the latter class of cases am I inclined to place J. B. Vuillaume.] As a member of the staff at a violin sanitarium, it came in my way to see somewhat of the pathetic side afforded by violin interior surface disintegrat- ion. Truthfully, the sight of such destruction of violin value became the stimulating agent for my later effort to find a safe means and a meth- od for its prevention. Often have I witnessed total ruin of tone-value for no other reason than that of leaving interior surfaces of the violin unprotected from disintegration by the action of heat and moisture. The question, ' 'Can such dis- integration be prevented without injury to tone," stood in large type before my eyes until I acted against a universal conclusion. Now, I am surpris- ed that injury to tone, from protection of violin interior surfaces, is but an imaginary bugbear. I do not state that tone cannot be injured by the in- terior protection. On the contrary, I well know

80

\section*{VIOLIN TONE-PECULIARITIES.}
that interior-surface protection can be the cause of tone-value ruin. I know full well that crude means, and crude methods of applying means to interior violin surfaces can injure tone when ap- plied upon exterior surfaces; but, not more; and, not less. But, in accounting for tone-injury, the following fact must not be overlooked. Thus: What ever injures tone when applied upon a single sounding-board surface, multiplies such injury by the factor 2 when applied to both surfaces. The point is, neither means, nor method of applying means for protection of either surface, resulting in injury to tone, should be used upon the violin. As musicians of experience, you well know that 'tis not what you play, so much as how you play, that wins the "encore." Had I sat down and cried quits after the first few attempts to apply protection to those interior surfaces, I yet might be contemplating this bug- bear from the usual distance. Had I then aband- oned such effort, I never would have received those marks of approval from the 60 owners of violins thus protected. Had the tone of violins thus treat- ed, been injured by such treatment/then, approval from such owners would have been withheld. In every case my fee was subject to approval. In no case have I lost my fee. I only mention these lat- ter facts as proofs that violin tone is not injured by protection of interior sufaces. [Perhaps 'tis proper to state that I worked upon violins for more than forty years with no intent, or expectation whatever of turning- such work to bus-

\section*{Si}
\section*{VIOLIN TONE-PECULIARITIES.}
iness account. I only charged a fee when violins came to me in such numbers as to demand much of my time.] It is my belief that to apply protection to either exterior, or interior surfaces of violin plates, with- out injury to tone, demands both experience and judgement only obtained by experience. "Doctor, how long will your protected violins last?" I cannot answer that question further than to say I know of no reason why the material em- ployed by me as interior-surface protection will not last equally with exterior surface-protection. For such interior surface protection, I employ gum copal as a basis, but, the hardness of copal is tem- pered down with much softer, tougher, and more elastic gums as mastic, elemi, sandarac, etc. [Details of this work will be given later.] Again I call your attention to the fact that sounding-board wood, yielding the best tone-quali- ty, is first to be ruined by disintegration from heat and moisture. I know how much of variety exists in the matter of violin-taste. I well remember my own different taste at different periods in my exper- ience. At the earlier periods, I demanded noth- ing so much as great power in violin tone. Quali- ty, in either single-stop tones, or double-stop tones, I willingly sacrificed to mere power of tone. As I now look backwards upon my earlier attempts at making earth to tremble by my great tone, I am staggered at the mountain of infliction enthusiasti- cally poured into the suffering ear of near friends. The word "near" only refers to proximity, and by

82

\section*{VIOLIN TONE-PECULIARITIES.}
no means refers to ' 'loving" friends. Today, when speaking of the sounding-board yielding best tone, I mean the greatest tone-beauty, or, the most agreeable tone, or, the sweetest tone; and I espec- ially mean great beauty of double-stop tones. As a sample of the destruction wrought upon the better grade of sounding-board wood by heat and moisture, I present this old violin. The maker of this violin is not known. Its worn appearance is such as to make needless any certificate of age. The peg holes are worn beyond the size of any vio- lin peg! The "hand" is deeply worn by shifting thumb and finger. On surfaces touched by chin and shoulder, the varnish has disappeared. At both upper and lower extremities of the sounding- board the soft cellular connective tissue, between denser part of grain, is worn away, leaving the denser part standing as ridges. Because of its great influence upon double-stop tones, I call espec- ial attention to the soft nature of connective tissue in this sample of sounding-board wood. By reason of long experience, I feel safe in stating that double-stop tones from this violin possessed marked beauty. [I employ the past tense because this violin is not in playing order.] The only statement in text-books of philosophy, in which I have found something of value to the selection of violin sounding-board wood, is contain- ed in the following sentence: "The sounding- board may be compared to a bundle of strings." I consider this statement to be of value. Evident-

83

\section*{VIOLIN TONE-PECULIARITIES.}
ly, the word "strings" in the text-books refers to the denser part of grain. To be of value for sounding-board use; the "strings" must be held together by a connecting medium. Evidently, the character of such connecting medium modifies ac- tion of those strings. Thus: Freedom in action of those "strings" is as rigidity of connective tissue. As previously shown by example, the rigid violin sounding-board cannot be made to yield agreeable double-stop tones, because rigidity in connective tissue interferes with independent action of con- tiguous fibers; ("strings," of the text books.) In this old sounding-bcaid, connective tissue is soft. The owner of this much worn violin is Mr. Aug- ust Wolfe, Music Director, Valparaiso College, Val- paraiso, Indiana. Prof. Wolfe brought this old violin from his Austrian home. When a small boy, this worn old instrument came to him as a present from his father, and was accompanied by the usual fallacious remark, "It'll do to begin on." [Considering the present condition of this violin, I believe that none other than a German-speaking boy could have survived to become a master

violinist.] In size, this violin is known as i. I ask you to observe the extremely delicate beauty of this old neck and head. Even Hebe's own is not more beautiful. This slender hand, these thin peg-box walls, those exquisite lines of fluting and scroll, powerfully appeal to our sense of the beautiful. Because the peg-holes are worn out, and a fracture extends down through the A peg-hole, and a piece

84

\section*{VIOLIN TONE-PECULIARITIES.}
is gone from edge of fluting, the Prof, suggests re- placing this worn old neck with a new one. I have heretofore done surgical work calling for nerve, but, at severing this beautiful head from its body, I halt. "Pis a case of "heart failure." I can eas- ily give days to the work of restoring the broken lines on this Hebe-like neck and head that those lines of beauty may remain to delight the eye of connoisseur. I am now to present something unusual. I call attention to the varnish on this old violin. This varnish is soft; rather too soft to agree with modern ideas; also, the amount of varnish originally applied, judging from the amount now found in the hollows, was much greater than is employed today. It is so soft that with but moderate pressure, my finger palp causes perceptible indent therein. As friction pro- duces the familiar odor of mastic, therefore I sup- pose this soft gum to be in excess. In this varnish, I am less interested in the gums than in the color- ing matter employed. As you observe, the prevail- ing colors are red and black, and mixed, in quanti- ty of each, to produce a deep red-brown, or brown- red shade. But, such dirty brown-red! No self-respecting modern violin maker would permit such dirty, muddy-looking color-work to appear out of his pile of waste- wood. "Doctor, isn't this a rare old violin?" S-u-r-e! [I gave notice of presenting something unusual.] I am now to present something yet more strik-

85

\section*{VIOLIN TONE-PECULIARITIES.}
ing-; more shocking; and more pathetic. I am to present an example of disintegration of sounding- board wood such as will burn into your memory; an example of slow conbustion of wood in the presence of heat and moisture; an example showing the down-hill route traveled by those priceless old gems of tone-art now stranded along the violin's last hundred-year path. This irreparable loss, and this pathetic scene, might easily and with certainty have been averted, I fully believe. With a thin spatula I easily remove this old sounding-board. Shade of Dante! [As an application in Hades, we may well fear that slow combustion heads the list. The truth in that old saw, "The devil's in the fiddle," now re- ceives confirmation. But, strangely, and unexpect- edly, this confirmatory evidence points to the fact he "is in the fiddle" for the purpose of destroying- the fiddle instead of destroying human beings. It is claimed that recent archaelogical discoveries bring to light proof of some error in Scriptural reading's. The evidence afforded by this example of slow combustion, when supplemented by other examples of like nature, may be sufficient to change orthodox reading's thus: "The devil is an enemy of the violin."] "Punk" may be described as a product of wood by slow combustion in the presence of heat and moisture. Upon the inner surface of the sounding- board we have an example of "punk." It falls as dust mixed with thin threads of wood-fiber as I touch it with scraper. This punk condition extends

86

\section*{VIOLIN TONE-PECULIARITIES.}
over the entire inner surface of the sounding-board, and its depth equals 1-32 inch. In places, the lin- ings on upper edge of the ribs are hanging in par- tially decayed shreds. You observe that both lin- ings and corner blocks are much lighter than those of the present day. Why all samples of sounding-board wood of equal age, do not present similar punk conditions is a matter beyond my ken. I can only say that differ- ent samples of wood possess different degrees of resistance to slow combustion. [In my younger days, punk, gunpowder, flint and steel, were our only matches. We took equal pre- caution in keeping punk and powder dry. Thus kept, punk would last a long time; but,. exposed to heat and moisture, as when found in the woods, it soon turned to dust. Not every decaying log yields a good quality of punk. The punk hunter must often make extended search before finding the best. At this distant day, to the best of recollec- tion, the best quality of punk was found in large, dead branches upon living trees; or within the trunk of trees dying while standing. In woodcraft it is well known that the greatest depth of color- changes are found in trees dying while standing; also, that lumber from trees thus dying sooner de- cays. In seeking for reasons explaining rapidi- ty of disintegration observable in the better grade of old violins, some writers suggest that the sound- boards thereof were obtained from trees having died while standing. Upon this point there ap- pears no positive evidence. The fact that some

87

\section*{VIOLIN TONE-PECULIARITIES.}
samples of sounding-boards have succumbed to the action of heat and moisture, while other samples are yet sound, remains unexplained.] I now hold this violin in such manner that you may look towards the inner surface of the back plate. I am thus particular to say "look towards" because you cannot ''look upon" this plate. Your gaze cannot penetrate through this deposit of earthy matter. The color of this deposit is black, and therefore suggestive of alluvium. But, alluvium being a deposit from water, we must therefore look elswhere for the origin of this deposit of aeri- al dust from the "Cremona period." [Some of you may entertain a different opinion. I admit that, your opinion is equally as good as mine, possibly better than mine. I only claim myself unable to estimate "Cremona dust" above its face value. I know that some old-violin experts, "just out for their health, " (!), claim no other dust has value. I'm not "out for health." I admit that Cremona dust has a penchant for settling within the violin. Indeed, assisting such dust to settle has improved the health, (!), of many experts. Good health is a valuable asset. For this reason alone, the old-violin "expert" undoubtedly will re- main "out."] Were I attempting the old violin "expert" role, I might continue thus: "Were a few boulders to be seen scattered about in this deposit, 'twould in- dicate that the age of this old violin is co-equal with the glacial epoch. But, as boulders are ab- sent, we may therefore safely continue march, in a

88

\section*{VIOLIN TONE-PECULIARITIES.}
backwards. Undoubtedly this deposit was once fly- ing particles of earthly matter dust and nebul- ous. Now we have it. The nebulous period is the limit. Gentlemen : Behold the first violin ! ' ' Having now removed sufficient punk from the sounding-board, and sufficient dirt from the back- plate, we are enabled to observe that both plates are cut from the "slab." This manner of working out violin plates seems to have prevailed to a great- er extent in earlier days of violin history than in modern times. In my own experience, the "whole" plate yields tone equally as good as the divided plate, provided that the grain of sounding- board stands at, or nearly at a right angle to the arching in the tone-producing area. But when such grain angle is oblique to the arching I have always found that tone to suffer loss in both power and brilliance. The same loss also occurs in the divided sounding-board when the grain angle is oblique to arching. I attribute such loss to dimin- ished fiber action. [Honeyman states with positiveness that when- ever a Nicolas Amatus sounding-board is found with thicknesses less than i, at center, and i at edges, such sounding-board has been subjected to re-graduation by some modern workman. For the sake of both Music and humanity, I consider such re-graduation to be fortunate. It is my experience with sounding-boards thus thick, and of average density, that even a set of guage 1 strings cannot overcome rigidity sufficiently to yield one octave of

89

\section*{VIOLIN TONE-PECULIARITIES.}
good tone upon each string.] As evidence bearing- upon the confusion found in violin history, I here quote from the Frenchman, N. E. Simoutre, book and charts, Paris, 1885. Simoutre submits charts of two Nicolas Amati violins without date: Thick-

ness, Plate 1, center of sounding-board, 360mm, (meaning three and five-tenths millimeters) down to 250mm, and at the edges, 400mm, (meaning 4 millimeters.) Thickness, Plate 2, at center 450mm, down to 200mm, and at the edges, 300wm. (Simoutre explains his figures by centieme des millimetres, meaning so many hundredths of a millimetre. )

Comparing the thickness at the center of the Nicolas Amatus sounding-board given by Honey- man with the greatest thickness given by Simoutre, we find a difference of seven-hundredths inch. Considering the depths of human affection for the violin, I am astounded at its abuse. Were this old violin once a masterpiece of tone-art, the abuse to which it has been subjected is ample for its ruin. Repairing this "worn old thing," in a way to extend its period of usefullness, is a severe tax upon patience. Brittle as glass, it demands careful handling. In this work, a money consideration cuts but little figure. Reward comes in listening to such

sweet tones as money cannot reproduce:

90

VIOLIN TONE-PECULIARITIES. Sweet old violin! Worn, Torn, scraped from e'n till morn! Work for thee has ever been To hand out joy since thou wert born. Dotage, now within thy form, Marks thee down a "worn old thing;" Yet, thy tone to us remains A gem beyond the reach of king.

91

\section*{VIOLIN TONE-PECULIARITIES.}
LECTURE VI. GENTLEMEN: Since our last session I have com- pleted repair-work, re-touching work, and now pre- sent this worn old violin in playing order. As you observe, the marks of age and wear are much less in evidence. From your seats, the patch-work up- on head and peg-box cannot be seen. The worn areas upon sounding-board, back, and ribs are re- colored, and the entire exterior is re-finished. In- deed, this old violin no longer appears as a "worth- less old thing." But, the exterior work, compared with the interior work, is a trifling matter. At the time of building this violin, its exterior surfaces were given a quality of varnish-protection sufficient- ly durable to last to the end of time. Gum mast- ic, while long in drying, and never drying hard, is the most elastic, and the toughest gum coming un- der my observation. In my hands, gum mastic, with oil required the time of one year for drying; and even after ten years, it is not dried hard. [As you know, violin varnish is a subject arousing end- less discussion. I do not intend to arouse discus- sion upon this point. I shall only state facts as they appear to me. It is a fact in my observation that any varnish whatever, which dries hard, proves itself to be a serious damage to violin tone. It is within my experience that varnish interferes with independent action of sounding-board fibers exactly in proportion to its rigidity when dry. As an experiment, many times have I "tied up" tone by application of rigid-drying gums to the sound- ing board. Equally as many times have I ' 'untied' '

92

\section*{VIOLIN TONE-PECULIARITIES.}
tone by removal of such varnish. I do not find that application of rigid-drying gums to the back plate and ribs results in any damage whatever to tone. I have known violin users who prefer the "tied up" tone. To satisfy the taste of such vio- lin users is an easy matter, and of certainty in at- tainment. Such certainty is due to the fact that varnish action, or rather, lack of action, is a con- stant factor. Because the action of wood is not a constant factor, therefore, the amount of varnish necessary to "tie up" the tone of any given sound- ing-board can only be determined by trial. It is a fact that in playing upon a violin having the "tied up" tone, less care is necessary in handling the bow. As a sequence, one may therefore more eas- ily pass as a skillful violin player. Of course such tone fails in the long-distance test.] The soft varnish upon this old violin, notwith- standing its age and use, would be to-day in per- fect condition had it but received proper care. As an example of correct care for the violin I know of none surpassing that of the renowned violinist, Bernhard Listemann. That Listemann is a gentle- man of the old school needs no more than an in- troduction. That he is an artist the world knows. That he is a connoisseur, and collector of old Italian violins, is in evidence at the moment he opens his fire-proof vault. As he lifts one of those violins, his care is in evidence thus: His left is holding the neck; his right thumb and index finger is holding the tail-pin. As he presents it to you, he politely says: "Please sir, hold it thus, and do not touch

93

\section*{VIOLIN TONE-PECULIARITIES.}
the varnish." Upon his violins you do not see a speck of dirt, nor a speck of rosin, nor a scratch from button or finger-nail, nor a greasy imprint from the finger-tips. Such is proper care of the violin. With such, care continued, the brick and stone walls of the Listemann mansion will be flying dust ere the exterior of those violins shows sign of disintegration. Now I'm not, in the least degree, going into ravings over Cremona varnish, nor Cre- mona colors in Cremona varnish. I never rave without at least 50 per cent, of absolute in mine. Thus when I do rave, it is for "good" reason. Again, because of peculiarities in my optic nerves, (for which I'm not in blame) "aerial dust" from the Cremona period does not affect my eyesight. Therefore, clear gum copal, tempered down with clear gums of softer nature, appear just the same to me whether lying upon the Maggini, the Guarneri, the Amati, the Stradivari, the Montagnani, the Francisco Ruggeri, or lying upon the Franklino Robinsoni, (Frank Robinson), but, in the matter of using colors, some of those old violin builders, (not all, ) did possess a rather unusual development of color-sense. Because color-sense is usually a matter of growth by cultivation it therefore fol- lows that this sense will not be possessed by all in an equal degree. Although all of the older makers might have employed the same gums, yet, in color- ing such gums, widest results might be anticipated. It is evident that the person who colored the varnish applied upon this old violin was largely de- void of color sense'. That his selection of gum was

94

\section*{VIOLIN TONE-PECULIARITIES.}
excellent is in evidence by the resistance to wear while yet being soft. (That any one could for a moment, entertain be- lief that gums can be made to permeate wood, upon which they are laid as varnish, is to me something astounding. Even were permeation by gums a possibility, nothing but ruin of sonority could fol- low. Such ruin is easy of demonstration. Many times have I ruined sonority of the sounding-board by applying thereon a material which does per- meate through the wood. After every such soak- ing, by whatever agent, the tone thereafter re- mains dead. Several violins, ruined by soaking the sounding-board have been brought to my notice. For their dead tone I know no remedy other than a new sounding-board. Work upon the interior ) surfaces of this old violin is of vastly greater im- portance to its tone. Those slender linings, and light corner-blocks are replaced by others having more than two times greater mass. That solidity of linings and corner-blocks adds to tone-power is also easy of demonstration. By removal of the linings and corner-blocks from any violin of good tone-power, the proof thus obtained affords ample evidence. To the great original thickness of this sounding-board may we attribute remaining tone value of this violin. Had its original thickness been reduced to equal thickness with the Strad of 1707, (given by Simoutre as 2 and 8-10, down to 2 and 7-10 mm, this sounding-board would be to-day ) in a hopeless condition. The amount of disintegra- tion from those destructive forces, heat and mois-

95

\section*{VIOLIN TONE-PECULIARITIES.}
ture, would now be sufficient for total ruin. Its re-graduation is copied after the Strad of 1707. First, I will give you opportunity to judge of its tone, and thereafter describe the grain and color within the sounding-board. As you observe at nearby distances, the power of its tone is not great; but, its other tone-qualities cannot be excelled. Particularly observe the power of harmonics. As you have observed, power of harmonic tones, from different violins, is a varying quantity. From dense sounding-board wood I have not been success- ful in securing great harmonic tone-power. Har- monics are of value to the solo violin. Indeed, than harmonic overtones and harmonics a basset, as I like to call them, resultant tones, as text-books call them, there cannot possibly be more beautiful sound. As previously shown, to these beautiful sounds must be credited the "rich" violin tone. In my experience, audible resultant-tones are diffi- cult of production. So difficult of production are they that, from many violins, I have not been able to produce them at all. Only from sounding-boards yielding absolute purity of musical sound have I produced them in marked degree. In all cases where those shadowy creations exist, I believe the tone of such violins to have reached the limit of tone-value. To make audible resultant tone requires two other tones in -exact chord. Exactness in chord of the generic tones is wherein lies the difficulty. At the least perceptible variation from exactness, those filnay sounds in- stantly disappear, and no coaxing whatever can

96

\section*{VIOLIN TONE-PECULIARITIES.}
induce them to re-appear until after exactness in the generic tones is established. (Text-books state that, as a rule, resultant tones are pitched two octaves below the generic chord; but, exceptions to such rule are also noted. Thus: The resultant tone, produced by the third and tonic above, is not two octaves below either third or tonic, but, is two octaves below the fifth of that particular scale. As an example I draw out the tones B, and its tonic g above the staff; therefore the scale is G major, and the fifth therefore is D. In this case the resultant tone is two octaves below D; and, to find its pitch, we divide the pitch of D, 600 by 4, equals 150. The harmonic overtones of B and g are respectively, one octave above. Thus: B equals 500; its harmonic, 1000; g equals 800; its harmonic 1600. Again I call your attention to the fact that the combination of such five tones of widely vary- ing pitch is the cause for the "rich" violin tone.) Yes, in truth, I have carefully scrutinized those sounding-boards yielding the "rich" tone. I have even re-opened such violins for no other purpose than to re-examine the sounding-board. As you know, my work upon the violin as been chiefly given to such as have been in use; therefore the ''finish" prevents accurate determination of physi- cal qualities from examination of outside surfaces. To the best of my ability, I will now describe such physical qualities and the color thereof. First: The grain follows a straight line, with no deviation whatever, nor wherever. Second: The yearly growths are neither widest nor narrowest: and

97

\section*{VIOLIN TONE-PECULIARITIES.}
measure on an average, 18 to the inch. (Extremes are 16 to 20.) Third: There is no appearance of fat whatever. Fourth: There is no sap-wood, nor black- wood spot. Fifth: The wood is brittle. Sixth: The whole grain is rather soft than dense. Seventh: A keen, smooth cutting tool is necessary to leave a smooth surface. Eighth: The color is the deepest yellow of butter with more red than in any sample of butter. Ninth: This depth of color extends through the sounding-board. From all evidences in wood-craft that I am able to muster, such sounding-boards only come from the older and larger trees. Both plates, and linings, and blocks of this old violin are now hermetically sealed. It is therefore expected that never again will those terrifically de- structive agents, heat and moisture, show their presence within this violin. Therefore, whatever of tone value it now possesses, it will possess cen- turies hence. That its tone is beautiful, after such interior protection, yourselves may verify. You may also verify to the beautiful tone of 59 other used violins having similar interior protection. Someone has stated that the average longevity of the used violin is 80 years. Upon this point the difficulty in securing accurate statistics is apparent. To the violin student, the average period of violin usefulness is a topic possessing much interest. As a premise, I will state that enthusiasm, plus power of bow-arm are potentialities affecting violin longevity not to be omitted from consideration. Thus when stating that I wore out a good violin

98

\section*{VIOLIN TONE-PECULIARITIES.}
sounding-board in 30 years you will first look at my eyes to see whether I mean it or not. Next, you will undoubtedly look at my 200 pounds of human- ity and mentally "size up" my 16-inch biceps, plus enthusiasm. Both will stand inspection. It seems to me that because my violin has noth- ing of the romantic, nor mysterious, nor spectacu- lar about its origin, therefore its history must be refreshing. Certainly 'twill be unique. As an im- portation, my violin came from Luxembourg. In this fact there's nothing extraordinary. Haller, himself a Bohemian, said my violin was undoubted- ly made by a herder of sheep on Tyrolean moun- tain. I say good for the sheep, good for the mountain, better for the herder of the sheep, and best of all for me, for that violin possessed a tone. That's why I wore that violin out. That's why Haller, himself my violin tutor, thrice counted out \$100 for it. That's why I wouldn't sell it. That's why I think all good violin makers don't hail from Cremona. That violin had tone-power to give away. But, its value lay not in tone-power. Its tone-intensity was marked to a rare degree. Its value was not in intensity. It was both brilliant and sweet. Not there its value. Other violins possessed power, intensity, brilliance, and sweet- ness of tone, but compared with my violin, were valueless to me. I am now come to something which is extraordinary. My violin possessed hu- man-like tone-quality. Its tone could weep in des- pairing sorrow, in contrition, in joy; it could pray in deepest devoutness; it could laugh in utmost

99

\section*{VIOLIN TONE-PECULIARITIES.}
abandon; it could sing sweetly as a bird. In these tone-qualities lay its value. Sell my violin? Starvation first! Even then I'd have worked the "hand outs" to the limit. In thirty years use that powerful, intense, bril- liant, sympathetic, human tone went down into piano. At first I did not discover the reason. I attributed the reason to first one thing and then another. I tried strings of different sizes; finger boards at different heights, of different densities, of different weight; bridges of different mass, of different density; posts, ditto; but all were of no avail. Hitherto in moments of triumph, my violin had always brightly looked up into my eyes for approval. I always granted all it asked of me. But now I noticed a change in its look. In place of its wonted brightness, there was sadness. Pity- ingly it looked up at me now. The deep despair on its sweet face brought enlightenment to me. Upon discovery that the caresses of my bow were crush- ing the life out of my pet, the pain at my heart was something I would like to forget. Gently I laid it in its case nor opened that case in two years. I couldn't. I am but an ordinary violinist. I never could play anything behind nor on top of the bridge, nor but little in front of it. With that violin I didn't have to do much myself. It did best when I med- dled least with its moods. Often I gave way to

100

\section*{VIOLIN TONE-PECULIARITIES.}
those moods. Then it was in delight. Then it would lead me a chase through the shadows and sunshine of melody in a way at once my despair of representation by notes. Years have rushed past. To-day memory rushes me back to days when I lived in the little city of C. In a large majority, its people came from the thither Atlantic shore. Each of such had exper- ienced the heart-ache in family partings. Each of such, loving liberty, firm in purpose, buoyed by hope, had left loved ones and set out for that bright western star whose symbol is known to the world as U. S. A. Each of such, knowing, the rustling, jostling, pushing, human activity, the crushing, whelming, human pain at departure, ask- ed my violin to paint that scene abandoned by art- ist-brush in lack of color for the pain of brother parting from brother, parting from father, parting from weeping, dear old mother; for the pain of sister, alone, braving the sea for a home with the free; for the pain of father his family leaving, for the sound of gong, for father's tender tone in "good bye, son," "good bye, daughter," for his heart-aching tone to wife, "Mother" he calls her, (gong), for the silence as his arms enfold her, "Mother," 'tis but a whisper, gone her voice, (last gong), for the smothered sob from father's pain-heaving chest, for his unsteady step, as blindly he follows the sound, for the thud of "Mother's" falling, for the hissing of steam, for ''all aboard," for the fainting shout of fading friend, for the crashing of pounding billow, for the

101

\section*{VIOLIN TONE-PECULIARITIES.}
storm demon's maddening shriek, calling for hu- man victims daring his fury, for the quiet of re- turning calm, for the relief of human tension in 3-4 rythm, for the hum of returning animation in rollicking 2-4, for the joyful shout, "Land, ho!"- yet, my violin accepted such invitation. I loved that violin. Poor Patti! She, the one bright star of us elder- ly folk, at 65 weeping weeping because her once thrilling tone has gone down into dotage! In public I have not wept. In private, the handkerchief in my right is yet convenient. Strange how this idol of wood can find its way down into deep recesses of human heart, thence defying all comers! Occasionally, in unexpected moments, and in un- expected places, one reads fragmentary statements about this thing, or that thing, able to restore vio- lin tone. After the tone of my violin had gone down into dotage, one such statement attracted my attention. The writer, without attaching signa- ture, stated that oil varnish, applied to worn out violins, had power to restore tone. At that time I neither could affirm nor deny such statement be- cause of having no experience with oil varnish. Although this statement bore little of sound logic, yet, being something not tried upon my violin, I therefore decided upon giving oil varnish a trial. From a friend, I procured a quantity of such varnish, together with instructions in the manner of its application. With strings and bridge in po-

102

\section*{VIOLIN TONE-PECULIARITIES.}
sition, I began the work of applying oil varnish by the "rubbing" process. As this method of apply- ing varnish was new to me, I could not determine how much varnish I was using, but thinking if there's good in little, there's good in more, I gave a whole day to rubbing it on. (At this moment, I think of the truth in, "How little we know without experience. ' ' I also con- fess to an attack of "quick return" fever. You know what "quick return" fever is. It is that "can't wait" feeling. Humanity in general may be attacked with this fever; but, its most virulent form is manifested upon such as handle the violin. ) Of course, at the moment of completing this days work, I was filled with anxiety, buoyed by hope, and pulled down by doubt. Hope in the pos- sibility of restoring tone to my de-throned violin made me over-anxious for returns. Therefore, I threw down the rubbing pad, picked up a bow, and drew from the G an octave of the "deadest" tone imaginable. (Not many of us go through life without seeing ghosts; at least we thought so at the time, which amounts to enough for a story. The point is, the ways of ghosts are sudden. Their coming is never announced by shrieking whistle nor clanging bell, nor grinding wheel, nor sounding horn. Suddenly they come. Suddenly we go; that is, so soon as our suddenly vanishing breath permits of going. The suddenness of ghostly appearance is not more startling than the phenomenon appearing upon the sounding-board of my violin as I removed it from

103

\section*{.VIOLIN TONE-PECULIARITIES.}
beneath my chin. ) This phenomenon consisted in a large number of crater-like openings, or upheavals, in the oil var- nish I had just put on. Some of those openings were 3-16 in diameter, and running from that di- ameter down to mere points. The crater-like open- ings were so near together that the edge of each touched that of its neighbor. In this varnish phe- nomenon, the point of greatest interest is its locat- ion. Doubtless you are familiar with the theory of the French philosopher, Savart, upon the quest- ion, "How the violin operates to produce sound." You remember those elaborate, scientific experi- ments, and his conclusion therefrom; that the top- plate, back-plate, and ribs, all join at once in strik- ing those blows upon contained air resulting in the production of sound? (From N. E. Simoutre's book, 1885, I infer that Savart is yet regarded in France, as an auth- ority upon this question. From other sources, I learn that Savart 's theories are considered as ex- ploded. Up to the moment of this varnish phe- nomenon, I was a follower of Savart. His drafts on science looked genuine, and, in no guide book could I find either affirmation or denial of his posit- ion. I therefore fell in with his conclusions with- out attempting any demonstration myself. Really, after German and English philosophers dismissed violin tone with sui generis, (self-generating,) thereby practically abandoning the problem as unsolvable, I had no hope of ever finding, or see- ing such solution. I do not claim such solution to

104

\section*{VIOLIN TONE-PECULIARITIES.}
exist now except as a partial solution. From my point of view, the evidence afforded by this phenomenon, together with other demonstrations, to be given later, are ample proof of error in Savart's theory of how the violin operates to produce sound. ) There is not a shadow of doubt about the cause of this varnish disturbance. The upheaval of this soft mass is wholly due to oscillations, or vibratory movements, (synonymous terms, ) of the sounding- board. I applied the bow only upon the G string. The point is, can you, or I, or anyone whatever, pre-determine the location of this disturbance by any existing theory of how the violin operates to produce sound? It is at once self-evident that such disturbance must exist over the entire violin body, did the entire body act with equal energy in pro- ducing sound. The varnish I applied was equally distributed over the entire body. The area of disturbance is limited. In length, this area equals 2 and inches; its width, 1 and i inches. The lar- ger craters are at the center of such area. From the center, the size of the craters diminishes down to mere dots. It is self-evident that the widest sounding-board oscillation occurred directly be- neath the wider craters. It is also self evident that this point of greatest sounding-board oscillation is the center of a certain area responsible for G-string tone. I call attention to the fact that the G-tone of this violin was once noted for its power. I call attention to the fact that the locat- ion of this area of varnish disturbance affords a

105

\section*{VIOLIN TONE-PECULIARITIES.}
valuable cue in tone-regulation of the violin; also, to the fact that the finding of the area responsible for G-string tone led to finding the areas responsi- ble for D, A, and E-string tone; also, to my conclu- sion that the violin sounding-board is wholly re- sponsible for violin tone. I mean that to the blows delivered by the sounding-board upon contained air do I now attribute violin tone. (This statement does not include violin tone- modifiers; as the strong back-plate; the weak back- plate; imperfect inner surface of the back plate; position and area of exits; model of violin; position, length, density and diameter of post; position, den- sity, height, and mass of bridge; height weight, and under surface of finger-board; diameter and quality of strings; amount and quality of varnish. These modifiers of tone will be considered at a later date.) I will now consider the location of varnish dis- turbance. The center of this area is half-way from the position or the bridge to upper edge of the sounding-board, and I inch to the left of the bar. Possibly none will view the location of this area of varnish disturbance as I view it. As I view it, this varnish phenomena, and its location, turned a flood of light directly upon the sounding-board areas responsible for tone of each violin string. During more than ten years of continuous work, since occurrence of this accident, (I call it accident No. 2.) I have used this index in tone-regulation of used sounding-boards without once meeting dis- appointment. Therefore am I convinced that this

106

\section*{VIOLIN TONE-PECULIARITIES.}
index is reliable. From this index I learned where to find, and how to correct sounding-board errors injurious to tone of all strings, or only injurious to any single string. Therefore I consider this index of exceeding value to the violin. I also acknow- ledge this index to be the triumph of accident. Science has no part whatever in its existence. Sci- ence can only come along with a post facto explana- tion. It is my belief yet, that science cannot build a violin. I mean that by no scientific formula can you, nor I, nor any person whatever, pre-determine violin tone throughout the list of violin tone-pecul- iarities. I call attention to the fact that this index great- ly simplifies the work of tone-regulation; also, to the fact that such regulation may be directed wholly to, and upon, the sounding-board. In this matter, the evidence afforded by accident No. 2 may be accepted as corroboration of the facts many years known to violin students. In the last ten years of my work, the back-plate has been wholly ignored as a tone-producing agent. I have only treated the back as a tone-modifier during the time mentioned. All I now ask of the back-plate is that its inner surface be absolutely perfect for the re- flection of sound-waves, and that its rigidity be suf- ficient to withstand, without a tremor, the charge of molecular movement originating at the sounding- board. I am aware of the fact that many good vi- olin makers continue to treat the back-plate as a tone-producing agent. Such violin makers may continue to produce good violins.

107

\section*{VIOLIN TONE-PECULIARITIES.}
(This question seems to possess a tenacity to life equal to the tenacity in the question of planting potatoes. The enterprising Circassian, although finding the American aborigine enjoying potatoes, failed to find traditional authority pertaining to the proper phase of the moon for planting potatoes. Such omission entailed a seemingly perpetual di- vision in Circassian ranks. Yet, strangely, both sides to the controversy grow good potatoes. ) As I view the location of this varnish disturbance, it also affords a demonstration for the power of sympathetic action. Briefly, as a law in physics, it may be stated that the power of sympathetic action diminishes inversely as the square of the distance. Therefore, that part of the sounding- board nearest the string must receive greater im- petus from string-action than the part farthest from the strings. Therefore we might pre-sup- pose the greater varnish disturbance to be located upon the upper-half of the sounding-board. As I held the violin beneath my chin I could plainly feel sounding-board vibration at that point, as was us- ual. Therefore, I know that the sounding-board acted in nearly its entire length at the moment of causing the varnish disturbance. But, that such action was of less power in the lower-half of the sounding-board is proven by the fact that no var- nish disturbance whatever appeared on the lower- half. (Within the limits of my observation, sym- pathetic action, excited in the sounding-board by string-action, has not heretofore received atten-

108

\section*{VIOLIN TONE-PECULIARITIES.}
tion. For that reason I am desirous of your opin- ion thereon. Indeed, in giving my conclusions in this matter, I expect to be favored by the opinions of other students of violin tone. Because this new question is brought before us by accident, I there- fore entertain no feelings of proprietorship. It is yours as much as mine. Nothing but the presen- tation of the question happens to belong to me. ) I do not find the laws governing sympathetic ac- tion to be complex. On the contrary, such laws are easy of comprehension. Up to this moment I can call up only two facts as bases for such laws: (1) Equal susceptibility to force. (2) Proximity of bodies. That equal susceptibility to force is a con- dition necessary to sympathetic action may be demonstrated by attempting to excite such action in bodies possessing widely varying mass, density, and rigidity. Such attempts are failures because such widely varying bodies are not excited to vibra- tory action by an identical force. As a ready-to- hand means for demonstrating unequal susceptibil- ity to force, I employ this violin. As you observe, each string on this violin yields powerful tone. These strings are carefully selected guage 2. The rigidity of sounding-board beneath each string has been carefully reduced until the force in the blow of each string, at concert pitch, is ample for pro- duction of vibratory action in the sounding-board and corresponding to the action of each string. I remove this G and replace it with this E. In its new position the E-tone is weak. Now, the sound-

109

\section*{VIOLIN TONE-PECULIARITIES.}
ing-board beneath the E is not susceptible to iden- tical force with the E. Hence there is no sym- pathetic action between these two bodies. In its proper position, the tone from this E is powerful; and such tone-power is due to the fact that the E, and the sounding-board beneath the E, are suscep- tible to an identical force. Hence sympathetic action exists, and its existence augments tone- power. (Often in re-toning work I demonstrated the value in "Let well enough alone." Thus, after deter- mining length of strings to be used, after deter- ming length by position of bridge, after adjusting length, mass, and position of post, after determin- ing mass and height of bridge, after testing the tone and finding slight weaknesses here and there, and positively knowing such weakness to be due to too much sounding-board wood here and there, after reducing thickness in such places with my utmost precaution, after securing even power for all strings, then idiotic-like, instead V)f letting well enough alone, I have lost my work by trying to do better than "well enough." The ' 'Elgin," adjusted to temperature, position, isochronism, and running "well enough," is not more susceptible to idiotic treatment than the finely adjusted violin. With not more safety can we file off metal at any point on the periphery of the Elgin balance wheel, than we can change the pitch of strings, diameter of strings, length of strings, height of bridge, mass of bridge, density of bridge, length of post, mass of post, position of post, and sounding-board rigid-

110

\section*{VIOLIN TONE-PECULIARITIES.}
ity of the ..finely adjusted, tone-regulated violin. "Tis no use to cry about spoiling my work. Crying but adds "baby" to ''idiot." With expressions more or less colored, I jump up, seize the red pen- cil and write upon the wall, "Let Well Enough Alone.") The familiar, fun-provoking broom-stick fiddle is not without a lesson of value. Although its one vipjin D-string arouses no '.'sympathy" neither from the audience nor from the broom-stick, yet, it has an opportunity to demonstrate its tone-pro- ducing power when totally unaided by sympathetic action. ' K '\textasciicircum\{\}'t-. Proximity, of bodies, .as a basis for law gdvern- ing sympatKetic action,, may be demonstrated in a variety of ways. As an imperfect way, and only because of Convenience, I remove all strings upon this violin to a greater distance from the sounding- board by replacing thjs bridge with another of an inch higher. As you. observe,, tone-power of all strings is perceptibly diminished. Although faulty, yet, this demonstration shows that sympathetic action diminishes as distance increases. It is evident that distance may annihiliate sym- pathetic action. Per contra, proximity augments sympathetic action. From these facts I conclude that sympathetic action causes wider amplitude of oscillation in the upper-half of the sounding-board; and, to such wider oscillation, with its augmented power, do I attribute the location of this varnish disturbance. From the same facts, I also conclude that the upper-half of the sounding-board,

111

VIOLIN TONE-PECULIARITIES. to a limited distance either side of the bar, is re- sponsible for G-string tone-power. From the same facts do I conclude that rigidity should be less in the lower-half than in the upper-half of the sound- ing-board. Distance certainly diminishes the force of sympathetic action as the square of the distance. (Knowing that every sentence in the description of this varnish phenomenon, together with my conclusions therefrom, will be subjected to the limits of scrutiny, I have therefore called upon my utmost ability for clearness of diction. Should you find lack of clearness upon any point, and should you desire further elucidation thereon, you need only to notify me. In advance, I request your opinions on my conclusions. I admit the fact that "two heads are better than one." I admit that your opinions may be better than mine; therefore your opinions will be received with pleasure. ) I have now completed the presentation of varn- ish phenomenon No. 1. Varnish phenomenon No. 2 will be presented upon a later occasion.

112

\section*{VIOLIN TONE-PECULIARITIES.}
LECTURE VII. GENTLEMEN: At this hour I present varnish phenomenon No. 2. As you remember var- nish phenomenon No. 1 occurred on the exterior surface of my violin. Varnish phenomenon No. 2 occurred upon the interior surface of another vio- lin. The latter violin is the first in a list of sixty wherein I applied an interior surface protection, and the latter phenomenon is due to two errors: (1) Crudity in application of material. (2) Assem- bling and testing tone before the material becomes dry. [Thinking that some reader may wish to do in- terior surface work upon the used violin, and de- siring that such person may not meet with disap- pointment, I am therefore careful to note my own mistakes. It is quite safe to assume that, in all lines of human activity, mistakes are made. In knowing . such mistakes, they may be avoided. Aft- er ten years of effort in protecting interior violin surfaces from disintegration and thinking that I succeeded in affecting such protection without in- jury to tone, yet, I do not know or claim, that the details of my surfacing work cannot be improved. I shall be heartily glad to know that my own efforts have received improvement.] Crudity in application of material is shown thus: (a) By failure in preparation of wood-surface, (b) By application of an excess of material. After my experience in this interior surfacing work after observing the intensity and brilliance of tone with- out loss of sweetness whatever, I find myself won-

113

\section*{VIOLIN TONE-PECULIARITIES.}
dering why I did not sooner understand that inter- ior surfaces of the violin ought to be equally per- fect with interior surfaces of the flute, clarinet and horn. In the latter instruments, intensity and brilliance of tone are modified, for better or worse, by the degree of interior surface perfection. I am now fully satisfied that the same condition affects violin tone. In every method of applying trans- parent finish to wood, swelling of grain is a diffi- culty to be met. Coming directly to the point, I know of nothing causing less swelling of grain than boiled linseed oil. But, I find there is a right and wrong way to apply oil. I do not find that oil alone, or in mixture with the slowest drying gums, pene- trates even the softest pine or white cedar when applied in such attenuated layers as is possible to the "rubbing process." I find that wood surfaces must be carefully prepared in order to secure the greatest attenuation of oil, either alone or in mix- tures. In used violins I have never found interior surfaces ready for the reception of finishing mater- ial. The interior surface of the used sounding- board usually presents a series of ridges and val- leys caused by the swelling of the connective tis- sues between the denser fibers. The depth and heighth of these ridges and valleys varies with the amount of water vapor absorbed. As previously shown, different samples of sounding-board wood absorb different amounts of moisture owing to variations in the caliber of sap carrying capiDary tubes. The interior surface of the hard wood back and ribs very often present uneveness from un-

114

\section*{VIOLIN TONE-PECULIARITIES.}
equal shrinkage in the transverse markings; the denser part standing as a low ridge, the tubular part as a shallow depression. The point is to level down all such ridges and remove all moisture be- fore applying any protecting material whatever. I find artificial heat quite often a necessity for re- moval of moisture from used violin wood; also a necessity for continuing the employment of artific- ial heat until after surface work and the applica- tion of finishing material becomes accomplished facts. For leveling down ridges, I prefer to use sand- paper over a block of wood rather than held over the fingers. The finger palps, being soft will cushion more or less, and thus continue cutting away valley surfaces, whereas, the block, being rigid, only cuts away the ridges. For such block H I employ any firm wood, i inch thick, inches in width and 2i in length, having one edge curved and slightly oval. Upon the sounding-board, the ridges are leveled down by -working across the transverse markings. Sometimes \textasciicircum\{\} this work upon the back plate requires considerable time to cut down the dense ridges. But, in my hands, the scraper is a dangerous tool for such work- because of tearing out pieces of wood, and especially dan- gerous in old,' brittle wood. For such reasons -I pa- tiently continue with the sandpaper until the brown surface of the valleys appears new. I next remove sandpaper marks with powdered pumice and pol- ishing pad of felt having one surface slightly moistened with boiled oil, only using sufficient oil

115

\section*{VIOLIN TONE-PECULIARITIES.}
to hold pumice to its work. The final preparatory work is done with the "bone. " The "bone" is a piece of hard ebony, of a size convenient to hold, having one end, or edge, slightly curved and oval, and polished to the limit. With this 4 'bone' ' the sur- faces are rubbed until they are of a mirror-like smoothness, whereafter they are ready to receive the protecting material. Upon surfaces thus pre- pared, oil, alone, or in mixture, may be rubbed on in coats of the greatest possible attenuation, and with imperceptible swelling of grain. The inner surface of the ribs, the blocks, and the linings, receive the same careful attention. In this work there are two important objects to keep in view. First to only apply sufficient mater- ial to cover the wood. Second: To produce a per- fectly smooth surface. For this work, the brush cannot secure the attenuation of the material that is easily secured by the rubbing process. Fluidity of protective material is something to avoid, espec- ially in the first coat. The dry wood will absorb moisture from the fluid applied by the brush, thus defeating the object of hermetically sealing up DRY WOOD. The case is different from that of ap- plying protection to the outside surface of the vio- lin plates. Should absorption of fluid occur at the outside surface, it may escape from the interior surface, and no precaution against admission of moisture into the wood can be too great. When hermetically sealed, moisture can neither enter nor escape from the wood. Hence the necessity for doing this work in a dry atmosphere and the

116

VIOLIN TONE-PECULIARITIES. employment of a material containing the minimum of fluid. I know of no protecting material afford- ing less of moisture absorption than boiled linseed oil. That this oil does not penetrate the softest wood, when applied by the rubbing process, and in attenuated coats, will be shown at a later moment. But, after observing that when a mixture of oil and gum copal, tempered down with softer gums, is employed, the oil, drying with less rapidity, and being largely forced out of the mixture, and lying upon the surface, I have since employed this mix- ture for interior surface-protection to the exclus- ion of all other agents. With the oil thus forced out, I have tried two methods for its disposal. First: Allowing it to remain and dry as part of the finish. Second: Removing it with a soft cloth at the expiration of twelve hours after application. The gums, then being semi-solid are not roughen- ed by careful removal of oil. The latter method has an advantage in greater smoothness of surface, and greater rapidity in drying. With surface oil removed, the gum will dry in 24 hours. With sur- face oil remaining, time required for drying will vary as to the amount of oil, from 4 to 7 days. For application of protecting material, the de- tails are as follows: In my experience, heavy, fine, long-knap cotton flannel has proven the best ma- terial for the rubbing pad. I cut out a piece 2x4 inches and fold it once, end to end and knap out- side. To insure the use of but a small quantity each of varnish and oil at one time, I resort to the following means: One 2 oz. phial is filled two-

117

\section*{VIOLIN TONE-PECULIARITIES.}
thirds with varnish; one 2 oz. phial is filled two- thirds with boiled oil. Slightly moisten one sur- face of the pad with oil, (allowing no swimming whatever of either oil or varnish,) then hold the center of pad over the mouth of the varnish phial, invert the latter, then rub the pad on the wood with a circular motion to avoid sticking and conse- quent roughness of the varnishing surface. The pad should be firmly held and with the end of the index finger pressing down upon its center. When this amount of material has become spread to its limits, then return the phials for more oil and var- nish, thus continuing until the surface of the plate is coated. The amount of material used for one coat may appear surprisingly small. It is small. It ought to be small. It ought also to be small up- on the exterior surface of the violin. In either case I only use sufficient material to cover the wood. So far as I am aware, this method of application commands the minimum of material and the maxi- mum of surface perfection. When carefully ap- plied, and the last coat becomes dry, then do final polishing with dry, powdered pumice stone, using a fresh pad of the same cotton flannel. Of course ex- perience is necessary for skillfulness in this work. I have placed correct details before erroneous de- tails for the purpose of making such errors stand out in glaring colors. To crudity in application of material is due the failure at my first attempt to protect interior violin surfaces; and to the same cause is due varnish phenomenon No. 2 as will now appear.

118

\section*{VIOLIN TONE-PECULIARITIES.}
Varnish phenomenon No. 2 contributes nothing to the value of the violin, and nothing of interest to the violin student further than a description of erroneous details in the work of interior surface protection. Knowing the amount of prejudice ex- isting towards protection of any kind whatever, for interior surfaces of the violin, I am therefore care- ful to minutely describe both the correct and in- correct methods of applying such protection. Af- ter ten years of experience in this work, I am fully convinced that existing prejudice against violin in- terior surface protection is due to erroneous details in application of protection material. As an argu- ment against protection, I remember hearing that it "sharps" tone. I now offer evidence that such protection may be applied in a manner causing "dead" tone. [Than the violin, I know of no historical subject offering greater confusion of evidence. For some occult reason, the violin neither seems nor sounds alike to two different persons. This fact is a phe- nomenon without explanation.] Rated upon tone value, the violin, in which var- nish phenomenon No. 2 occurred was worth \$25. This violin had been used five years at the time I selected it for application of interior-surface pro- tection. For access to its interior surfaces, the sounding board was removed. No preparation other than brushing off the dust, was given to those surfaces. Upon them I poured a quantity of boil- ed oil and transparent, tempered copal varnish. I spread these materials about with a mop. As is

119

\section*{VIOLIN TONE-PECULIARITIES.}
well known in housekeeping physics, the mop is caused to act by reciprocal motion; circular motion being employed only upon parties tracking the floor." Not having to contend with the latter dif- ficulty, I brought my mop into action by reciprocal motion. Of course the mop "stuck" at each re- verse of motion, while I "stuck" to my job by pour- ing on more oil and varnish. After "mixing it" thus until the mass looked as if it were "laid on" (oh!) I pronounced those surfaces amply protect- ed. My glue being hot, assembling work was quickly accomplished. Two hours thereafter, ' 'quick return fever germs had driven me to insan- ity. That 'can't wait" feeling soon became mast- erI applied the bow. [Sheridan's ride to Winchester his changing defeat into victory, has been pictured by able his- torianseven poets have tuned their lyres to the tatoo of his horse's feet. Sheridan's fame spread around the world in a few hours. In my case, sev- eral years elapsed before I dared look into the face of either historian or poet.] I can't say much for the tone of this violin. Truthfully, it has not any tone worth mentioning. Its tone is not only "dead" but 'tis "buried. " The only value in its tone is an example of how not to do interior surface work and as a warning against quick return fever. To apply the bow before protecting mater- ial upon either surface of the violin plates becomes dry is an act of insanity. Had the historian re- ceived access to the "sticky" facts clinging to the inside of this violin, then the violin student today

120

\section*{VIOLIN TONE-PECULIARITIES.}
would be reading that interior-surface protection for the violin is a failure. Had I myself abandon- ed effort after this disaster, I would now believe that such protection causes "dead" tone. I did believe it during the next half-year. Then in a moment of idleness I picked up this discarded vio- lin and again applied the bow. Surprise was lying in wait for me. In the qualities of intensity and brilliance, its tone now surpassed former rating. Quickly I again removed the sounding-board. You who have read the description of varnish phenomenon No. 1 may think yourselves prepared for a description of varnish phenomenon No. 2. Let me warn you that you are not pre- pared. You remember that the varnish dis- turbance in No. 1 took the form of crater-like openings. Only in one point is there similarity be- tween these two phenomena both are located up- on the sounding-board. As I view this fact, it af- fords further proof of the dominant part taken by the violin sounding-board in the production of sound. Both of these varnish disturbances point to greater amplitude, and greater power in sound- ing-board oscillation as the cause for their exist- ence. Varnish phenomenon No. 2 is in the form of pen- dant drops, and hanging from the inner surface of the sounding-board like stalacities from the dome of Mammoth Cave. These drops are not found ov- er the entire sounding-board. They only exist in the central area, and within lines drawn over the center of the exits. The graduation of this sound-

121

\section*{VIOLIN TONE-PECULIARITIES.}
ing-board! is 2 and 8-10 mm down to 2 and 7-10. Why these' tv/6 phenomena vary so diametrically in physical appearance is beyond my understanding. From this crude beginning I continued until feel- ing\textasciicircum\{\}that defeat is turned 'into victory. In the ex- perimental stage of this work I tried various sub- stances for- interior surface protection. One of these substances or combination of two substances, although discarded as valueless, yet possesses in- terest because of showing the effect upon violin tone from bending the sounding-board. This mix- ture is composed of boiled linseed oil and transpar- ent glue. The union of these two substances is but mechanical. The presence of oil causes flexib- ility in the mass when dry. So great is the flexibil- ity that, when dried in thin lamina, they may be in- definitely bent upon themselves without fracture.' Because of such physical quality, I gave extended trial to this mixture. In this combination the pro- portipnate quantity is 25 drops of oil to one fluid ounce of glue,.having the proper consistence for be- ing applied with a brush. There are two .objections to this mixture as a material for violin interior-surface protection. Either of these objections is sufficient to condemn its use upon the vioffru- First: Moisture from- the glue is drawn into the- wood. Second: , In drying, this mixture shrinks with sufficient power to bend the sounding-beard. By repeated coatings, the -"cuppings" of the sounding-board may be deepen- .ed until arching is raised 1-4 inch. With the ) sounding-board thus bent, the tone undergoes- a

122

\section*{VIOLIN TONE-PECULIARITIES.}
a surprising change. Such changes are manifest- ed by short duration and, by strident quality, (as the croupy voice, ) and, by increased brilliance, or distinctness in passage played rapidament. There are violinists pleased with the shortened duration of tone from the bent sounding board. These are players whose forte lies in rapidity of execution. As is well known, there are violins on- ly producing a confused jumble of sound in rapid passages. Comparing such violins with those hav- ing this strident, but brilliant tone I prefer the lat- ter. But, the strident and brilliant tone, by no means, compares with the natural and brilliant tone, as I view this matter. The strident is a thin tone, and therefore an unnatural tone. For some occult reason, different violin tone- tastes more nearly approach infinity than tone-tastes for any other musical instrument. My amateur violin-making friend, J. D. O'Brien, Pittsburg, Pa., during his last European tour, had the pleasure, or rather the pain of seeing five worn out Stradivari left in care of Hammeg, violin mak- er, Berlin. One of these violins was once present- ed to Dr. Joachim, by his London 'Why admirers. In late years I have asked, 4 do we never read of the worn out Amati, or the worn out Gua- neri, or the worn out Maggini, or even the worn out Da Salo?" I receive no answer thereto. Sup- position seems to be the only source for an an- swer. As heretofore stated, that sounding-board yielding the most agreeable tone is the first in suc- cumbing to the disintegrating forces of heat and

123

\section*{VIOLIN TONE-PECULIARITIES.}
moisture. Of course the violin yielding the more agreeable tone is likely to receive the greater amount of use in any given period of time. But alone, legitimate use cannot account for rapidity of disintegration. Were my name Hammeg, before one hour older, those worn out sounding-boards would be subjected to tests for density of hard fibre and connective tissue between. Today I entertain conviction that every worn out sounding-board, taken in time and given correct interior-surface protection, might have retained all original tone- value to an indefinite period. Upon this point, I have made ample, 'practical demonstra- tion, and have satisfied myself that such conclusion is based on safe reasons. I have given both correct and erroneous details for interior-surface protection that you may satisfy yourself with much less trouble. You may ask yourself the question, "What logic is there in protecting one surface of violin plates while leaving one surface without any protection whatever?" You will receive various answers to this question. In all probability you will receive positive replies from parties who prac- tically know nothing about this subject of interior- surface protection for violins, that such protection locks up secretions in the wood. Certainly it does lock them up that is, if there is any secretion therein. If you have read me closely, you observ- ed that I am careful to repeat the fact of only ap- plying interior-surface protection to used violins violins having had time to complete the process of shrinking after leaving the builder's hands; and

124

\section*{VIOLIN TONE-PECULIARITIES*}
that I am careful to further dry out the wood with artificial heat when necessary. Should you do interior-surface protection work* carefully following correct details, surprise will lie in wait for your bow. Augmented intensity and brilliance of tone will be something new to your old violin; and, if its tone was valuable before the in- terior surfaces were protected and made perfectly and permanently smooth, you will enjoy the fact of knowing such increased tone-value to be also permanent. Permanence of violin tone-value is a very desir- able quality. It is yours for the asking. Possibly you may not entertain sentiments iden- tical with mine regarding the inanity in certain statements concerning the Cremona violin varnish. As students of the violin, your attention has been called many times to this subject. You have read of it in books; in booklets; and in periodicals; even in newspaper ' Interviews. ' ' Possibly, such reading did not make the same impression upon you as up- on me. Nothing in life is more patent than the fact that mental impressions may vary as the num- ber of minds. In early life I read something about the beauty of Cremona violin varnish. I then thought the writer re ally, referred to gums com- posing, the varnish. In middle life, I began to read that to Cremona varnish is due the superiority of Cremona violin tone. In advanced life, I began to read that Cremona varnish, in some way, permeat- ed the wood; also, that Cremona violin varnish sud-

125

\section*{VIOLIN TONE-PECULIARITIES.}
denly, mysteriously, disappeared from commercial channels. Strangely, I've read nothing about the disappearance of other furniture varnish. Tis passing strange that fate should select violin var- nish for annihilation! Beyond all degrees of strangeness is the fact that but a limited number of Cremona violins are, or ever were celebrated for superiority of tone. Pointing directly to the two most famous Cremonese violin builders, Stradivar- ius and Joseph Guarnerius, They reliable "experts" have said to me. ' did not make every violin of equal tone- value. ' ' "Why not?" They had access to Cremona varnish! Really\textasciitilde{}we may now expect an "interview" from some London "expert" (trade promoter rather,) wherein he states that Cremona varnish was so dear (oh!) that none could afford to use it except in a limited way. . , Perhaps he's right. . The use of a varnish which soaks into the wood must require a lot of it. Italians are now, and for a long time have been meritoriously noted for superiority of color-sense. But, even in violin color-work, they did not make all violins equally beautiful. I by no means pose as a color artist, but, from my experience in color- work upon wood, I am confident that no one can make what.is called a beautiful violin without hav- ing beautiful wood as a pre-requisite. Even then failure is easy. Six times have I removed color- work from a single violin before finding the right

126

\section*{VIOLIN TONE-PECULIARITIES.}
combination to match that particular sample of wood; and a long interval may elapse before that combination will be equally effective upon another sample. People exclaim, "What beautiful varnish!" The beauty is not all in the gums. 'Tis all in the combination of colors. Not one iota would I subtract from the rightful merits earned by the "old masters." Neither can I give them unearned credit. For those empty statements, intended only for "trade promotion" in Cremona violins, I entertain nothing but con- tempt. Doubtless my contempt for those London ' 'inter- views" and for those London authors whose ani- mus is plainly "trade" is quite paralleled by their contempt for the Americans. The Londoners es- timate of the Americans is clearly shown in the fol- lowing words credited to Kipling. 'Twas on board a sailing vessel from Calcutta to London. One morn- ing, when off the west coast of Africa, the watch re- ported being passed at daybreak by the largest sea- serpent on record. The watch stated that their attention was first attracted to the monster by a steam-like hissing as it came up under the lee counter. Its forked tongue was eight feet long. Its head was ten feet in length from the point of the jaw to the single fiery headlight eye. Its head was carried fifteen feet above water. Its body was longer than the ship. While listening to this sailor yarn, a New York newspaper man, busy with pencil and note-book, remarked to Kipling

127

\section*{VIOLIN TONE-PECULIARITIES.}
that upon arrival at London he would report this remarkable incident to the press. "Don't," said Kipling. "Why not," replied the New Yorker. "My friend," replied Kipling, "let me tell you something which you don't seem to know. England was gray headed before you were born. Keep this story until you reach home. ' ' We must admit the truth in Kipling's remark. England does hold a monopoly of gray heads. 'Twill not be her fault should a single gray head ever appear in any other country. Something ov- er a hundred years ago she received a hint from Americans that we desired to live our natural lives whether or not that life might be long enough to reach the gray head stage. From signs of the present, there is hope for us. From Portland to Seattle, from New Orleans to Duluth, there are those who no longer swallow the London expert's "varnish permeation-of- the- wood" story. Today the London expert, "interviewing" himself bewails (!) the prohibitive prices for the few old violins permeated by the lost Cremona varnish. Judging by the past, what harm in prophecy for the fut- ure? Entertaining conviction, I confidently pre- dict that, within the present century, for the best violin, Europeans will search throughout America. Why not? Cremona violins, penetrated by varnish applied only upon a single surface, will soon be out of the market. It is certain that a few American violins, penetrated by varnish applied to both surfaces, will then be on the market. Why not picture de-

128

VIOLIN TONE-PECULIARITIES. scendants of the present London expert as then interviewing themselves about prohibitive prices for American double-permeated violins? There is one American gray- head who would immensely en- joy that picture.

129

\section*{VIOLIN TONE-PECULIARITIES.}
LECTURE VIII. GENTLEMEN: We are approaching problems in violin tone which are not only difficult of solution, but are also difficult of enunciation. Silently I sat at my bench many years, thinking, working, but never writing. During these years I was favored at long and brief intervals only by the companion- ship of persons who cared for the philosophy in- volved in violin tone. As is well known, much talking on any subject greatly facilitates the selec- tion of precise words to make our meaning stand out in a clear light. Being deprived of such bene- fit, and receiving little or no direct benefit from either text-books of philosophy, or from books of general reading, I therefore find this task of precise diction to be a matter of difficulty. My age and infirmities also handicap my pen. Therefore, should I fail in the matter of precision, 'twill cause no surprise to myself. I am comforted by the as- surance that each succeeding generation becomes brighter in intellect, hence more capable of precis- ion in enunciation of both principles and solutions. Because that problem in the arching of violin plates securing maximum concentration of molecular movement at the exits is abandoned by physicists as insolvable, by no means do I lean back and say that this problem never will be solved. I firmly believe a solution for this problem yet will be forthcoming. To deny such possibility is equally as absurd as saying that the violin reached perfec- tion 200 years ago. Within my observation, violins are made today of greater tone power, and greater

130

\section*{VIOLIN TONE-PECULIARITIES.}
evenness of tone than at any other period in violin history. [Galilleo did not discover the Americas, but, Galilleo's enunciation of principles made their discovery possible. After Columbus had demon- strated that Galilleo's principles were correct, it is related that a vaporous party at the dinner table jumped up to remark that this discovery was an easy matter any one could have made it nothing in it except the matter of sailing west until stop- ped by land, then sailing back again. Calling for an egg, Columbus asked each person present to make the egg stand upon its smaller end. After all had failed, Columbus, striking the egg upon the table with sufficient force to partly crush the shell, made the egg stand firmly upon its smaller end. This mythical story is here told for a purpose. Should I succeed in even cracking the shell of some difficult problem in violin tone, then that party who later succeeds in making such problem "stand on end" I shall designate as Columbus if he but hints that my name is Galilleo.] 'Tis well enough to smile while having a chance, for soon we shall enter a territory wherein dryness precludes the possibility of smiles. \_ After crossing this "dry district" we will try to find moisture for our parched lips and duly celebrate our "passover." Uniformity in violin tone- values is considered one of the unsolvabe problems. We can truthfully say that this problem has successfully defied solution during 400 years. That this problem ever will be solved, we can only hope. In attempting a solu-

131

\section*{VIOLIN TONE-PECULIARITIES.}
tion of this problem, I accomplished something, but not enough. Only in a limited number of violins have I succeeded in securing such uniformity of tone- values as to prevent me from determining the tone of each one when played while hidden from sight. In this respect, violin tone is as the voice of singers. Only at rare intervals do we hear two singers possessing identical vocal qualities. In the solution of this problem, what immeasur- able value to the violin! Tis a value defying arithmetic. Feeling certain that this problem is of interest to the violin student, I therefore describe details in my work for the production of uniformity of tone- values. First, I select violins as nearly uniform in dimension as possible. It is at once apparent that much lack of uniformity in cubic capacity must defeat uniformity of tone qualities. Thus, when two violins vary much in length of perpendicular air columns within the body, it is a difficult matter to bring their tone to uniformity in pitch. [I do not mean difficulty in tuning the strings to what is called "unison," but do mean the difficulty in preventing the ear from distinguishing the pres- ence of two sounds varying in tone-pitch.] Second, I select violins having wood as nearly uniform as possible in maturity, density, and width of grain. These latter points exert much influence in the the quality of violin tone. After thus selecting a half-dozen violins, I remove the varnish therefrom. Such removal is necessary be- cause both quality and quantity of varnish operate

132

\section*{VIOLIN TONE-PECULIARITIES.}
to raise tone-pitch and to diminish volume of tone. These violins are next opened and necessary work is applied to their interior surfaces. This work may vary in each violin. Some of them may need re-barring; some need re-graduation; some need both re-barring and re-graduation; possibly some may need heavier linings and corner blocks; and in all, the interior surface is made perfectly and per- manently smooth by means of elastic varnish ap- plied in such attenuation as is only possible to the ' 'rubbing' ' process previously described. Of course in the re-barring and re-graduation, account of varying degrees of density in sounding-board fiber must be considered. After assembling, comes ad- justment of the finger-board. Its hollow under surface, beginning at the base of the neck, is made so straight that the straight-edge touches at all points in the length of the hollow. [The philosophy for thus fashioning the under surface of the finger-board will appear later in the discussion of tone-modifiers.] The under surface of the finger-board is placed at a uniform height throughout the half-dozen in- struments. The weight and density of the finger- boards are as uniform as possible. While these violins are yet in "the white," bridges of equal maturity, density, and width of grain, and guage 2, six-strand, hard-twisted strings, each of even diameter are adjusted, and the preliminary test for uniformity in tone values is made. This test is made prior to varnishing the exterior surface for the purpose of retaining the opportunity for work-

133

\section*{VIOLIN TONE-PECULIARITIES.}
ing upon such surface to lower the tone-pitch where found necessary. Thus, after applying the bow upon each violin, I select one having an aver- age tone-pitchnot highest nor lowest in the lot then raise the tone-pitch of those below, and lower the tone-pitch of those above this selection. In the work of raising tone-pitch, I now have three, or, counting varnish, four means for assistance, thus: (1) Enlarging the exits. (2) Position of bridge. (3) Position of post. (4) Quantity of varnish. For the work of lowering tone-pitch, I have three means, thus: (1) Diminishing sounding-board thickness. (2) Position of bridge. (3) Position of post. [As a matter of fact, a fifth means for raising tone-pitch may be found by counting diminished depth of ribs; (See Rule VI) but, being unncessary, work except in extreme cases, I do not here find myself compelled to employ it.] The philosophy involved in this work of lower- ing tone-pitch is reserved for a later date; there- fore at the present moment it is sufficient to state that after the tone-pitch of these violins becomes so nearly identical that my sense of hearing fails to distinguish the pitch of one from another when played while hidden from view, they are then dis- mantled and varnish, of equal quality and quanti- ty, is applied upon the exterior surface and by the "rubbing" process. When again ready for use,

134

\section*{VIOLIN TONE-PECULIARITIES.}
they are taken to the open field for their intensity test heretofore described. Before giving results of the long-distance test for intensity, (carrying power) I will mention the fact that in the qualities of tone-pitch and intensity of tone is shown the nearest proximity to equality of tone- values in these violins as now prepared. At near-by dis- tances there is yet a notable difference of harmon- ics a bassa and consonant harmonic overtones. Thus, their tone varies in ' 'richness." Again, there is a difference in response to bow-pressure; and particularly observable in greater or less dis- tinctness of tones played rapidament. I now return to the long distance test in the open. Prior to the herein described treatment, the record shows intensity of tone in these violins as straggling along from 400 to 800 feet. After prep- aration, as described, the record shows an increase in intensity of tone up to 1,100 and 1,200 feet, av- eraging 1,175 feet an increase of 80 per cent in distance; and an increase, in uniformity of tone- value upon this isolated tone-quality, of 62 per cent. I know of no way to describe precisely the per cent of increase for such tone-quali- ties as volume, evenness, richness, freedom from dissonant overtones, sympathy in concert, response to bow, agreeability of double-stop tones, brilliance of tone, and power of harmonic tones. In a gener- al way, and in my opinion, the latter tone qualities are increased equally with intensity, and in the majority, I believe the per cent of increase to be greater than the increase of intensity.

135

\section*{VIOLIN TONE-PECULIARITIES.}
There is yet another tone-value in this lot of violins. It is a value difficult of accurate descrip- tion. I refer to the effect produced by playing these violins in unison. I believe that my experi- ence in listening to violins thus prepared has no parallel. This belief causes me to hesitate before attempting a description of such effect. In this description, I cannot find support by way of com- parison. I am obliged to find words descriptive of a new impression upon both the sense of hearing and the sense of feeling. The first word to come out of the mist is solidity. When saying that any certain object is "solid," we mean that such object possesses weight. As a rule, the word solid is used only to describe objects which can be seen. Yet, I am inclined to use that word in describing the combin- ed tone of this half-dozen violins. Certain it is that while listening thereto, and while at a distance of more than 1,000 feet, I seemed to feel sound- waves striking against my person. As you remem- ber, the long-distance test is made in open air, under a cloudless sky, winds at rest, temperature from 70 to 90 degrees Fah. water vapor at, or be- low the normal point, (never , above) and at the hours between 10 a. m. and 4 p. m., (hours where- in sound-waves are propagated with the greatest difficulty upon any given date) and, upon level ground, and with no surrounding objects capable of reflecting lines of molecular movement. Under such circumstances, you at once observe that the position occupied by these violins, during such

136

\section*{VIOLIN TONE-PECULIARITIES.}
tone-test is at the center of a circle whose radius is greater than 1,000 feet. You also understand that, under the above conditions, sound-waves from these violins reach the periphery of this circle at all points. I fully realize at the present moment, that my statement concerning the "solidity" of the tone from these violins, prepared as herein set forth, has no corrobation. I therefore ask you to make such trial yourselves. In all matters of practical appli- cation to violin tone, my work is finished. Your- selves, being alive to the fact that in certain con- ditions where sound-waves are propagated only with difficulty, and knowing the complaint of not hearing first-violin tone, will take up this work where I leave it and demonstrate its value in both the small and large orchestra. In this work, I cannot give assurance that any one can succeed without experience. As I glance backwards at the years I've devoted to violin tone- peculiarities, and note the difficulties upon either hand of the violin-builder's path, I am constrained to say that success depends largely upon experi- ence. The spirited horse can be goaded into mad action by a galling ill-fitting harness. The violin sounding-board can also be driven into mad action by an ill-fitting sound-post, by an ill-adjusted fing- er-board, by an ill adjusted bridge, by an ill-assort- ed set of strings, and, by a worthless bow. When all of these ills combine, imagination fails. So long as one-hundredth of an inch in variation of sounding-board thickness, in string diameter, in

137

\section*{VIOLIN TONE-PECULIARITIES.}
bridge-thickness and position, in sound-post thick- ness and position, in area and position of exits, in depth of rib, in diameter of bow-hair, affects violin tone, so long will success defy the hand-ax. Not more can a carpenter succeed in violin-making than a blacksmith can succeed in watch-making. Do not understand me to say that a carpenter may not become a violin maker, nor that a blacksmith may not become a watch maker. That is another prop- osition. But, I do say, emphatically, that both must serve an apprenticeship. To illustrate the truthfulness in the latter state- ment, it is necessary only to call up a single diffi- culty encountered by the violin builder. This dif- ficulty lies in variations in sensitiveness of sound- ing-board wood. These variations are wide enough to cause not only profound surprise, but also, profound disappointment. I know of one sounding- board, and but one for that matter, which is not affected by five times the amount of varnish neces- sary for mere protection. Again, I have handled sounding-boards so sensitive that but a single atten- uated coat of varnish plainly caused damaging ef- fects upon tone. In view of these facts, it is evident that only such judgement as must come from ex- perience can secure best results. My experience causes me to believe that every good violin, built upon hard-and-fast rules, is but an accident. It is my desire to be of assistance to all such as attempt the work of violin tone-regulation. To this end, I can think of no way more effective than pointing to such facts as endanger success. 'Tis of

138

\section*{VIOLIN TONE-PECULIARITIES.}
value to know how not to do .any certain kind of work. In the matter of applying protecting material to interior surfaces of new violins, I have no experi- ence to offer. In my judgement, this matter hinges upon the degree of dryness and completion of the shrinking process existing in each sample of wood. Manifestly, it is inadvisable to hermetically seal up violin wood prior to the completion of dry- ing and shrinkage. In our climate, the presence of water vapor in air. operates to prevent both com- plete shrinkage and drying of wood. When sur- rounded by water vapor* the capillaries of violin wood are bound to draw in and hold more or less water. It is also manifest that the presence of water in the capillaries diminishes resonance of every violin. No other conclusion is possible. Again, it is manifest that hermetical sealing up of violin wood effectively prevents further entrance of water. Thus, violin interior-surface protection, aside from its value as a perfect reflecting medium, is also valuable in permanently maintaining reson- ance, c I repeat the statement that, prior to her- metical sealing, violin wood should be absolutely dry by artificial means. In our climate, we may not expect absolute dryness of wood when left to nature. Observation of unequal shrinkage in vio- lin plates, subsequent to leaving the builder's hands, leads me to the opinion that, for the reason of certainty in new violins, the plates, nearly re- duced to correct thickness, should thereafter be allowed one or more years to complete the process

139

\section*{VIOLIN TONE-PECULIARITIES.}
of shrinkage before being assembled. By no means do I desire that you understand me as advising uninterrupted effort to produce violin- tone of maximum volume and intensity of tone. The necessity for .violins ,of "big tone" only exists in a limited way. The total number of auditoria, creat- ing necessity for "big tone," is a small number in comparison. As the object of all musical interpre- tation is the pleasure of the listening ear, it there- fore follows that both disagreeable tone quality, and the tone so feeble as not to be heard at all, de- feat the object of musical effort. Thus, the mis- take of employing the violin possessing "big tone" for. studio, parlor, or the small auditoria is only parallelled by the mistake of employing the weak tone for the large auditoria. In small rooms, the ."big" tone is painful; in the large room, the weak tone is disappointing. This proposition has two sides; one, aesthetic, the other, "business." It is quite safe to assume that all patrons of music halls. expect pleasure from first-violin tone. It is also quite safe to assume that either disagreeable first- violin tone, or inability to hear first- violin tone operate to reduce patronage. In this matter, the patron is nowise to be blamed for withdrawal. He parts with wealth for expected pleasure and finds but disappointment. H It .is evident upon either hand that the great ma- jority of violins are used only in comparatively small rooms; In view of this fact, 'tis wisdom to build and tone-regulate the great majority of vio- lins minus maximum volume and intensity

140

\section*{VIOLIN TONE-PECULIARITIES.}
of tone. To-day the skillful violin maker can come nearer producing violin tone "to order" than at any other period. I present the following factors as proving them- selves reliable in diminishing both volume and in- tensity of violin tone: (1) Control of sounding-board action by var- nish.

(2) Control of sounding-board action by bend- ing.

(3) Control of sounding-board action by thick- ness of wood. (4) Control of sounding-board action by arch- ing of plates. (5) Control of sounding-board action by the bar.

(6) Control of sounding-board action by the post. (7) Control of sounding-board action by the bridge. (8) Control of sounding-board action by the finger-board. (9) Control of sounding-board action by the di- ameter of strings. (10) Diminishing volume and intensity of tone by the condition of interior surfaces. (11) 2 Diminishing volume and intensity of tone by the area and position of exits. (12) \textasciicircum\{\}Diminishing volume of tone by the depth of ribs. Either of these factors, alone, is capable of pro- ducing perceptible diminution of both volume and

141

\section*{VIOLIN TONE-PECULIARITIES.}
intensity of violin tone. ( 1 ) Control of sounding-board action by varnish , in all cases except the one case herein mentioned, has proven to be possible, and quite easy of accom- plishment. The only difficulty that I have found in diminishing volume and intensity of violin tone by means of elastic varnish was wholly due to in- herent variations in the action of sounding-board wood. These variations are so wide as often to defeat tone-regulation by "hard-and-fast" rules. In my hands, the greatest effect from varnish, in controlling sounding-board action, has been shown upon wood of soft fiber; and the least effect has appeared upon wood of dense fiber. Thus, in the employment of a certain quantity of varnish upon each violin, its effect upon tone varies as the de- grees of density in sounding-board wood. I have not found that any amount whatever of varnish, applied upon the back, operates to dimin- ish tone in any degree whatever. On the contrary I have observed a slight increase in tone-power from heavy coats of varnish applied upon back plates whose thickness had been reduced to that point causing tone-weakness. Without doubt, any amount of varnish increases rigidity of violin plates; and, upon the sounding-board, any amount of var- nish also operates to diminish independent action of contiguous fibers. Thus, varnish, upon the sounding-board, operates to diminish both volume and intensity of tone. Often have I observed vio- lins having greater power in single, than in double- stop tones. In all such cases, I have found increase

142

\section*{VIOLIN TONE-PECULIARITIES.}
of power in double-stop tone following removal of varnish from the sounding-board. To demonstrate that varnish, upon the violin sounding-board, oper- ates to diminish both volume and intensity of tone is a matter easy of accomplishment. In its effect upon violin tone, varnish possesses an esthetic side. There are many musicians who prefer the natural tone from untrammeled wood. When such tone possesses beautiful quality, there can be no objection from any one. The beauty in such sound is impossible of enhancement by any means known to me. But, herein lies the difficulty. Inherent capriciousness in the action of sounding- board wood produces tone from the extreme of beautiful down to the extreme of coarseness. This matter wholly lies in the domain of nature. Man can only modify it. There are a limited number of musicians who profess pleasure in coarseness of violin tone; but, the great majority prefer such coarseness to be modified by varnish, and content themselves for the loss in volume and intensity by the loss in disagreeable tone-quality. (2) Bending the sounding-board into position is a powerful factor to diminish, not only volume and intensity of tone, but also, to diminish duration of tone. Violinists, whose forte lies in rapidity of ex- ecution, are sometimes satisfied with the shortened and lifeless tone from the bent sounding-board. The influence upon volume, intensity, and dura- tion of tone from bending the sounding-board is in exact proportion to the degree of bending. For diminishing these tone qualities, my own taste does

143

\section*{VIOLIN TONE-PECULIARITIES.}
not approve of the bent sounding-board, because of a certain "dead" tone-quality. I prefer to di- minish the area of the exits, while leaving the ac- tion of the sounding-board free; thus, preserving liveliness of tone. The difference in tone-value is vastly in favor of the latter method. I do not find that bending the back into position operates in any degree whatever to diminish power of tone; but, on the contrary, bending a back plate, already too light in wood, increases the tone-pow- er. But, such increase in tone-power does not hold good indefinitely. Here is the reason. Upon opening the violin with bent back, and within but a few years after leaving the builder's hands, I find the bending transferred to the sounding-board. This result might be expected did we but give it- thought; and the two factors, causing loss of tone- power to accompany such transference of bending, might also became apparent with thought. The trouble is to keep up to the point of hard thinking. 'Tis exhausting. Facts are easiest made known by stumbling upoir them. Stumbling costs nothing but picking one's self up, and remarking. (3) Control of sounding-board action by thick- ness of wood is a factor of great importance. This factor has received more attention than any other single item in the list of violin tone modifiers. Only the effect of diminishing volume and intensity of tone by thickness of sounding-board wood will receive attention. [Upon a later occasion, and in the discussion of maximum evenness of tone, this important tone-

\section*{\textasciicircum\{\} 144}
\section*{VIOLIN TONE-PECULIARITIES.}
modifier will be given further attention.] It is evident that the thickness of the violin sounding-board must be governed by the diameter of strings to be employed. When not otherwise specifically mentioned, I refer only to guage-2 strings. [As previously stated, my experience compels me to the belief that the function of the violin sounding-board is to originate those sound-waves eventually reaching the ear as musical tone. Pre- vious to varnish phenomenon No. 1, I directed work upon the back plate as if it also were a tone-pro- ducing agent. To such work I now attribute ruin of tone- value to a number of violins. After confin- ing the work of tonerregulation wholly upon the sounding-board, I met with no failure attributable to my work. By placing new, strong back , plates upon violins having suffered serious injury to tone from attempted tone-regulation upon that plate, I succeeded, to a satisfactory degree, in restoring lost tone- value. I do not wish to be understood as saying that loss of tone-value may not also follow erroneous work upon the sounding- board. On the contrary, I do say that that tone- value may be easily lost by erroneous work upon the sounding-board. What I do say with positive- ness is that, after confining the work of tone-regu- lation to the sounding-board, I met with vastly greater success in securing tone-value in such tone qualities as intensity, brilliance, and sweetness. Volume of tone, at nearby distances, I could ob- tain with certainty by work upon the back plates

145

\section*{VIOLIN TONE-PECULIARITIES.}
alone; but, such tone invariably fell far short in the long-distance test out in the open; whereas, with a strong, unyielding back plate, I could se- cure equal volume of tone and vastly greater in- tensity and brilliance of tone by work directed to the sounding-board. This fact I have demonstrat- ed by many repetitions of the long-distance out-of- doors test. As previously stated, this test affords evidence based upon actual measurement, in lineal feet, of intensity (carrying power) of violin tone. It is a test settling this question beyond the shad- ow of doubt. ii , In diminishing volume and intensity of tone by thickness of sounding-board by no means have I found the same difficulty as is presented in the work of securing the maximum of those tone qual- ities. In securing the maximum of volume and in- tensity of tone, the work, from a certain point of thickness, must proceed with caution in the mat- ter of measurements of thickness, and removal of wood; and because of the reason that in every sam- ple of sounding-board wood there is a certain degree of thickness yielding such maximum, and, below this degree of thickness means weakness of tone, whereas, for the production of diminished volume and intensity of tone, the thickness may run from 16-64> down to 10-64, or even lower with wood of dense fibre. These great thicknesses invariably operate to diminish both volume and intensity of tone, but not in equal degrees; the greater diminu- tion appearing in volume of tone. These great thicknesses also operate to raise tone-pitch.

146

\section*{VIOLIN TONE-PECULIARITIES.}
It is a fact in my observation, that every inherent disagreeable tone-equality in sounding- board wood developes and becomes apparent as thickness of wood diminishes to the point produc- tive of maximum volume of tone. I regard the knowledge and consideration of this fact to be of value to both the violin builder and those violin us- ers who prefer quality of violin tone before quantity of tone. In my observation, there are many violin users who entertain such preference. It is also a fact in my observation, that for all occasions and circumstances, a violinist can become equipped only by having at command a number of violins possessing varying degrees of volume and intensi- ty of tone. (4). The arching given to violin plates defies the pencil of that physicist who would write down its solution. Yet, 'tis an easy matter to demon- strate the fact that such arching may add to, or diminish both volume and intensity of tone. I know of no single item in the list of violin tone- modifiers affording greater interests to the violin student than the item of plate-arching. In practi- cal application, there is no diminution of interest from the perfectly flat plate on up through every degree to a height of arch equaling 1 inch. Other dimensions remaining equal, power of tone increas- es, from no arch at all, up to a certain height of arch; and from this certain point on up to that height equaling 1 inch, power of tone steadily di- minishes. Other dimensions not remaining equal, as increasing the area of exits, increasing sounding

147

\section*{VIOLIN TONE-PECULIARITIES.}
board thickness together with an increase of string diameter, power of tone may follow to a higher de- gree of arching than in the first proposition. It is apparent that the diameter of strings may safely govern the height of arching. In my obser- vation, using the guage-2 string, and distributing the arching equally from either end of the plates to the position of the bridge, the greatest volume of tone is reached at that height of arch equalling f inch, while the greatest intensity of tone is reached at that height of arching equalling i inch. I desire that the above statement be taken with the understanding that the area of exits, in either of the above heights of arching remains pre- cisely equal, and that such area is neither large nor small, and that the width and length of the body and the depth of ribs be not greater than that of "full size"; also, that the distance between the ex- its, at their upper extremities, be exactly 1 and 4 inches. These conditions are absolutely necessary in determining the influence of arching upon vol- ume and intensity of tone. - I present this violin as an example of the influ- ence exerted upon tone-power by plate arching. You observe that the exits are of medium area; and that their position is> neither "high", nor "low." As you look at this violin in profile, you observe that its waist line is aldermanic; that its height of arch, at the position of the bridge equals \& inch; that the distribution of the arch is equal from the ends of the plates to the bridge; that the quality of wood, varnish andwork-

148

\section*{VIOLIN TONE-PECULIARITIES.}
manship is good. As I apply the bow, you observe that it has a tone-power equalling the tone-power of a pewee; and even that tone is "all inside". In attempting an explanation for this great loss of tone-power, due to such formidable height of arch, we are left with nothing but supposition. We know the fact, but positive assertion can go no further. Although knowing mere surmise to be uninteresting, yet, I'll take the risk of presenting my surmise for an explanation of this defiant prob- lem if you'll not charge me with positive assertion in the matter. I fully understand that assertion, without corroboration, proves nothing. Therefore, I begin by saying thus: I think loss of tone-power from high-arched violin plates, other dimensions remaining equal, is due to the fact that the exits are not in the line of wave sound movement.

"Well, now comes the defiance." If, by possible means, we could place ourselves within the body of this violin while the bow excites the strings to action, and if our vision were quick enough to follow the line of travel taken, by any single wave-movement as it originates upon the in- terior surface of the sounding board, then we would know why such movement indefinitely con- tinues to travel "all inside;" then we would know why the great majority of such movements fail to pass out through the exits made and provided, especially for the egress of sound waves. I can safely say that those "if's" have used up the last grain of plumbago in many a pencil.

149

\section*{VIOLIN TONE-PECULIARITIES.}
LECTURE IX. GENTLEMEN: At this hour the presentation of violin tone-modifiers, operating to diminish volume and intensity of tone, is resumed. At the close of the preceding hour we were yet considering the great influence upon tone from the arching given to the plates. In continuation, I present this violin as another example of arching, which invariably operates to diminish tone-power. As you view this violin in profile, you observe that the entire rise of the arch is given to the first two inches at either end of the plate; that between such abrupt elevations, the plates present a straight line; that other dimensions remain as usual. Application of the bow demonstrates the fact that the tone-power of this violin, while perceptibly greater than the tone-power of the higher arched violin, yet lacks much of the maximum. As I view this model of arch for violin plates, it is a model affording but a a few degrees of greater tone-power than no arch at all. Were it not for the slight concentration of sound-waves at the exits due to the lateral, or cross-section arch, there would be no difference whatever in tone-power. Thus the greatest dim- inution in volume and intensity of tone due to arch- ing, other dimensions remaining equal, is found at the extremes; that is, the perfectly flat plate, and the enormorous arching previously indicated. In stating my belief concerning the cause of such diminution of tone-power, I acknowledge inability to furnish proof. Were there but a single arch in the violin plates, the proof would be forth-coming

150

\section*{VIOLIN TONE-PECULIARITIES.}
in but a few moments. Tis the presence of the lateral arch that adds complication to this problem. With but the greater arch, (from end to end of plates) to contend with, the line of travel followed by sound-waves, originating at the sounding-board, could be accurately traced on paper. The physical laws explaining sound-wave movements, are easily comprehended. > The two laws concerning this problem are: ; (a) Sound-waves, at the moment of origin, travel at a right angle from the agent producing them, (b) Sound-waves, after striking a reflecting surface, travel therefrom at an angle equal to the angle of incidence. It is obviously an easy matter to draw on paper a line at a right angle to any given point upon the interior surface of any given arch. It is also obvious that a line of sound-wave movement, originating from any point on the arched violin sounding-board, will not strike upon the back at a point perpendicular to its origin,, but, will strike the back at a points nearer the exits. So far 'tis .easy.; but, the next move o the sound-wave is wherein lies the difficulty. 'Tis now the lateral arch makes its presence known. 'Tis plain, that rafter striking the back, the wave will be reflected; but, in what direction will the reflection travel? It is the intention of the arches to direct and con- centrate wave movement at the exits. It is also the intention to place the exits in the line of wave movement. It is evident that such intentions are defeated by enormorous degrees of lateral arching, and also by no arching at all. It is apparent that

151

\section*{VIOLIN TONE-PECULIARITIES.}
when no arch whatever is given to the plates, there can be no direct progression whatever of wave movement toward the exits. Hence, the tone from such plates is diminished in power. It is evident that sound-wave lines of travel, within the high- arched violin are directed away from the exits to a 'Where great extent. In this case the question is, ' should the exits be located?" [Inasmuch as the violin is but a product of exper- ience, therefore we may continue expecting violin tone-improvement only to follow experience. With this idea uppermost in my mind, I call up some experience connected with a barrel of new cider inadvertently left out in the hot sunshine dur- a whole day. Being employed as an inspector of this particular barrel, and, the barrel being tightly closed up, I directed a moderate tap from a ham- mer to the vicinity of the bung; whereupon, and without warning, that bung flew into my face aUegro con fuoco et cider-o-so. ftow (vengefully) I do suggest cutting a bung hole in the sounding-board between bridge- feet as the proper thing for all violins having an aldermanic waist-line.] (5) Control of sounding-board action by the bar is easy of accomplishment, very easy; I feel like saying, "Too easy." Experience in wrestling with the problem of the violin bar may cause equal disappointment with experience at any other "bar" whatever. In either case, "loading up" to\textasciicircum\{\}heavily is bound to prove disap- pointing. The violin bar is a powerful factor in

152

\section*{VIOLIN TONE-PECULIARITIES.}
modifying G, and D-string tone. The bar may be too great in length, too great in thickness, too great in depth, and of too great density in fiber, and too light in mass. The G and D-string tone, from the finest sounding-board, may be easily marred by; any one of these faults, and completely ruined by their combination. The position of the bar also may seriously injure G and D-string tone. [Desiring that my statements carry proof with them, and also desiring that benefit may follow such statement, I therefore call up the case of a certain violin suffering loss in tone-value from un- bearable "wolf" on open G, and caused by mal-po- sition of the bar. This violin is made of valuable , wood, and, in all other respects, its tone-value is high. Its builder is A. F. Anderson. In my opin- ion the mechanical work on this violin cannot be excelled. The varnish is of the toughest, and most elastic variety. In my opinion, . judging from the mal-position of bar, and the graduation of the plates, Mr. Anderson tljs builds violins upon hard-and- fast rules. As Mr. Anderson is a high-class work- man, therefore it is hoped \#nd expected that he will be pleased at having an error pointed out in his otherwise faultless work.] Details are premised by the statement that I objected to opening this violin. My disinclination was due to the fact of uncertainty as to the cause of this ' 'wolf' ' tone. By no means could I assume a cure in this case. Never before had I encountered a wolf on the fundamental tone of the G-string as an isolated tone-fault. Many times I had been suc- cessful in driving a "pack of wolves" from violin

153

\section*{VIOLIN TONE-PECULIARITIES.}
tone, but "wolf" in an isolated tone is a different proposition for solution. Once I had been asked to remove "wolf" from the E-string of an old violin. 'Twas a case of "wolf" on g, and g sharp in alt. This old violin bore ample evidence that genera- tions had devoted their time to moving the sound- post. In this work, the right edge of the right exit has been worn away i inch; and both plates were bruised in a circle 1 inch in diameter by the shifting ends of the post. Not being able to de- termine the cause of "wolf" by use of the bow, I removed the sounding-board. To all appearances, the workmanship within this old violin was fault- less. I could not see any cause for "wolf;" yet, I positively knew "wolf" existed, and, as above in- dicated. Being curious to know the graduation of this sounding-board, I employed the calipers, and thereby found reasons for "guessing" at the cause of "wolf" tone in this case. Here is what I found by the use of the calipers: The builder, exercising the extreme of caution in reducing thickness of wood beneath the E-string, had gone beyond the limit of safety in such work; but, only a very little beyond. Under such meteoric conditions as here- tofore have been carefully and repeatedly describ- ed, the intensity of tone, belonging to the E-string of this violin surpasses all other E-strings which I have subjected to the out-of-doors test. The tone of this E-string carried to the great distance of 1480 lineal feet. [Bear in mind those meteoric conditions, because a change in the single item of greater amount of

154

\section*{VIOLIN TONE-PECULIARITIES.}
water vapor present at the moment of making the long distance, out-of-doors test, might double, and even quadruple this distance.] The greater power of tone possessed by this E- string concentrates interest in two factors. (a) Graduation of sounding-board beneath the E-string. (b) Enlargement of the area of the right exit. Briefly, I will describe the graduation beneath this E-string: At the position of the bridge, thick- ness equals 7-64 inch; thence gradually diminishing down to 4-64 at a point half-way from the bridge to the upper end of the plate; from this point to the end of the plate, thickness gradually increases to 9-64. Evidently, this form of sounding-board results in the production of two tapering springs between the bridge and the upper end of the plate. Practically, these tapering springs are attached to each other at their thinner extremities. It is evident that the spring farthest from the bridge is the stronger. It is well known that the stronger spring acts with greater rapidity than the weaker spring. It is evident thdt action of either of these two springs must excite action in the other spring. Because of differing degrees of thickness, it is evi- dent that such action must vary in numbers per second; therefore, these two springs, striking upon contained air at differing number of blows per sec- ond, produce two simultaneous tones; and, these two tones, being pitched at inharmonious keys, op- erate to produce "wolf." There are two methods available for remedy: (a). Equalizing spring

155

\section*{VIOLIN TONE-PECULIARITIES.}
action by reducing thickness of the stronger spring, (b). Equalizing spring-action by short- ening the weaker spring. Because of the danger of weakening tone-power by the first method, therefore I chose the second. Should the second method prove unsuccesssul then I would know the necessity for employing the first. But, the second method proved successful. Practical application of the second method consisted in gradually mov- ing both the bridge and sound-post upward, or for- ward. A curious phenomenon appeared upon E- string tone as the bridge and post approached their final resting place. That phenomenon was present- ed by shifting of the "wolf" from g, and g sharp to g sharp and a, thence, to a and b flat, and, at the next move the "wolf" disappeared. [Enlargement of the right exit in this violin, not being in point, need receive no further notice at this moment.] I return to the Anderson violin. Upon opening this violin, the cause for "wolf" was not apparent at the first glance. Without employment of the calipers, I could determine that graduation of the sounding-board was not the cause; yet, upon the sounding-board must the cause be found. While looking at the interior surface, and feeling quite uncertain of success in this case, I noticed an unusual obliquity in position of the bar. [Sometimes the doctor's diagnois of an obscure malady must depend upon negation; 'tis not this, nor that, nor the other, and so on by exclusion un-

156

\section*{VIOLIN TONE-PECULIARITIES.}
til there remains but few, or but a single probable cause. But to the doctor, there's a difference be- tween the fiddle patient and human patient. In presence of the former, the doctor need not trouble himself to "look wise."] By measurements, I found that the upper end of this bar crossed the line of the G-string, crossed the line of the D-string, and reached a point be- tween the D and A-strings. Because of not find- ing any other detail appearing as a probable cause for "wolf" upon open G, I therefore removed this bar and placed another bar with the center of its thickness directly beneath the center of the left bridge-shank, and with an obliquity equaling obliq- uity of the G-string, not greater nor less. When again in playing order, there was no "wolf" whatever upon any string. In attempting an explanation for this case of "wolf," I can submit nothing more than an opin- ion. Thus: Believing that certain sounding-board fibers, beneath each string, act to produce all pos- sible tones upon each string, therefore, in this case, I think the great obliquity of the bar caused "wolf " upon the open G tone, and my reasoning is as follows: The graduation of this sounding- board, being a modification of Stainer, therefore greatest thickness was at the position of the bridge, and thence thickness diminished down to 4-64 at all points near the edges and at the ends of the plate. Thickness at the bridge equaled 10-64. Therefore the thinnest part of this sounding-board, beneath the G-string, lay above the bar; not above

157

\section*{VIOLIN TONE-PECULIARITIES.}
the upper end of the bar, the upper end of the bar lay at a point between the D and A-strings. Bear- ing in mind the great obliquity of this bar, it is evident that sounding-board fibers beneath both G and D-strings are shortened from left to right. Thus, the longer fibers are beneath the G, the shorter fibers are beneath a point between D and A. It is evident that when these varying lengths of sounding-board fibers are aroused to action producing sound, that there must be more than one sound, or one tone. I know no way to determine the exact number of these sounds. I surmise that some of them, being at an inharmonious pitch with open G, operated to produce "noise," or "wolf" tone thereon. The cure was complete. In my experience, increasing the length of the bar beyond 10 inches operates to diminish tone- power, also, to diminish duration of tone. This statement is made with the understanding that the amount of wood in both bar and sounding-board has been accurately adjusted for the production of the maximum of tone-power. Keeping this fact in view, I state that shortening the bar, other di- mensions remaining equal, operates to lower tone- pitch, to increase volume of tone, and to diminish intensity of tone; also, conditions as above, adding to either thickness, or depth of bar, operates to raise tone-pitch and to diminish tone-power. The position of the bar may operate to increase tone- power of the G-string, and, at the same time, to diminish tone-power of the D-string. Thus, when

158

VIOLIN TONE-PECULIARITIES. the distance between the upper end of the exits equals 1 and 5-8 or 1 and 11-16, then, placing the left side of the bar flush with upper end of the ex- it operates to increase G-string power while dimin- ishing power of the D-string.

159

\section*{VIOLIN TONE-PECULIAKITIES.}
LECTURE X. GENTLEMEN: There are violin users who prefer such unequal tone-power. It is my experience that too great prominence of G-string tone is not most agreeable to the listening ear, especially when ap- pearing upon the solo violin. I deem this point to be of sufficient importance to warrant repetition of a previous statement that the beautiful harmony from evenly balanced strings, played in double-stops, is the chiefest attraction in solo violin music. (6) . The post is at once the most aggravating, powerful, indispensable, innocent appearing thing within or without, around or about the violin. It does not fall down so often as it makes us fall down. Had it but ability to laugh, 'twould be the chief "monkey." To the violin user, the post is the chief object of solicitude. Its power for good is angelic. Its power for evil is Satanic. Can this innocent-appearing thing command the power to diminish both volume and intensity of violin tone? Let it but fall down there you are! Because I enjoy frequent shots at this angel-imp, therefore, as a subject, it will not be exhausted at once. At this moment I will call up but three points connect- ed with the post which operate to diminish volume and intensity of tone: (a) . Mass of post. (b) . Position of post. (c) . Length of post. In the ordinary condition of the right exit, the greatest effect upon volume and intensity of tone, due to mass of post, cannot be demonstrated. En-

160

\section*{VIOLIN TONE-PECULIARITIES.}
largement of this exit affords an opportunity for entrance of a post having great mass. Although mass of post does not diminish tone-power to the same degree produced by some other factors, yet, such diminution is perceptible. I know of no hard- and-fast rule governing the mass of post. Like many other factors affecting violin tone, that mass of post which is best for each violin can only be determined by trial. (b). Position of post exerts greater influence upon tone-power than mass of post; that is, upon any mass which I have tried. In early life I re- member reading that the proper position for the post is i inch below the right foot of the bridge. In later life I learned to unlearn that lesson. I had to learn that the position for the post varies as the variations in sounding-board rigidity. All effects upon tone, due to the post, are manifested princi- pally upon the A and E-strings; and such effects may be largely directed upon either of these strings by position of the post. When that position is found which equally supports the A and E-strings, then, moving the post to the left operates to di- minish the tone-power of the E; and, per contra, moving the post to the right, diminishes tone-pow- er of A. Again, moving the post downwards, oper- ates to diminish tone-power of both A and E. Again, placing the post above the bridge operates to diminish tone-power of the A and E-strings in an amount equal to the re-inforcement due to sym- pathetic action between these strings and that part of the sounding-board beneath them. Placing the

161

\section*{VIOLIN TONE-PECULIARITIES.}
post above the bridge operates to transfer sound- ing-board action to the lower-half of the sounding- board; per contra, placing the post below the bridge operates to cause sounding-board action in the upper-half. I desire to be understood as refer- ring only to that sounding-board action aroused by action of the A and E-strings; also, that in either of these two positions for the post there is no per- ceptible change produced in the tone of the G and D-strings. [Upon a later occasion will appear a practical demonstration for the fact that the lower, right- quarter of the sounding-board does not act to pro- duce sound-waves when the post is placed below the right foot of the bridge.] Placing the post directly beneath the right foot of the bridge operates to diminish tone-power. In this position of the post, it is evident that a blow from the strings is expended equally upon both plates; yet, as I view this matter, the sounding- board continues to strike the greater blow upon contained air because of its greater proximity to the strings. This view is based upon the fact that sympathetic action diminishes as the square of the distance, or stated in the reverse way, proxim- ity augments sympathetic action. [To the violin tone-regulator there are experi- ences with humanity that possess more than pass- ing interest. Indeed, some of those experiences leave an impression upon memory equally as unfad- ing as the impression on "burnt wood." Thus: After you have carefully adjusted sounding-board

162

\section*{VIOLIN TONE-PECULIARITIES.}
thicknesses to respond to the action of the guage-2 strings, after the interior surfaces are carefully prepared, after the area of exits has received at- tention, after the finger-board has been carefully prepared and adjusted, after selecting and testing different densities in different bridges, after select- ing the choicest strings, after determining the ex- act position for bridge and post, after testing the tone for evenness of power in two octaves on each string, testing quality of double-stop tones, testing brilliance of tone, testing pizzicato tones, testing harmonics simple, harmonics melodic, harmonics a bassa, after pronouncing your work completed to the limit of your ability: Enter Mr. Addlepate. (Mr. A.) "I'm looking around for a first-class violin for my own use. Didn't know but I might find one at your place." Rather confidently, you place in his hands a com- pleted violin, telling him that this one is about as good as you can turn out. (Mr. A.) "May I take it home and try it? I'll take good care of it." "Certainly." After a month or two, re-enter Mr. A. (Mr. A.) "Say, this violin hasn't got just the tone I'm wanting think Sam Jones has got one a little better'n yours." As you glance at that violin your breath makes a gasp. Your carefully adjusted sound-post is gone. Standing in toward the center, and leaning backward, is a white-oak post, whittled down with a dull jack-knife; a deep notch is cut around near

163

\section*{VIOLIN TONE-PECULIARITIES.}
the upper end; in that notch is a doubled-and-twist- ed dirty cotton string tied fast with a square knot; either end of the string, passing out through either exit, is carried around behind the back, there tied in a bow knot; thence hangs 1-4 yard of festoon. "What makes you gasp so?" "Can't be that white-oak post?" "N-n-o!" "Possibly 'tis that festoon?" Under such provocation, there is but one way to manifest patriotism. That one way consists in di- recting an enthusiastic kick to the seat of Mr. Ad- dlepate's understanding. You can't miss it.] (c). Length of post, when so great as to bend the sounding-board upward, diminishes volume and duration of A and E tone, but, does not diminish intensity -in an equal degree. In fact, the post of too great length frequently operates to increase in- tensity of tone from these strings. Such augment- ed intensity of A and E-string tone is considered valuable by some violin users, especially by first- violin players in positions where sound-waves are propagated only with difficulty. The post of great length also operates to diminish power of G and D- string tone. Thus, the post manifests its power for good and for evil. (7. ) The bridge is a powerful factor in dimin- ishing both volume and intensity of tone; but, its greatest influence is manifested upon volume. The bridge possesses seven features of vast interest to the violin student:

164

\section*{VIOLIN TONE-PECULIARITIES.}
(a). Height. (b). Thickness. (c) . Span between shanks, or pedestals. (d). Density of fiber. (e) . Scroll work. (f). Maturity of wood. (g) . Position upon the sounding-board. (a) . Height of bridge may be either so high, or so low as to diminish both volume and intensity of tone. The problem of the bridge I am unable to solve without a partial solution for the finger-board problem. For the sake of clearness, it is first nec- essary to consider the finger-board insofar as prox- imity to the sounding-board is concerned. If the finger-board did not extend over the sounding- board, then the problem of bridge-height could be solved by itself. That part of the finger-board ex- tending over the sounding-board, may operate to diminish both volume and intensity of tone. This fact can be easily demonstrated. Thus: Upon a violin, having a height of arch equaling f inch, lowering the finger-board down to | inch from the sounding-board causes a weak and thin tone, even with a corresponding diminution of bridge-height. As I view this phenomenon, the reason for such loss of tone-power is as follows: Lowering the finger-board brings its surface line nearer to a par- rallel with the line of the sounding-board. It is apparent that, were these two lines perfectly par- allel, then all sound-waves, originating beneath the finger-board, would be reflected directly back upon the sounding-board; and, because of only traveling

165

\section*{VIOLIN TONE-PECULIARITIES.}
a short distance, those reflected sound-waves must strike upon the sounding-board with their initial force; a force great enough to diminish the ampli- tude of sounding-board oscillation. It is apparent that such diminished amplitude operates to dimin- ish the force of succeeding blows delivered by the sounding-board upon contained air. Hence the loss of tone-power. The problem of the finger-board involves the line of its under surface, and the placing of such line at the nearest distance from the sounding-board possible while preventing the reflected sound-waves from striking upon the sounding-board in such a manner as to diminish amplitude of its oscillation. As previously stated, I find the best lines on the under surface of the finger-board to be a straight line at the bottom of the hollow, and that the hol- low extends well towards the base of the neck. The latter point should be governed by the gradu- ation of the sounding-board, and the arching. Thus, when the graduation places the thinnest part of the sounding-board at a point half way from the position of the bridge to the upper end of the plate, then the hollow of the finger-board must not extend to the base of the neck, but should extend to a point 2 inches, or slightly more, from the base. [Upon a later date this form for the under sur- face of the finger board, together with this form for sounding-board graduation, slightly modified, will be described in the production of maximum tone-power.] It is plainly apparent that the set of the neck

166

\section*{VIOLIN TONE-PECULIARITIES.}
may operate to place the finger-board near to, or distant from the sounding-board. It is also appar- ent that when the point of widest sounding-board oscillation is nearer to the end of the plate, the greater must be the distance to the under surface of the finger-board, because the longer the hollow, the less the obliquity of the line of the hollow, the more direct the return of sound-waves to the sounding-board. There yet remains consideration of the high- arched plates. In these cases, the under-surface line of the finger-board, and its distance from the sounding-board must be quite the reverse; the low- er end of the finger-board must approach much nearer the sounding-board, while the under surf ace should be flat. At first thought, these facts seem incredible. I confess to astonishment when first seeing and playing upon a Carl Johann Flick- er. As you know, this violin is built upon enorm- orous waist-lines. The hollow at the lower end of the finger-board, (guessing at it) was less than i inch distant from the sounding-board. The height of the bridge was correspondingly lowered. Before applying the bow, I expected to hear a weak, thin tone. But, its volume was great enough, and, intensity seemed to be of the average degree for the arch. I had no opportunity to determine the height of arching given to the Flicker plates. I could see that the long arch was quite equally dis- tributed from the ends of the plates to the position of the bridge. Upon studying this situation, it be- comes apparent that to raise the lower end of this

\section*{VIOLIN TONE-PECULIARITIES.}
finger-board to a height enabling its under surface to reflect sound-waves towards the bridge is im- practical, not impossible, simply impractical; and because the prodigious height of the strings would cause too great difficulty in playing. The Flicker was not difficult of playing. Thus is shown the fact that the height of the bridge must depend upon the height of the finger- board, and the height of the finger-board should depend upon the height of plate- arching. In the case of the Flicker violin, the flat under surface of the finger-board, resting upon a high neck, operated to direct sound-waves towards the neck, and apparently, such is only the practical plan for violins having the greatest height of arch. The Flicker exits were unusually large, and appear- ed to be unusually close together. The spring of arch began directly within the purfing. To these facts do I attribute the unusual tone-power of this violin in comparison with the tone-power of other violins in its class. It is evident that increasing width of the finger- board increases the number of sound-waves reflect- ed'back to the sounding-board; therefore, such in- crease in width operates to diminish tone- power. The finger-board is a necessity, yet, with best'possible adjustment, the finger-board operates todiminish i violin tone-power. Thus, the width, the reflecting, under-surface line, and the distance of such 'line "from the sounding-board, are potenti- alities of the finger-board demanding the closest attention from the violin tone-regulator. Lack

168

\section*{VIOLIN TONE-PECULIARITIES.}
of such attention brings penalties. In the matter of distance between the sounding-board and under- surface line of the finger-board, hard-and-fast rules cannot apply only upon the condition that the arching and graduation of the sounding-board govern the rule. The desired height of the strings from the fing- er-board, and the position of the bridge upon the sounding-board govern the height of the bridge. Height of strings above the finger-board may quite safely be allowed some latitude to accommodate different taste. My own choice for this height is -b> clear, throughout. Some prefer a greater height for the G, and less height for the E, while some prefer a greater height throughout. I once knew a proficient violinist, having an unusual length of hand and fingers, who placed the strings i inch above the lower end of the finger-board. I asked whyfore? He replied, "Because of two reasons; one being the fact that my violin yields greater power of tone; the other being the fact that pizzi- cato tones are clearer." I do not find his first reason to hold good in all cases; but, on the con- trary, I have known tone-power to be weakened by thus raising the strings. Clear pizzicato tone is very desirable. The snap- ping pizzicato tone due to string-oscillation receiv- ing interference from the finger-board, is some- thing inadmissable. To secure clear pizzicato tone, while yet leaving the strings at or less, from , the lower end of the finger-board, I resort to the following treatment of the upper finger-board

169

VIOLIN TONE-PECULIARITIES. surface. With a small block-plane I dress away this surface until its line presents a slight curve from end to end. The least wood is removed from beneath the E, and the greater amount from be- neath the G, and because of the fact that amplitude of string-oscillation becomes wider from E to G. The lowest point in this curve, determined by the straight-edge, is not placed higher than c in alt, and because of the fact that shortening a string oper- ates to diminish amplitude of its oscillilation. Further consideration of the bridge will be re- sumed at the next hour.

170

\section*{VIOLIN TONE-PECULIARITIES.}
LECTURE XI. GENTLEMEN: Diminution of volume and inten- sity of violin tone by the bridge is resumed. (b) Thickness of the bridge may be so great as to become a potent factor in diminishing both vol- ume and intensity of tone. A demonstration for this fact is very easy of accomplishment. Although this fact comes within the daily observation of vio- lin students, yet, because only a small minority of such students trouble themselves with the philoso- phy involved in tone-modifiers, and, because ac- quaintance with certain physical laws is of value to the violin tone-regulator, therefore it seems ad- visable to give consideration to such laws. Certain physical laws are valuable to violin tone insofar as human ingenuity can apply them. In such appli- cation lies a difficulty. In all my long-time appli- cation there has been but a single result of marked value to violin tone following a priori reasoning. I find the problem of bridge-thickness to be com- plicated with four inevitable factors: (a) Height of finger-board (b) Height of arching. (c) Diameter of strings. (d) Density of fiber. (a) As previously shown, the height of the bridge is governed by the height of finger-board, yet, this statement is subject to modification. Be- cause two violins have precisely similiar height of finger-board, it does not follow, as a necessity, that the bridges have precisely similiar height. This fact is due to dissimilarity in height of arching

171

\section*{VIOLIN TONE-PECULIARITIES.}
at the position of the bridge. Doubtless many, if not all violin students have observed that some makers place the highest point in the arching at the position of the bridge, while other makers place the highest point upwards, or forwards from the po- sition of the bridge. It is evident that this dissim- ilarity necessitates dissimilarity in bridge-height. [Although not here in point, yet, fear of omis- sion makes me call attention to placing the highest point of the arch, (longitudinal arch) at position of the bridge, Within my observation, placing the highest point of arching at position of bridge, op- erates to increase tone-power' From my view point, the reasons for such increase are two in number, and thus: First: Diminishing height of bridge permits diminution of bridge- thickness; hence, diminution in the muting effects from a greater mass of wood in the bridge. Second: Placing the highest point in plate arching forwards of, or upwards from position of the bridge operates to diminish the amplitude of sounding-board oscil- lation beneath the strings; hence, diminished pow- er of tone.] It is evident that thickness of the bridge may operate to mute violin tone. It is also evident that diminution in bridge-height permits diminution in bridge-thickness. It is observable that diminution of bridge-thickness, down to the point of sustain- ing string-pressure without bending, operates to deliver upon the sounding-board a greater force from blows of the strings, Hence, both diminish- ed height and diminished thickness of bridge op-

172

\section*{VIOLIN TONE-PECULIARITIES.}
erate to diminish tone-power. [At a previous hour it was shown that proximity of vibrating bodies, susceptible to identical force, operates to augment sympathetic action. For this reason, loss of sympathetic action between the strings and that part of the sounding-board be- neath the strings might be considered as one of the modifiers operating to diminish volume and inten- sity of violin tone; but because the list of such tone- modifiers already reaches the number 12, and, be- cause of failure to find two more such modifiers, therefore I decline to permit sympathetic action, (with its number 13) a chance to "hoodoo" my work. (c) Diameter of violin strings largely, (not wholly) governs downward pressure upon the ' bridge. It is apparent that the smaller string-dia- meter causes less downward pressure upon the bridge by reason of less tension demanded in tun- ing. Therefore, to produce maximum tone-power from smaller strings requires diminution of bridge- thickness proportionate to diminution of string- diameter. The exception to string-diameter gov- erning downward pressure upon the bridge lies in plate-arching. It is evident that, with no arching whatever, downward pressure by the strings is at the minimum; and, because of the fact that the strings are at the least practical distance above a straight line from saddle to nut. At such straight line, the downward pressure being zero, it follows that every degree of bridge-height above this line operates to increase downward string-pressure.

\section*{VIOLIN TONE-PECULIARITIES.}
Therefore the higher the plate-arching, the greater is downward string-pressure upon the bridge. Therefore arching becomes a factor in determining bridge-thickness. Thus is shown the fact that hard- and fast rules cannot apply to violin bridge-thick- ness, (c) The span between bridge-pedestals, or shanks, while not being a potent factor in diminish- ing tone-power, yet it is a factor demanding atten- tion. By being either too great, or too small, this span operates to diminish amplitude of sounding- board oscillation, and therefore diminishes tone- power. There are two factors governing bridge-span: Position of bar. Distance between exits. According with my observation, failure in plac- ing the center of the left pedestal over the center of bar-thickness operates to diminish D-string tone- power. Again, placing the center of the right pedestal over sounding-board fibers cut off by the right'exit operates to diminish E-string tone-pow- er. The bridge, without pedestals whatever, possesses interest, thus: In those cases wherein a hollow tone of great volume, but of weakened in- tensity, caused by too great reduction of sounding- board rigidity in its central area, the bridge with- out pedestals, carefully fitted to lateral curve of the sounding-board, operates to augment A and D- string tone-power; not greatly, but perceptibly. Upon that sounding-board possessing sufficient rigidity in its central area, I do not find that the bridge without pedestals operates to either aug-

174

\section*{VIOLIN TONE-PECULIARITIES.}
ment or diminish tone-power of any string. In this respect I find the bridge without pedestals to operate precisely as the re-enforce block glued transversely upon the inner surface of the sound- ing-board at the position of the bridge. [In experimental work, I have placed such re- enforce block upon 17 different sounding-boards having varying degrees of rigidity in their central areas. The .benefit to tone therefrom was only manifested upon such sounding-boards as were reduced in thickness sufficient to cause weakened A and D-tone. Upon such sounding-boards as pro- duced strong A and D-tone, the re-enforce block neither operated to diminish nor augment tone- power of any string.] Density of fiber is a powerful factor in the bridge. Either softest fiber, or densest fiber operates to di- minish tone-power. Without doubt, this fact is due to failure in transmission of force from strings to the sounding-board. In the case of the violin, it is apparent that force from the strings can be only communicated to the sounding-board by the vibratory action of connecting media, as the bridge, and air. Such transmission is not confined to the bridge by any means. Sympathetic action between strings and sounding-board is wholly due to the presence of air as a connecting medium. The po- tency in sympathetic action I find to be eminently worthy of consideration by the violin student. I know of no easier demonstration for such potency than as follows: Remove the bridge, and, over the tail-block, place a bridge of sufficient height to

175

\section*{VIOLIN TONE-PECULIARITIES.}
maintain the strings at their usual height above the finger-board. [To make a precise test for the potency in sym- pathetic action, it is a necessary condition that the strings be not attached to the violin. In either method, tone-power from sympathetic action is surprisingly great when sounding-board rigidity corresponds to string-diameters, in other words, to string-force; but, when such rigidity is too great, then the power of sympathetic tone is diminished, and diminished because maximum sympathetic act- ion between two contiguous, vibrating bodies de- mands equal susceptibility to force. Such sympa- thetic tone is also diminished by area of the exits.] The greatest density in any bridge I have tried is found in one made of bone. The density in this bridge modifies tone in a peculiar manner; both volume and intensity of tone being greatly dimin- ished, while the little tone remaining is remarkably thin> and of an excruciatingly stinging quality. The bridge of softest fiber to which I have given trial, is selected from soft Michigan pine. This bridge operates to diminish both volume and inten- sity of tone, but, it also operates to diminish dis- agreeable quality of tone. Upon the violin of noisy tone, I have not observed failure of improvement in disagreeable tone-quality to follow employment of this soft wood 1 for the bridge. The thickness of such bridge cannot be reduced to the same degree as the maple bridge because of greater ease in bending under string-pressure. Asperity of tone may be greatly reduced by employing such soft

176

\section*{VIOLIN TONE-PECULIARITIES.}
wood for both bridge and post. Scroll-work upon the bridge has a more import- ant mission than mere ornament. It is my exper- ience that the "whole," or "solid" bridge, of any wood whatever, operates to diminish tone-power. Could thickness in the "solid" bridge be diminish- ed until its vibratory action becomes equally sus- ceptible to force as the scroll-cut bridge, then scroll-work would become merely ornamental. But, such diminution of bridge-thickness is imprac- ticable because downward pressure of the strings operates to bend the "solid" bridge when thus re- duced in thickness. It is a fact that bending the bridge operates to diminish the amplitude of its os- cillation; hence bending the bridge operates to di- minish tone-power, and, in precisely the same way as bending the sounding board. Obviously, thick- ness of the bridge must be great enough to hold the bridge erect under string-pressure. Because string-pressure varies with height of plate-arching and diameters of strings, and, because different samples of bridge-wood present differing degrees of rigidity, therefore hard-and-fast rules for bridge-thickness cannot apply. In this matter, I know of no successful rule other than the usual " rule, ' 'Cut-and-try. Experience demonstrates that the "solid" bridge, thick enough to stand erect under 25 to 28 pounds of string-pressure, possesses mass sufficient to appreciably mute the tone. Experience also demonstrates that some part of such mass may be safely removed from cer- tain parts of the bridge without diminishing rigid-

177

\section*{VIOLIN TONE-PECULIARITIES.}
ity to the point of bending. Such diminution in mass is secured by the famil- iar scroll- work: The location and extent of scroll-work upon the bridge by no means should be left to chance; and because mal-position and extent of such work may operate to cause unevenness in tone-power. Thus, after every other minute detail in construction, op- erating to produce evenness of tone-power, has re- ceived the limit of attention, yet, mal-position and extent of central scroll-work on the bridge may re- main to defeat the most skillful violin builder who- ever felt the impulse of ambition. As I hold this bridge up to your view, you ob- serve the location and extent of the central scroll, and also the scroll-work at either end. Near to either end, you observe an isthmus, or narrow neck between the central and end scrolls. As you ob- serve, these necks connect the upper and lower- halves of the bridge. It is apparent that all vibra- tory action in the bridge, aroused by action of the strings, must travel downwards through these nar- row necks. It is evident that diminution of these necks operates to modify transmission of vibratory action. It is also evident that inequality in the di- mensions of these necks operates to cause unequal susceptibility to force. It is also evident that force in violin strings varies as their diameter and weight. Hence, it becomes obvious that equality in the dimensions of these necks operates to dimin- ish tone-power of the smaller strings. Therefore, for production of evenness of tone-power, dimen-

178

\section*{VIOLIN TONE-PECULIARITIES.}
sions of the neck beneath the A and E-strings should be less than dimensions of the neck beneath the G and D-strings; and, such diminution in mass should be in the same ratio as the diminished diam- eters of the strings. Here is a dozen high grade bridges. With a machinist's rule I determine dimensions of their isthmi between central and end scroll-work. I find thickness of those isthmi to be i; but, in their width, I find variation, from equality, to be 1-16 inch. It is plainly apparent that such variation in mass may be turned to good account in the work of se- curing evenness of tone-power. Thus: Placing the larger isthmus beneath G and D-strings operat- es to equalize transmission of force. It is also ap- parent that diminishing thickness of these isthmi permits the bridge to bend under string-pressure; but, diminishing their width cannot cause such dis- astrous result. It is also apparent that the dimen- sions of those isthmi should be governed by string- pressure; and, because such pressure varies with varying degrees of plate- arching and string-diam- eter, therefore, for such dimensions, hard-and- fast rules cannot apply. (f ) Maturity of bridge-wood, (maturity of tree before jbeing felled) is a powerful factor in the list of tone-modifiers. Such maturity is of equal im- portance with maturity in sounding-board wood. In the absence of maturity from either lies defeat of tone quality. The reasons for taking sounding- board wood from that part of the log between heart-wood and sap-wood apply equally to bridge- wood.

179

\section*{VIOLIN TONE-PECULIARITIES.}
LECTURE XII. GENTLEMEN: (g) Position of bridge is an im- portant tone-modifier. So important is it that mal- position may utterly annihilate "richness" of vio- lin tone, even when all other factors are at their best. I recall that case of "wolf" which was com- pletely cured by position of bridge and post. I also recall the fact that harmonics a bassa, or resultant tones of the text-books, tones assisting in produc- ing "rich" violin tone, will become audible only when length of strings, (from bow to nut, not from bridge to nut) and the active length of sound- ing-board are equal. Thus: Because the bow practically shortens the strings by one inch, and because the greatest length of sounding-board ac- tivity, productive of audible sound, equals 12 inch- es, therefore, for production of the "rich" violin tone, the length of strings, from bridge to ' nut, must equal 13 inches. But, in tone regulation work, the above rule cannot be interpreted as a hard-and-fast rule. Within my observation, there is one point, and but one, for that position of the bridge yielding greatest richness of tone; and, by no means am I able to precisely pre-determine such position. Experiment on each sounding-board must be depended upon for precise determination of the best bridge-position. I close discussion of the violin bridge with the unqualified statement that all bridges of immature wood, (always sap-wood) are best disposed of as firewood. (8) The finger-board problem comes next in the

180

\section*{VIOLIN TONE-PECULIARITIES.}
list of factors operating to diminish volume and in- tensity of violin tone. Because of my inability to solve the problem of the bridge until after solving the problem of the finger-board, therefore the fing- er-board was given priority in presentation. At this moment I cannot call up anything further to say upon this tone-modifier. (9) The strings are a very important factor in diminishing violin tone-power. Their importance is clearly shown by the fact that sounding-board rigidity must be governed by string-diameters and weight. In my statement concerning violin strings, the wire string receives no consideration further than that it is an evil made necessay by evil situa- tions; that is, situations wherein excess of water vapor operates to quickly ruin gut strings. In all cases where sounding-board rigidity is determined by large strings, then substitution of smaller strings operates to diminish tone-power. It is my obser- vation that in all cases wherein sounding-board rigidity is precisely reduced to correspond with force in guage-2 strings, then substitution of larger strings operates to augment volume of tone while diminishing intensity of tone. Thus, at nearby distances, tone is greater; but is a failure at long distances. I find a satisfactory explana- tion for this phenomenon to be a matter of difficul- ty. As an aid to such explanation, I call attention to the fact that larger strings, played in double- stops, cause violent trembling of such sounding- board. It seems reasonable to suppose that trem- bling of the violin may be so great as to weaken

181

\section*{VIOLIN TONE-PECULIARITIES.}
tone in view of the fact that violent trembling of the horn operates not only to weaken tone, but also, perceptibly to lower tone-pitch. It is a fact of more or less frequent observation that the over- forced tone from any horn in a brass band is out of tune with other horns, and out of tune because of being- lowered in pitch. It seems to me that such lowered pitch is caused by the violent trembling of the instrument. That such tone is weakened in intensity is clearly perceptible. Again, in orches- tra ensemble, tone from the violently trembling violin may be totally annihilated by sound-waves from harmony instruments. For this fact there appears no reasoning so plausible as diminished in- tensity in the tone of such trembling violin. Again, in the chorus, when any member indulges in vio- lent tremolo, it is at once apparent that his, or her, (too often "her," more's the pity) voice is not only weakened, but, is so much out of tune as to be- come sickening to the musically trained listener. Again, strings of too great diameter operate to ac- centuate noisy tone-quality. Thus, whatever may be the correct explanation, the fact remains that employment of strings., having too great force for sounding-board rigidity, or for back-plate rigidity either, is but inviting disaster. From my view point, 'tis safer to err in the opposite direction be- cause smaller strings operate to diminish noisy tone-quality. The "twist" and number of "strands" in the strings may operate to diminish both volume and intensity of tone. Thus, the string of but a single

182

\section*{VIOLIN TONE-PECULIARITIES.}
strand operates to diminish tone-power. The rea- ' son is obvious. In such string, flexibility being at the minumum, therefore rapidity in winding and unwinding under bow-pressure is greatly diminish- ed; therefore diminished force in blows delivered upon the bridge must follow. Flexibility in violin strings operates to augment rapidity in winding and unwinding; hence the flexible string delivers a blow of greater force. Weight of strings may also diminish tone-power. To be equal in weight, the silk string must be given increased diameter, while the steel string must be given diminished diameter. In all cases evenness of tone is diminished by dis- proportionate diameter and weight of strings. As a violin tone-modifier, the subject of strings pos- sesses vast interest to the student of tone peculiar- ities. The tone of the finest violin ever built may be easily ruined by the strings employed. In my observation, the choice in metal for winding the G- string should be governed by the tone-peculiarities inevitable to each individual violin. In equal quan- tity, copper, silver and gold vary in weight. Whenever I have found a copper- wound string pre- cisely fitting the tone-peculiarities of a certain vio- lin, then substitution of either silver, or gold- wound G has proven to be a disappointment. In precise adjustment of the G-string, it is my method to make trial of such strings as have differ- ent diameters. After determining that diameter giving best results in double-stops, then I select one of slightly greater diameter and give it time to stretch. Then, while in position, and in tune, I

183

\section*{VIOLIN TONE-PECULIARITIES.}
proceed to diminish both weight and diameter with an inch-square piece of emery-cloth. With precis- ion as an object, the emery-cloth, folded once around the string, and firmly held by the thumb and fing- er, is made to travel from bridge to nut, then, slightly turned and made to travel back to the bridge. This operation is frequently interrupted by application of the bow, and for the purpose of knowing the moment when such reduction in diam- eter and weight has reached the desired point. When carefully done, this method largely removes certain rough inequalities of tone inevitable to all new G-strings. When this method is applied upon a G-string having correct diameter, the result is disastrous. Considering the vast value in beauti- ful G-string tone-quality, no amount of work thereon is too great. [Possibly I may be liable to the charge of fre- quent repetition; but I'll take such risk by again stating that beautiful double-stop tones are the sal- vation of the violin soloist. With salvation in view, the G-string must not be offensively prominent. It is possible that I may be hard to please in the matter of violin tone-quality; but, whatever I say upon this point is said only as the opinion of an in- dividual, and said with the acknowledgement that my opinion may be in error. In life I know of no fact more potent than the fact that musical taste varies as the number of people. In my opinion, the tone from a set of steel violin strings is only worthy of anathema. (10) That the condition of interior violin-surfaces

184

\section*{VIOLIN TONE-PECULIARITIES.}
may operate to diminish both volume and intensity of tone requires only a thought. Did I recommend spreading a carpet over the inner surface of the violin back, every reader would doubtless pronounce judgement upon me thus: ' 'That Castle is a fool, " or worse yet, "Castle is a fiddle crank." The worst of it is, the "coat would fit." Doubtless many readers think ' 'the coat fits" snug enough when I say that both volume and intensity of violin tone are augmented by a perfectly smooth interior surface. I do say that, and say it with all the earnestness at my command. This problem in vi- olin tone is solved to my satisfaction. If any good quality of tone were injured by the perfectly smooth interior surface, then by no means would this problem be solved to my satisfaction. The question of solving this problem to your satisfac- tion rests entirely with yourselves. I'd never give you the details for interior-surface work were I "out just for health." On the contrary, I'd keep to myself the secret of changing \$5-fiddles into \$100-violins. If necessary, I can give the names of several quite competent violinists whose \$3-fid- dles, (wholesale) are now valued by them at sev- eral hundred dollars. Do not understand me as claiming all this change in tone-value to be due to a perfectly smooth interior-surface; but, you may understand me to say that, without such permanent- ly smooth interior-surface, the owners would not have attached those high prices. I can give the name of one such owner, living within 30-minutes ride of one of the world's largest violin markets,

185

\section*{VIOLIN TONE-PECULIARITIES.}
who positively declined to attach a price to his originally \$10-flddle. "Did I give details to those owners?" "Certainly not." "Why not?" "Because of the very prejudice you entertain at this moment." Here's a "pointer" intended for your benefit: When any person guarantees the tone of his violin to "carry" 1250 measured feet out in the open and under identical meteoric conditions, and in identi- cal hours of the day, and with identical precaution against the presence of sound-wave reflectors that have been heretofore described, just trouble your- self to ascertain whether or not, the interior sur- face of that violin is permanently and perfectly smooth. 'Again, in all old violins which have not been "cleaned out," and many, very many used violins not yet old in years, I assure you of finding their interior surface covered with a "carpet." Such carpet will be found composed of wood-fiber, wood- dust, and dirt. Here is your chance to test aug- mentation in volume and intensity of tone without injury to any tone-quality, of which I have given assurance in these pages. To all violin users desir- ing-a violin possessing maximum tone-power, I rec- ommend careful application of the details for inter- ior surface*protection, together with other details herein submitted. With careful, painstaking work, I am confident of your success. (11) Area and position of the exits are the most

186

\section*{VIOLIN TONE-PECULIARITIES.}
powerful factors operating to diminish volume and intensity of violin tone. Area of exits alone, is more potent than all other factors combined. Startling as this statement may be, yet its correctness is of easy demonstration. Without exits whatever, the finest violins ever built will not, nor cannot yield more tone than the broom- stick fiddle. In such situation, the strings are left with nothing but unconfined air upon which to ex- pend their force. Hence, instead of concentration of sound-wave lines of travel, there is the widest possible dispersion of those lines. Again, in such situation, string force receives augmentation neith- er from direct sounding-board action nor from sym- pathetic action. Without exits, the sounding- board is motionless. Its power is not sufficient to compress confined air. Without exits, the violin descends to the broom-stick level. Hence, without exits without violin. But, a minute opening in the walls operates to permit some action of the sounding-board. In- stantly there follows perceptible augmentation of tone-power; and such augmentation follows placing the opening upon the- back, upon the ribs, upon the sounding-board, wherever you please, but such augmentation does not follow in equal degrees. Those curved walls operate to direct sound-wave lines of travel. As you look at a violin of good model, you instantly perceive that those interior curved walls cannot direct the bounding ball towards the ribs. Were those ribs of glass, no bounding ball would ever break them. Wherever

187

\section*{VIOLIN TONE-PECULIARITIES.}
the bounding ball goes, there, also, the sound-wave movement must go. But, there's a difference in the action of a single ball and countless number of balls touching each other. Air molecules, being balls possessing immense elastic energy, and not moving away from that position occupied at the instant of receiving a blow, operate to produce sound-waves by communicating elastic energy from one to the other; and such communication contin- ues until force in the original blow becomes ex- hausted. It is evident that increasing force in the original blow operates to increase the distance traveled by sound-wave movement. It is also evi- dent that sound-wave movement may be either dis- persed or concentrated; also, that concentration of such movement operates to augment both volume and intensity of tone, but, in unequal degrees; in- tensity being augmented in the greater degree. It is apparent that interior walls of the violin body not only prevent dispersion of sound-wave move- ment, but also because of longitudinal and trans- verse arching of the plates, operate to concentrate such movement. It is apparent that the point of sound-wave concentration within the violin must vary with varying degrees of plate-arching. To precisely determine such varying points is the des- pair of scientists. It is evident that placing the exits at points of greatest sound-wave concentra- tion becomes a powerful factor in the production of maximum tone-power; also, that placing the exits at a distance from such points of concentration op- erates to diminish tone-power.

188

VIOLIN TONE-PECULIARITIES. "How may we find those points? "Don't ask the scientist. " Not more can he tell you than Sam Jones can tell you where heaven is. By reason of experience alone, the violin builder will cut out the exits obliquely across the sound- ing-board and trust to luck for hitting those points. Thus, obliquity in position of the exits becomes of immense value to violin tone power. Could those points of greatest sound-wave concentration be pre-determined, then exits in parallelogram form might be employed to augment tone-power. As previously shown, the exit of small area per- mits but a limited degree of sounding-board action, therefore it follows that increasing the area of ex- its operates to permit increased amplitude of sound- ing-board oscillation; hence an increased force of blow upon contained air; hence an increased tone- power. We are now arrived at a point in violin construct- ion possessing intense interest to two classes of violin users; the one desiring quality of tone with moderate volume; the other desiring great volume of tone regardless of quality: From my point of view, interest in violin exits equals interest in sounding-board wood and model of plate-arching. Further consideration of the exits will be con- tinued at the next hour.

189

\section*{VIOLIN TONE-PECULIARITIES.}
LECTURE XIII. GENTLEMEN: In resuming presentation of violin exits as tone-modifiers, I take occasion to repeat the fact that the exits are more potent than all other factors combined. Therefore, to violin exits attaches intense interest. "Without exits, with- D\textasciicircum\{\}' violin," is equally truthful as, "without sun, with- out light." From pole to pole, all animate objects love the sunshine. From southern habitable limits to northern limits, humanity loves sweet sounds. All over the habitable world, the area and position of violin exits contributes to the sum of human happiness. In diminished area of exits the aesthet- ic violin lover may find solace. In enlarged area of exits, the lover of great volume may find enjoy- ment. In power to command beautiful tone-quality, vi- olin exits stand supreme and alone. Their power to suppress "noise" surpasses all other tone-modi- fiers combined. Blessed are they! Mr. Builder, in cutting out those exits, I pray thee to use occasionly that keen blade with a sparing hand. Sweet Music, keenly watching your work, will sweetly sing your praises in the ear of your cus- tomer. There are yet violin users who will give you an ounce of gold for every pennyweight of wood between the small and the large exit. [There is a large class of violin users preferring quality of violin tone above mere quantity of tone. Without hesitation, I confess membership in such class; but, do not understand me as condemning

190

\section*{VIOLIN TONE-PECULIARITIES.}
quantity of violin tone. On the contrary, I hold to the belief that quantity of violin tone is a necessi- ty in certain situations; and, in such situations, quantity of tone is equally valuable with a gun "out west;" either being a necessary evil for sal- vation. What I do condemn is the employment of "big tone" in all situations and upon all occasions. Next to the noisy violin, the violin of "big tone" becomes offensive when employed in other situa- tions than where it belongs; that is, in the large auditoria. Offensiveness of tone is a serious mat- ter, whether or not offense is due to too much or too little tone. It is my desire to assist in prevent- ing employment of the violin in situations arousing contempt for both performer and instrument. Thus: To carry an old violin, whose tone has gone down into dotage, into the larger auditoria and at- tempt to force its once willing tone to farthest expectant ear is not only an act of inexcusable idi- ocy, but, is also a display of heartless cruelty. Pos- sibly the later withdrawal of patronage may grow to become an efficient method for preventing con- tinuance of such distressing displays.] The most wonder-exciting phenomena connected with violin tone are due to area and position of the exits. Even the constancy of these phenomena ex- cites wonderment. There are seven of these phe- nomena, and four out of the seven do not depend upon inherent quality of wood for existence. Nor hard fiber, nor soft fiber nor no fiber at all exert any influence whatever upon those four phenomena. The existence of these four phenomena depends

191

\section*{VIOLIN TONE-PECULIARITIES.}
wholly upon air. The seven phenomena connected with violin ex- its are: (a) Position of exits may diminish or increase tone-power. (b) Increasing area of exits augments volume and diminishes intensity of tone. (c) Diminishing area of exits augments intensi- ty and diminishes volume of tone. (d) Diminishing area of exits lowers tone-pitch. (e) Enlarging area of exits raises tone-pitch. (f) Diminishing area of exits diminishes noisy

tone-quality. (g) Enlarging area of exits accentuates noisy tone-quality. In view of this display of facts, there is little wonder that great interest attaches to violin exits. Because of not having a satisfactory explanation for the phenomenon in raising and lowering tone- pitch by enlarging or diminishing area of exits, therefore none is offered; but, in advance, I offer thanks to anyone for such satisfactory explanation. (a) Position of exits may diminish or increase tone-power. This fact is susceptible of conclusive demonstra- tion. Thus: Placing the exits in the ribs of the middle bout, other conditions remaining as usual, operates to diminish tone-power. The reason for such diminution is apparent; the exits thus placed are not in the line of sound-wave concentration. This conclusion is proven by closing such exits, and placing others in the usual position. The

192

\section*{VIOLIN TONE-PECULIARITIES.}
striking increase in tone-power following such change in position is the measure of sound-wave concentration in that particular violin; but, is not a measure of such concentration for another violin having a greater or less height of plate-arching. Again, placing the exits nearer to the edges of the sounding-board operates to diminish tone-power, but, the degree of diminution varies with the height of arching; the lower the arch, the less di- minution in tone-power. It is obvious that increas- ing height of arching operates to increase sound- wave concentration in the direction of the center join of the plates; therefore, the higher the arch, and the nearer the exits are placed to the center join, the greater is tone-power. In practice, I find a variation of i in distance from center join to the exits operates to change tone-power when height of arch equals f or more; below f, the change in tone-power is less marked. In this matter, the pattern for exits should vary in width of curve at the ends as the degree of tone-power desired. Thus, with diminished width of curve, the exit may be placed nearer the center- join; width of curve increased, places the exit at a greater distance from the center-join. For position of exits, hard- and-fast rules cannot apply. (b) Increasing area of exits augments volume and diminishes intensity of tone. My explanation for this phenomenon begins upon a well made vio- lin, of good model and good wood, and having no exits whatever, but having other conditions as usual.

193

\section*{VIOLIN TONE-PECULIARITIES.}
In this condition, the plates remain motionless and toneless under vigorous bow pressure. Its ad- mirable sounding-board cannot act. Contained air will not permit it to act. Elasticity in air molecules within this violin body is a resisting force vastly greater than force in the strings. Only by per- mitting escape of part of this resisting energy can this sounding-board be excited to action. For the purpose of observing the effect upon this sounding- board, I shall permit escape of such resisting ener- gy in gradually increased degrees. If my reason- ing is sound, then sounding-board activity will increase as escape of resisting force increases. With this object in view, I outline the exits in their usual position, and cut out the smaller round at their upper extremities. Application of the bow now produces sound, but it is a sound of small vol- ume, of low pitch, and of marked intensity as com- pared with volume. It is clear at once that small volume of this sound is due to escape of but a small amount of molecular, elastic energy, but reason for the intensity of this tone is not so clear, while rea- son for its low pitch drives me into the region of mere supposition. I know the supposition of phil- osophers upon this point, but supposition lacks much of being satisfaction. That the varying lengths of air columns confined yield sound of equally varying degrees of pitch is amply demon- strated in organ pipes, horns, steam whistles, and all wind instruments; but, after enlarging these exits to one half usual area, how may we account for the greatly raised tone-pitch following? It is

194

\section*{VIOLIN TONE-PECULIARITIES.}
clear that the greater volume of sound following such enlargement is due to greater escape of mo- lecular, elastic energy; also, that such increased escape of energy, by diminishing resistance, per- mits wider amplitude of sounding-board activity; hence greater force in the blow delivered upon con- tained air; hence greater volume of tone. After enlarging these exits to their usual area, application of the bow determines another great addition to volume of tone. Again, after enlarg- ing these exits beyond their usual area, volume of tone is yet increased, tone-pitch is raised, but, in- tensity of tone is diminished, and noisy tone-qual- ity is accentuated. How far such increase in area of exits may be extended without complete ruin to tone-value is one of the demonstrations I have not made. It is evident that such increase has its practical limits. We know that total absence of confining walls per- mits dispersion of sound-waves equally in all direc- tions; that such dispersion operates to diminish in- tensity of sound to the minimum. It is apparent that area of the exits may 'be so great as to permit an amount of dispersion greatly dimiminishing intensity of tone. It is a well known fact in physics that the greatest volume of sound occurs in open air where no confining walls are present. In view of these facts, it is evident that diminishing the area of exits operates to diminish volume of tone, to increase intensity of tone, and to diminish noise. In the work of artistic violin construction, I know of no factor possessing greater importance than

195

\section*{VIOLIN TONE-PECULIARITIES.}
area and position of the exits. Considering good material and good model to be present, then area and position of the exits must be relied upon to satisfy the varying tastes of violin users. Thus, that builder who only desires to please the taste for great volume of tone, will place exits of large area as near to the center join as is practical; and, that builder who only desires to please the taste for diminished volume of tone, increased intensity of tone, and diminished noisy tone-quality, will place exits of small area at a greater distance from the center join. From my view-point, it seems but wisdom to be prepared for all varieties of tone- taste; therefore in the work of tone-regulation up- on violins intended for such as prefer quality be- fore volume of tone, it is necessary to cut the exits too small at first, and to increase their area only after application of the bow. In my experience, it is a fact standing out in strong relief that hard- and-fast rules for the area of exits cannot apply. Thus: If a sweet, intense tone is desired, then the exit of large area operates to defeat intention. Again, differing degrees of arching operate to cause differing lines of sound-wave travel; therefore, po- sition of exits must approach to, or recede from the center join as the variations in height of arch- ing; therefore, hard-and-fast rules for position of exits cannot apply. In this matter, as previously stated, the violin builder must rely wholly upon experience and ob- servation. There is no alternative. In this mat- ter, that colossal thing called "science," isashelp-

196

\section*{VIOLIN TONE-PECULIARITIES.}
less as an infant. In the case of the bar- rel-shape violin, we may rest assured that the exits are not in the line of sound-wave concentration. Such assurance is based upon the fact of feeble tone-power invariably existing in violins of such model. The tone-power from this rotund violin is but the tone-power of an infant; whereas, the tone-jjower from this flat-model violin, its every line a line of beauty, merely suggesting rotundity, sets every beam of the house into a rollicking two- step. Without a shadow of doubt, the exits of the latter violin are in the line of sound-wave concen- tration. I believe that continued trial might de- termine the precise distance of the exits from the center-join for each degree of arching. With the intention of producing maximum tone-power, such figures would possess value. Oft repeated demonstrations have firmly estab- lished the fact that increasing the area of exits op- erates to raise tone-pitch. I confess this phenome- non seems paradoxical. I cannot divest myself of an impulse to think that the contrary result to tone- pitch should follow enlarging the exits. [Here is a demonstration for the fact that pre- conceived ideas may lead one into error. Often upon arising in the morning at some strange place for which we had formed preconceived ideas of the different points of the compass, the sun seems not to rise in the east. It may appear to rise in the north, and even in the west; and, it may obstinate- ly refuse to rise in the east so long as we remain at that point. Thus am I affected by the phenom-

197

\section*{VIOLIN TONE-PECULIARITIES.}
enon of raising tone-pitch by enlarging area of exits. Again I offer thanks in advance for a satis- factory explanation for thus raising violin tone- pitch.] The small amount of such enlargement necessary to increase volume of tone, and, raise tone-pitch is a matter for profound surprise. Removal of but a single shaving from the inner edge of an exit is sufficient to change these two qualities of tone. The philosophy for diminishing or accenuating "noise" in violin tone by diminishing or increasing area of exit is not difficult of explanation. The cause for "noisy" violin tone, as previously shown, is due to small areas of greater thickness of wood located in sound-producing parts of the sounding-board; and, because of such limited tone- producing area, it is apparent that the tone there- from possesses but feeble power. It is also appar- ent that any agent operating to diminish violin tone-power, might operate to annihiliate the feeble noise- wave. Hence, diminishing area of the exits augments sweetness of violin tone; also, when the cause for "noise" exists, increasing area of the exits operates to accenuate such noisy sound. Feebleness of the noise- wave is easy of demons- tration. Take the noisiest violin out in the open and apply thereto the utmost vigor of the bow, while the listener retires to a distance. As dis- tance increases, noise diminishes, while musical sound yet remains distinct. Position of exits ex- erts considerable influence in producing that rum-

198

\section*{VIOLIN TONE-PECULIARITIES.}
bling character in violin-tone so aptly described as "tone-all-inside." Although faulty interior lines of the plates are the chief factors in producing this character of tone, yet, position of the exits may add to or subtract from such rumbling sound; but, in cases where such rumbling sound is marked, the position of exits cannot be relied upon to effect complete cure. Indeed, there seems to be no cure for serious cases other than building new plates having correct interior lines; that is, lines which continually, and with precise regularity, direct sound-wave movement in the direction of the exits instead of directing such movement away from the exits. To produce the maximum tone-power, when interior lines operate to direct sound-wave move- ment from the exits, is an utter impossibility; also, with such lines absolutely perfect, placing the exits away from lines of sound-wave travel defeats max- imum tone-power. In the production of maximum tone-power, interior lines of the plates become prodigious factors; whereas, exterior lines exert no influence whatever. Therefore, in producing maximum tone-power, it becomes necessary first to establish interior lines of the plates, and thereafter reduce thicknesses by work on exterior lines. Thus, the highest point in the longitudinal arch, viewed from the interior, may be placed at the po- sition of the bridge; thus, sound-wave movement cannot be directed away from the exits when exits are placed at the nearest practical distance to the center- join. (12) Depth of ribs is the last factor in the list of

199

\section*{VIOLIN TONE-PECULIARITIES.}
tone-modifiers for consideration. The chief effect upon violin tone from depth of ribs is manifested upon tone pitch; yet, the effect of this factor upon volume, intensity, and brilliance of tone is worthy of consideration. Other dimensions remaining equal, diminishing depth of ribs operates to demin- ish volume of tone, while intensity and brilliance of tone are increased. In experiment upon this factor I have diminished depth of ribs from 1 and i inches, down, by degrees of 1-16, to a depth equalling I inch; and from this depth, have rebuilt such ribs, by additions of i, to the depth of 1 and f inches. Such changes in depth, being given to one particular violin, afford conclusive proof that shortening length of perpendicular air columns within the violin body operates to raise tone-pitch, and vice versa, to lower tone-pitch. In diminish- ing depth of ribs to I inch, violin tone undergoes remarkable changes in character; volume of tone is greatly diminished; tone-pitch is greatly raised; and intensity and brilliance of tone are greatly in- creased. Next to the area and position of exits, depth of ribs is the most potent tone-modifier in the list. At the extreme depth of 1 and k inches, volume of tone is greatly increased, while intensity and brilliance of tone are greatly diminished, and tone-pitch is greatly lowered. In my observation, variation of 1-32 inch in depth of ribs perceptibly operates to change pmMtiUn these tone qualities. In cases of hollow, weak tone, caused by too great diminution of sounding-board thickness, I have added greatly to tone-value by simply dimin-

200

\section*{VIOLIN TONE-PECULIARITIES.}
ishing depth of ribs; and, the amount of such di- minution is always governed by the degree of tone- weakness in each case. Such diminution may equal 1-32 or 1-16, or i inch when arching equal and depth of ribs equals 1 and i inches. In this , work, the point is to diminish volume of tone, and to increase intensity of tone, brilliance of tone, and to raise tone-pitch. Such effects upon tone follow with certainty, even while all other dimensions re- main equal. It is my observation that "noise" dis- appears from violin tone as weakness of tone in- creases; and no matter whether weakness is caused by too great thickness or too great diminution of sounding-board wood, or by diminished area of the exits, or by application of a mute to the bridge. Thus the hollow, weak tone is never a noisy tone. [There are situations where volume of tone may be of greater value than intensity of tone. Such situations are found where music is drowned by noise from shuffling feet, from loud conversation, from clinking glasses, from popping corks, and from clinking sound of silver coin; situations where music is wasted on desert air I mean "smoky air" situations wherein the only hope of aestheticism is centered in the snowy crown of carbonic oxide overtopping the graceful schooner. In such situa- tion you are advised to increase the area of the ex- its.] Extraneous factors operating to diminish violin tone-power are: (1) The mute. (2) Bow-hair.

201

\section*{VIOLIN TONE-PECULIARITIES.}
Philosophy involved in the mute is the only feat- ure of interest that the mute possesses for the student of violin tone-peculiarities. The mute clearly demonstrates the amazing value of the bridge to violin tone; and such demonstration places in a clear light the fact that the bridge, in trans- mitting force to sounding-board, operates by vibra- tory action. Because the mute diminishes such vi- bratory action, therefore, the mute is of value in demonstrating the fact that rigidity of bridge- wood fiber may be so great as to diminish tone- power. The mute is also of value in demonstrat- ing the fact that the noise- wave disappears first in diminished tone-power. Thus, no matter what the degree of accent given to noise by any violin, appli- cation of a mute to the bridge operates to annihil- ate such noise. Therefore, when a case of "nerves" exist in the neighborhood of a noisy vio- lin, the mute is raised to the degree of benefactor. Vivid recollections suggest that two benefactors might be better than one. (2) Bow-hair, and the stick itself, may operate to diminish both volume and intensity of tone. I recollect an astonishing occurrence connected with worn out bow-hair. A certain player gave himself the trouble of taking a day's journey to consult me about an unaccountabe loss in the tone-power of his violin. As he presented his violin and bow to me, he remarked, ' It used to have a good, strong tone." Upon attempting to draw tone from his violin, I at first thought his bow-hair had been smeared with grease, for it slipped across the

202

\section*{VIOLIN TONE-PECULIARITIES.}
strings without producing as much tone as follows the employment of the mute. Nor did increased pressure help the matter. Examination of the hair disclosed two conditions operating to diminish tone- power. First: Diameter of the hair was of the smallest variety. Second: The barbs or scales upon which dependence is placed for exciting string-action, were worn entirely away. Laying aside his worthless bow, I applied one of mine. His violin did possess a good, strong tone. Ques- tioning brought out the fact that he had given much employment to steel E's. How any one, having sufficient intelligence to read music, could fail to know the cause for loss of tone-power in this case is a problem in human peculiarities for which I offer no explanation. Without doubt, had this party paid a good price to a good violin maker for a good violin, he would now be condemning such maker as a fraud. Un- der such provocation, the "old masters" might be excused for restlessness. To-day it seems scarcely necessary to state that bow-hair should be coarse and strong; that the best hair is obtained from the male equine; that, because the barbs thereon point in one direction, therefore, to secure equal force to up and down strokes, one-half of such barbs should point up- ward and one-half downward; that, when those barbs become worn by friction upon the strings, and worn by friction upon lumps of "rosin," then the "hank" should be turned over; that, when again worn, the whole "hank" should be thrown

203

VIOLIN TONE-PECULIARITIES. away. There is truth .in saying, ' ' Tis more difficult to find a good bow than to find a good violin." In approaching dotage, I repeat that saying. But two times in fifty years use of the violin do I re- collect of holding in hand what I call a good bow. During those years, I have held in hand many good violins. What I call a good bow is one that bal- ances at a point seven inches from the frog; that springs back into the position of rest with the ce- lerity of tempered steel; that has the lowest point of the "cambre" in its upper third, and well up toward the tip; one that seems instinctively to hug the strings. At the next hour, the subject of maximum even- ness of tone-power will be presented.

204

\section*{VIOLIN TONE-PECULIARITIES.}
LECTURE XIV. GENTLEMEN: I am now to present maximum evenness of violin tone. I confess to no little dread in approaching this problem. Perhaps dread causes me to defer its presentation until near the close of our course. Perhaps the fact that my solution for this problem has received but a single demonstration operates to increase my dread. I know full well that one swallow does not make summer. I also know full well that the most misleading writers upon the violin are such as write up tone-peculiarities of but a single violin. To present the tone-peculiarities of a single violin as facts existing in all violins ought to make the writer liable to criminal prosecution upon the charge of conspiring to defraud. Only such factors as depend upon the action of air may be relied upon as constant factors influenc- ing violin tone; whereas all such factors as depend upon the action of wood are not reliable, not con- stant, but, on the contrary, are capricious in action from first to last. 'Tis such capricious action that prevents application of hard-and-fast rules to violin construction. 'Tis plainly evident that the market would be glutted with "best" violins could hard- and-fast rules apply to violin construction. 'Tis capricious action of wood which makes the violin stand in a class by itself. 'Tis quite within possi- bility to tone-regulate all other musical devices in such uniformity as makes the tone of one precisely like the tone of another; whereas, no degree of human skill can make two violins sound precisely

205

\section*{VIOLIN TONE-PECULIARITIES.}
alike except at rare intervals. Hence, whoever writes up the details of a single good violin as being details infallible in tone-results is chargeable with fraud, whether intentional or not. Again I state that only such violin tone-mod- ifiers as depend upon the action of air can be relied upon as infallible factors. Bearing this fact in mind myself, I ask that you also bear the same fact in mind when either study- ing or applying my details for maximum evenness of violin tone. That these details may fail at times and succeed at other times is as much of a certainty as the capricious action of wood. With no desire whatever to enlarge upon my own achievement, I state that these details for evenness of violin tone- power were traced on paper as a result of a priori reasoning. This fact operates to humble my pride. The fact that I devoted a lifetime to the study of violin tone and only succeeded in securing but a single beneficial factor thereto by \& priori reason- ing is humiliating. Yet, from no authority what- ever, can I find another factor beneficial to violin tone which was traced on paper prior to the fact. All along the violin's 400-year path no fact is so prominent as the fact that every factor beneficial to violin tone is due directly to experiment alone. Hitherto, the tracing on paper of every such factor has been a result of post facto reasoning. Practi- cally, to the violin there is no difference in favor of either method for tracing beneficial factors. The only benefit lies with the violin student; and, such benefit consists wholly in encouragement. If

206

\section*{VIOLIN TONE-PECULIARITIES.}
thought can secure one factor of benefit to violin tone, then thought may secure two such factors. There's the point! To say, "The violin reached perfection 200 years ago," is profoundly discour- aging to all who accept it as the truth. Happily it is now generally understood that this saying, in late years, originated with the old-violin-trade- promoter. Application of details for maximum evenness of tone-power was accomplished upon only a single violin; since when, this three- wheel-chair, whence I address you, has occupied my attention. So marked was evenness of tone-power in this single demonstration that I feel encouraged to ask that you take my solution where I leave it; and, should my solution prove correct, then will I be amply rewarded. Without diagrams, I am confronted with an un- usual degree of difficulty in the selection of words to describe accurately the details for maximum evenness of tone-power; therefore, you may find concentrated thought necessary to an understanding of their description. In such case, I know of no other way than repeatedly going over the ground. Such repetition is work, but, for such work I offer you encouragement. As a means for assistance in making the details for evenness of tone-power stand out in a clear light, it is first necessary to present the lines of reasoning which led up to the fact. Only by slow steps and concentrated thought could I follow those lines to a conclusion; and even thus proceeding, I

207

\section*{VIOLIN TONE-PECULIARITIES.}
found necessity for bridging one yawning chasm with a mere "guess." As you may have observed in emergencies, ' 'guessing' ' quite often becomes a necessity. Every musical device possesses a fundemental, or lowest tone peculiar to itself; and, in all musical devices, except the strings, provision is made for the production of tones, higher in pitch than the fundemental tone. In all wind instruments, such provision is based wholly upon the action of air; whereas, with the strings, such provision is bas- ed partly upon the action of the air, and partly upon the action of wood and strings. This difference places the string family in a class by itself; also, this difference immensely complicates the problems involved in production of string- tones; also, this difference defies human skill to produce precise uniformity in string-tone values. Although the piano is a stringed musical device, yet; because every one of its tones is a fundament- al tone, while the violin has but four fundamental tones, therefore there exists but a remote relation- ship between these two devices. Because piano tone depends upon the direct blow of a hammer, therefore, the piano belongs in the class of percus- sion devices, notwithstanding employment of a sounding-board to augment tone-power. The tone of the violin depends upon the winding and un- winding of strings, and employment of a sound- ing-board for augmentation of power. Hereat ter- minates similarity in these two devices. The sounding-board of the piano must be lengthened,

208

\section*{VIOLIN TONE-PECULIARITIES.}
shortened, widened and regulated in thickness to accommodate each and all of its fundamental tones, whereas, that particular part of the violin sound- ing-board, augmenting the fundamental tone of each string-, must also augment all other possible tones upon each string. This dissimilarity oper- ates to permit evenness of tone-power in the piano and to cause unevenness of tone-power in the violin. Concentrated thought and experiment, directed to the piano scale, (comparative length of strings and sounding-board) has resulted in a quite satisfy- ing evenness of tone-power. In this quality of tone, the violin sounding-board yet remains faulty. Nat- urally, the question arises, "Can concentrated thought and experiment improve the violin sound- ing-board in the matter of greater evenness of tone-power?" In solving this question, it is first necessary to definitely point out errors in dimensions of the vio- lin sounding-board. Naturally, those who believe that the violin reached perfection 200 years ago will decline to admit that imperfections exist in the 200-year sounding-board; or, that imperfec- tions exist today in sounding-boards precisely sim- ilar to the 200-year sounding-board. Possibly, no evidence whatever can change such belief. As a successful method for maintaining error, there's none quite so successful as refusal of* evi- dence. In calling upon the piano sounding-board for evidence of error in the violin sounding-board,

209

\section*{VIOLIN TONE-PECULIARITIES.}
there's bound to be some who will decline to re- ceive such evidence for fear of being convinced against their will. Here follows such evidence: (1) The piano sounding-board is heaviest be- neath the larger strings. (2) The piano sounding-board is lightest be- neath the smaller strings. (3) The piano sounding-board is longest be- neath the larger strings. (4) The piano sounding-board is shortest be- neath the smaller strings. In the Strad sounding-board there's no difference in length of sounding-board activity beneath the strings. In the Strad sounding-board there's no difference in thickness beneath the strings. From these facts, it is clearly evident that the violin sounding-board of 200 years ago had not reached perfection. [In the matter of producing evenness tone-power, some of the Joseph Guarnerius sound- ing-boards reached much nearer perfection than any sounding-board of the Stradivari.] Because the piano sounding-board of today pro- duces the greater evenness of tone-power, it is evident that error lies in the violin sounding-board; and, such error becomes apparent in attempting to produce two octaves of tone, having even pow- er, upon any of the four violin strings. Even with the aid of increased bow-pressure, such attempt is a failure. The reason for such failure is apparent; the same length and thickness of sounding-board, engaged in augmenting the fundamental tone,

210

\section*{VIOLIN TONE-PECULIARITIES.}
must also be employed to augment the tone from the shortened, and consequently weakened string. For this fact, there is but one conclusion; that is both length of sounding-board activity, and rigidity of sounding-board are too great for the weakened blow from shortened string. Starting with the fundamental tone of any violin string, and count- ing half-intervals, there are twenty-five tones within two octaves to be augmented by an identi- cal set of sounding-board fibers. Upon the piano, each of these successive tones is augmented by a sounding-board of diminished length, and diminish- ed thickness. Upon the violin, such diminished length \_ and diminished rigidity of the sounding- board, for each successive tone, is manifestly im- possible. Indeed, equal evenness of violin tone- power with the evenness of piano tone-power would be a forlorn hope were it not for the aid from increased bow-pressure. It is apparent that production of even tone pow- er throughout two octaves upon each violin string is a problem beset with difficulties. It is also ap- parent that a solution for such problem is of value to violin tone; and in all conscience, to claim im- provement upon the best sounding-board methods of Stradivarius and Joseph Guarnerius is apparent- ly sufficientl cause for hesitation. During 200 years, those methods have been faithfully copied by the most ambitious violin builders throughout the civilized world; and, in all the world, I've known of none, nor heard of none claiming im- provement upon the best from those two famous

211

\section*{VIOLIN TONE-PECULIARITIES.}
experimentalists. Yet, notwithstanding the fact that none are found making such claim, I know of modern violins equally as beautiful in tone as any Strad I've ever heard, and, while equally sweet in tone, possess greater evenness of tone-power. Why not claim the due? I cannot see any harm in claiming such due when the claimant has the goods to show! Tis true, 'tis best to never claim more than is due; and, 'tis best to be modest in making any claim for onesself but, 'tis of no avail to hide one's light under a bushel. ; The world is ever ready to accord merit at the moment of conviction that mer- it is due; but, the world rightfully demands proof before judgment. By all considerations, submit the proof. In submitting such proof, there's none to fear except the old- violin-trade-promoter; and, it is now quite generally understood that he is booked for that port where willful liars are consigned. As we elderly students of the violin vividly re- member, Ole Bull could draw four octaves of quite even musical tone from each string of his '"Jos- eph." As all students of the violin know, 'tis a fallacy to claim such merit to lie wholly with the performer. No matter who the performer, he first must hold the goods in hand. It *is my own observation that the "Joseph," among old violins, remains unequaled as a solo violin in the larger auditoria; and, such opinion is based upon the fact that the "Joseph" possesses both greater pow- er and evenness of tone-power. 'Tis but natural for

212

\section*{VIOLIN TONE-PECULIARITIES.}
'What the student to ask, ' causes greater power and greater evenness of tone-power in the "Joseph?" During many years, I have tried to answer this question, but, for same reason left to conjecture, precise data for the ''Joseph" are kept as a valu- able personal asset. Human nature yet remains very human. Concerning data for the "Joseph," Honeyman states thus: "The "belly" is thickest at the edges, and thinnest throughout the central areas." From such indefinite description, but lit- tle can be learned. Even a writer, so thoroughly equipped by education, by thorough training in violin construction, and, by opportunity for obser- vation as Ed. Herron Allen, refers his readers to the Strad data for thickness of the "Joseph" ta- bles, although claiming to have had M. Sainton's Joseph from which to obtain data. Strange! Here's one yet more strange. I sent to Paris for M. Simoutre's book and charts because among those charts was one of the Joseph. I received that book with the Joseph chart cut out. "Diable!" From a friend, I obtained a copy of a chart of the "Joseph." The original of this chart was made in the studio of Vuillaume. From the chart I received, and from Honeyman' s indefinite de- scription of the Joseph, I gave much time to re- peated experiment. Today I'm not mourning the loss of Simoutre's chart; and, because of the opin- ion that is impossible for any one to give more time to experiment upon the possible varieties of violin

213

\section*{VIOLIN TONE-PECULIARITIES.}
plate. thickness than I have given. Charts for the Stradivari, of several different varieties in plate thickness, are quite easily obtained. I have not only given repeated trial to all Strad varieties in plate thickness, but have also given many trials to the obtainable data for the Joseph. Between these two great genii, I unhesitatingly place the credit for greater tone-power, and greater evenness of tone-power to Joseph Guarnerius. For the greater duration of tone, and for the greater sweetness of tone, I place the credit to Stradivarius. The method of plate-thickness finally adopted by these two successful experimentalists is presented for the purpose of making my lines of reasoning, leading to a partly new method, stand out in a clear light. Thus: In experiments upon the Joseph method, it was clearly demonstrated that greater tone- power followed placing the thinnest point in the sounding-board half-way from position of the bridge to ends of the plate. Observe the two following facts in the methods of Strad and Joseph: (1) Thinnest point beneath all the strings is equally distant from position of the bridge, regard- less of difference in weight and diameter of strings. (2) Equal thickness of the plate beneath all strings, regardless of the difference in weight and diameter of strings. At this moment I call only brief attention to the differences between this method for sounding- board thickness and the method, (my method) for

214

\section*{VIOLIN TONE-PECULIARITIES.}
yet greater evenness of tone-power. The favorite method of Stradivari is: (1) Thinnest point in the sounding-board is at the greatest practical distance from the bridge for all strings without regard to differences in weight and diameters of strings. (2) Sounding-board thickness is equal beneath all strings without regard to difference in weight and diameter of strings. Comparing the new method, thus: (1) Thinnest point beneath the G-string is at greatest practical distance from position of bridge. (2) Thinnest point beneath the D-string is less distant than for the G. (3) Thinnest point beneath the A-string is less distant than for the D. (4) Thinnest point beneath the E-string is less distant than for the A. (5) At position of the bridge, greatest sound- ing-board thickness is beneath the G. (6) At position of bridge, less sounding-board thickness for D than for G. (7) At position of bridge, less sounding-board thickness for A than for D. (8) At position of bridge less sounding-board thickness for E than for A. This new method is based upon both fact and theory. The fact is: Strings vary in weight and diame- ter; hence, force in string action varies. Theory is: Sounding-board thickness beneath each should vary as the diameter and weight of

215

\section*{VIOLIN TONE-PECULIARITIES.}
strings; and, length of sounding-board activity be- neath each string should vary as the pitch of tones demanded from each string. For this fact, no proof is needed; it is self-evident; but the theory needs proof. As such proof, I first offer in evidence the fact, and a demonstration for the fact, that each string largely (not wholly) de- pends upon sounding-board activity directly be- neath for augmentation of tone. For such demon- stration, I offer this violin as a sacrifice. True, 'tis offering up but the one for the benefit of the many, yet, somehow I feel very much like a ' 'wood butch- er," that is, granting any feelings to the "wood- butcher." With this thin blade, I proceed to split the sounding-board of this violin in various places, and, shall continue such work until the tone of all strings is completely ruined. (1) Beginning at the lower extremity of the right exit, I split the sounding-board down to the purging. Application of the bow shows no damage whatever to tone of any string as following this split. (2) Beginning at a point near to, and below the post, I split the sounding-board down to the lower end-block. Application of the bow again shows no damage' to tone upon any string. (3) Beginning near to, and below the bridge, I split the sounding-board along the center-join down to the lower end-block. Again, no damage to tone of any string. (4) Beginning at the lower extremity of the left! exit, I split the sounding-board down to the

216

\section*{VIOLIN TONE-PECULIARITIES.}
I purng. Again, no damage to tone. (5) Beginning at point even with the bridge, and to the left of the bar, I split the sounding- board down to the lower end-block. Although not what might be called complete ruin, yet, both G and D-tone now show serious damage. (6) Beginning below the bridge, and to right of bar, I split the sounding-board down to lower end-block. Damage to G and D-tone is increased, but yet no injury to tone appears on A and E. (7) Beginning at a point even with the bridge; and one inch to the left, I split the sounding-board upward to the purnhg. From this split, there is only slight additional damage to G-tone, but, there follows no additional damage to D-tone. (8) Beginning at the bridge, near to, and to the left of the bar, I split the sounding-board upward to the end-block. G-tone is now completely ruin- ed, while D-tone, although injured, yet, is not to- tally ruined. (9) Beginning at the bridge, and beneath the D, I split the sounding-board upward to the end- block. D-tone is now completely ruined, while A-tone, although injured, is not totally ruined. (10) As before, splitting beneath A completely ruins tone of that string, while E-tone is but par- tially ruined. (11) As before, splitting beneath the E com- pletely ruins tone of that string. Thus is clearly demonstrated the areas of sound- ing-board activity upon which each string depends for augmentation of tone-power; and their definite

217

\section*{VIOLIN TONE-PECULIARITIES.}
location proved to be of immense value to the pro- duction of even tone-power. It is evident that such location of areas pemits of precision in reduc- tion of sounding-board thickness beneath each string proportionate to the diameter and weight of each string. Thus far, proof is conclusive; but, the next step is a step in the dark. The question is, "What shall be the ratio of lengths for sounding-board activity best augmenting the tones of each string?" The ratio for shortening the piano sounding- board cannot apply to the violin; nor in the entire range of musical devices can I find a ratio which may apply to the violin. In stumbling through a text-book devoted to the philosophy of musical sound, I found a ratio for fifths of the major scale. This ratio at once at- tracted my attention, and, because it not only points out the difference in the number of vibra- tions per second for two or three consecutive fifths, but, for all possible fifths in the major scales. It is a constant ratio, That's the kind of ratio I now am wanting. The book states that multiply- ing the number of vibrations for any tone by 3-2, finds the number of vibrations in its fifth above. Wishing for certainty in this matter, I proceed to test the constancy of this innocent-looking ratio, 3-2. Helmholtz states you know Helmholtz he must be German I think so anyway because he tells us more about the philosophy in musical sound than all other philosophers from all other countries jumbled together Helmholtz states that open G,

218

\section*{VIOLIN TONE-PECULIARITIES.}
violin at concert pitch, vibrates 200 times per sec- ond; therefore, if 200 be multiplied by this easy ratio, 3-2, 'twill find the fifth above open G; 200, times 3-2 equals 300; 300, times 3-2 equals 450; 450, times 3-2 equals 675. These figures do represent the open tones of the violin, G, D, A, E. How easy! This philosopher business isn't much after all. Anybody can run it. Yes, if 3-2 works for music, it certainly ought to work all right for the violin; yet, I do recollect cases wherein there appeared no visible ratio of any size between mus- ic and the violin. But, as I'm wanting a constant ratio for lengths of sounding-board activity be- neath violin strings, and, as no other ratio than 3-2 appears taiit the violin, therefore I'm bound to try 3-2. Thus : Considering the greatest length of sound- ing-board activity beneath G to be 12 inches, there- fore all that's needed to find such length for D is to use this innocent-looking ratio, 3*2,.-- upon , 12; and, 12 multiplied by 3-2 -equals 18 what! 18 inches of sounding-board activity for D! Wish I" could see Helmholtz for a minute. - In life there is a thing called misplaced confi- dence. Its location may be with others at times; at other times it may be with ourselves. The latter is the worst kind we can't cuss Jones. ' Yet, notwithstanding the fact of being led away by 3-2, I gave many a'long month to search for a ratio which would apply to the length of sounding- board activity beneath the strings. But, in both waking and sleeping hours, wherever and when-

219

\section*{VIOLIN TONE-PECULIARITIES.}
ever I looked for such ratio, there stood that 3-2 like the proverbial ghost. It would neither down, nor get up, nor go away; but remained as unruffled as a wooden Indian. I grew to hate 3-2; but, the more I hated it the more I had of it. In despera- tion, I searched for precedents giving relief from phantoms. I recalled successful treatment of oth- ers who were suffering from "phantoms," but, mine was not that kind. Late one sleepless night my long-looked-for precedent came. It was the precedent afforded by the captain of a Mississippi scow. This captain hearing about a new-fangled invention which was capable of fortelling approach of storms, invested good money in a barometer, and installed the same near the steering gear. Soon thereafter a great windstorm threatened to send this particular scow to the bottom. The bar- ometer, not being in the wind-storm business, calm- ly rested in an indifferent manner. To the captain, the difference between a wind-storm and a thund- er and lightning rainstorm cut no ice. Feeling himself the victim of misplaced confidence, the captain seized upon that impassive barometer and turned it up-side down. Here was my preced- ent. Because that captain thereafter made a successful "tie-up," therefore I jumped up, seized upon this impassive 3-2 and turned it up-side down. Then I slept. I confess that this new ratio, 2-3, is a mere whim; but, the violin itself is a product of whims. Dur- ing its first 200 years of development, the violin owes everything to whims; then; because two

220

VIOLIN TONE-PECULIARITIES. lucky whim-ers became genii, luck seems to have abandoned the class; that is, did we but trust the old violin promoter. I think by the 200-year prec- edent, 'tis now time for more luck.

221

\section*{VIOLIN TONE-PECULIARITIES.}
LECTURE XV. GENTLEMEN: Maximum evenness of violin tone -is resumed. It is a wise provision of law that proof must precede conviction. It is sometimes wise to suspend judgement upon mere assertions until after presentation of proof, except those cases wherein proof is self-evident. Thus the statement that errors in graduation of the violin sounding- board operate to cause uneven tone-power is an assertion needing proof, because proof is not self- evident. The experienced violin tone-regulator understands such proof; but, as there are many interested in the violin without having any experi- ence in tone-regulation work, therefore such proof is here presented. In connection with the statement and proof that errors in sounding-board operate to cause uneven tone-power, it is interesting to note the wide vari- ations in such thicknesses employed by the cele- brated builders of 200 years ago; also, to note the wide variations by different writers concerning such thicknesses. From the evidence thus obtain- ed it is clear that extremes in plate thickness were represented by Stainer, for greatest thickness, and by Stradivarius, and Joseph Guarnerius for least plate thickness, thus: Stainer, at the position of bridge, 5 mm, down near the edges to 1 mm. Stradivarius, at the po- sition of bridge, 2 and 8-10 mm, down near the edges to 2 and 7-10 mm. These figures are given by Simoutre, Paris, 1885, violin maker and collec- tor. But, it is clearly in evidence that we should

222

\section*{VIOLIN TONE-PECULIARITIES.}
not consider these figures for plate thicknesses by Stradivarius as being the only figures employed by this tireless experimenter. Simoutre's figures are for a Strad of 1707; whereas Honeyman gives fig- ures for plate thicknesses of a Strad, 1708, as 1-8 throughout the entire plate. It is thoughtful of anyone, when referring to a particular Strad violin, to give its date, because the date assists in making clear to the reader both model and plate thickness- es. To the experienced tone-regulator, it is at once apparent that the difference in plate thickness- es of the Strad, as given by Simoutre and Honey- man, operates to cause a difference in tone-pitch, and, in duration of tone. Thus, from Simoutre's figures, tone-pitch will be the lower, duration of tone greater, also, permitting use of lighter strings. Although this difference in plate thickness is but slight, yet it is ample to cause perceptible change in tone-pitch, duration of tone, and volume of tone, because plate thickness is down near to the limit of safety. By "limit of safety" is meant that degree of thickness below which weakness of tone follows. There is another 200-year fact which seems worthy of presentation with the lighter 200-year sounding-board. Thus: Two hundred years ago, concert pitch varied according to locality, from A, 405 vibrations per second at Paris, to 451i at Milan. It is but natural to connect the fact of low concert pitch at Paris with the fact that at Paris the Strad violins at once acquired a degree of favor never acquired in Italy. I mean violins of Strad 's third period. This statement is based up-

223

\section*{VIOLIN TONE-PECULIARITIES.}
on evidence of Italian musicians given to me per- sonally. "Tis but little more than a decade since "normal" pitch, or diapason normal, or in- ternational pitch was established with A at 435. It is a noteworthy fact that reducing concert pitch to A, 435, greatly accommodates those worn, weaken- ed, originally light 200-year sounding-boards yet in use. 'Tis but natural for the student to ask, "When those old sounding-boards are gone, will concert pitch be raised?" This question possesses interest for every violin tone-regulator. When a violin is carefully adjusted in sounding- board rigidity, bridge rigidity, quality, mass, and position of post for guage-2 strings tuned to A, 435, thereafter tuning A to 450, is but inviting disaster to tone- values. Again, when a violin sounds equally well at either pitch, such violin is not tone-regulat- ed with care. From any point of view, 'tis but wisdom to settle down upon a universal pitch for all concert instruments, because then the violin builder might do tone-regulation with precision, thereby avoiding unmerited adverse criticism for putting out violins which will not stand at various pitches at the caprice of various conductors. Neither in any 200-year sounding-board, nor in any later sounding-board have I found a method for graduation based upon the self-evident fact that force in violin strings varies as their diameter and weight. Considering evenness of tone-power to be a valuable feature of violin tone, and know- ing that no 200-year violin builder made provision in sounding-board rigidity for diminished force in

224

\section*{VIOLIN TONE-PECULIARITIES.}
the lighter strings, I therefore discredit the oft- repeated statement that the violin reached perfec- tion 200 years ago. In proof of correctness in my position, I now present a method for sounding-board graduation based upon two facts, thus: 1. Force in violin strings varies as string diam- eters and weight. 2. Augmentation of high-pitched tones demand both shorter and lighter sounding-board fibers be- neath them. The violin world has stood still, and with its gaze fixed upon Cremona, while the piano sounding- board has developed such evenness of tone-power and such quality of tone as threatens to rob "the king" of its title. In this connection, the piano sounding-board possesses interest for the violin student. Its shortened fibers and diminished thick- ness from beneath bass strings to beneath the lightest and shortest strings is a broad hint to such violin devotees as can see nothing but "Cremonen- sis faciebat." At a previous hour, a practical demonstration was submitted as showing that each violin string largely depends upon sounding-board fiber directly beneath for augmentation of tone. These three violins, designated by M. N. R. incidentally afford corroboration for the above fact; but at this mo- ment, these violins are submitted as proof that un- evenness in violin tone-power is caused by errone- ous graduation of the sounding-board. Each of these sounding boards is graduated after a diff er-

225

\section*{VIOLIN TONE-PECULIARITIES.}
ent method. For a purpose, the sounding-boards on M and N are given slightly exaggerated reduc- tion in thickness. The sounding-board of R is giv- en a thickness of i throughout. Violin M, graduated thus: At bridge, beneath G and D, 9-64; thence down at ends to 4-64. At bridge, beneath A and E, 4-64; thence down at ends to 4-64. Violin N, graduated thus: At bridge, beneath G and D, 4-64; thence down at ends to 4-64. At bridge, beneath A and E, 9-64; thence down at ends to 4-64. Violin R, graduated thus: Thickness of sounding-board throughout i. I will first apply the bow upon A and E strings of violin M. Beneath these strings, sounding- board rigidity is greatly reduced from position of bridge to ends of plate. It is apparent that a great- er length of fiber-activity cannot be given to the 14 inch sounding-board. You observe that the fun- damental tones of these strings are characterized by unusual volume, by low tone-pitch, by feeble intensity, and by freedom from noise. These ef- fects upon tone are expected by the experienced tone-regulator; and, the explanation is found in the fact that diminishing thickness of a tone-pro- ducing agent, other dimensions remaining equal, lowers tone-pitch; that lengthening perpendicular, confined air columns lowers tone-pitch; that, as volume of tone increases, intensity of tone dimin-

226

\section*{VIOLIN TONE-PECULIARITIES.}
ishes; that as power of tone diminishes, noise dis- appears. But, the next tone-phenomenon on these strings goes begging for an explanation. I now draw out two octaves upon each of these strings. As you observe, each successive, higher tone is characterized by increasing loss of power. "Mr. Tone-regulator, please explain?" "No?" I know not where others may turn for this ex- planation; but, I turn to the modern piano sound- ing-board, and to the modern piano designer. I ask him, "Why do you constantly shorten fiber- activity beneath succeeding tones of higher pitch?" I am not surprised at his elevated eyebrows as he asks me, "Are you from Cremona?" At this moment, G and D- string tones, violin M, have no further interest than baritone character and fair duration. Violin N, sounding- board graduated the reverse of M, greatly changes tone-character of its G and D. The tone from these strings is also character- ized by low pitch, by great volume, by feeble in- tensity, and by freedom from noise. Of these three violins, R possesses much the greater tone-value. Only its altissimo tone-pow- er is here in point. As I draw two octaves from this E, you observe perceptible diminution in pow- er from successively higher tones. From the tone of these three violins, I reach the following con- clusions: 1. In violin M, length of fiber-activity beneath A and E-strings is too great, and rigidity too great-

227

\section*{VIOLIN TONE-PECULIARITIES.}
ly reduced. 2. Ditto, violin N, beneath its G and D. 3. Violin R, rigidity too great beneath its A and E. From my point of view, these violins afford conclusive evidence that erroneous graduation of the sounding-board causes unevenness in tone- power. These violins also afford conclusive evidence that each string depends upon fiber-activity directly be- neath for augmentation of tone; also, that length of fiber-activity should vary as the pitch of tones; also, that rigidity of the sounding-board fibers should vary as the force in strings. Standing out in a clear light, these reasons led me to work out a method for sound-board graduation intended as a practical demonstration for its value in securing maximum tone-evenness. As previously stated, this method was applied to but a single violin be- cause my time for work came to an abrupt termin- ation; but, you have my assurance that the results are encouraging in a high degree. I give assur- ance that, as a solo instrument, there is a vast dif- ference in tone-values favoring that sounding- board having provision in rigidity for varying force in strings, and having provision in length of fiber- activity for tones of higher pitch. By no means do I claim that the following details reach perfection; and I not only grant permission, but also request younger students to improve upon them. I do not claim 2-3 to be the best ratio for diminishing lengths of fiber activity beneath violin strings. I only claim 2-3 to be a ratio stumbled

228

\section*{VIOLIN TONE-PECULIARITIES.}
upon in an effort to banish that 3-2 ratio for fifths in our present-day major scale. Having presented the principles leading up to details for maximum evenness of tone-power, I now present such prin- ciples and details in condensed form. These principles are: 1. Greatest thickness beneath G. 2. Thickness diminished from G to D. 3. Thickness diminished from D to A. 4. Thickness diminished from A to E. 5. Length of fiber activity greatest beneath G. 6. Length of fiber-activity diminished from G toD. 7. Length of fiber-activity diminished from D to A. " 8. Length T of fiber-activity diminished from A to E. \textasciicircum\{\}1 Thus, one of these principles refers or is applied to sounding-board thickness beneath each string; the other principle applies to length of fiber-activi- ty beneath each string. In practical application of the latter principle, it is evident that fiber-activity beneath G, must be determined first. Here is an- other chasm across this path which I bridged with surmise or "guess." As you may have observed, there are emergencies besetting life's path for which the only recourse is ' 'guessing. " Of a truth, such emergencies may be found in the violin with- out use of the field-glass. In the matter of deter- mining the precise length of fiber-activity in each violin sounding-board, I know of nothing reliable. From the evidence afforded by varnish

229

\section*{VIOLIN TONE-PECULIARITIES.}
phenomenon No. 1, I "guess" that one inch at each end of the sounding-board does not act with suffi- cient energy to produce audible sound. Basing conclusion upon this guess, the greatest possible length of fiber-activity is 12 inches. Taking this length of fiber-activity beneath G as a starting point, and, applying thereto the ratio 2-3, finds length of fiber-activity beneath the D, to equal 8 inches. Again, applying 2-3 to 8, finds length of fiber-activity beneath A, to equal 5, and 33-100 inches. Again, 2-3 of 5 and 33-100 equals 8 and 77 100 inches as the length of fiber-activity beneath E. Condensed thus: Fiber-activity beneath G, 12 in. Fiber-activity beneath D, 8 in. Fiber-activity beneath A, 5 and 33-100 in. rr Fiber-activity beneath E,-8and 77-100 m. -3 loo In using the term "fiber-activity" there is diffi- culty in making my meaning clear. It is clear that 12 inches of fiber-activity beneath G, is practically the limit; and, to secure this length, sounding- board thickness must gradually diminish from bridge-position to ends of the plate, and diminish down to a thickness which is the limit of safety. In the well known Stainer method, thickness at this point is reduced to 1 mm, or, 1-25 inch, and this point he places equally distant from the bridge beneath all strings. In my details herein present- ed, thickness at this point is only reduced to 4-64, or, ,1-16 inch; and, this point varies in distance from the bridge by the ratio 2-3. Therefore, by the

230

\section*{VIOLIN TONE-PECULIARITIES.}
term "fiber-activity," I mean the distance above and below the bridge between two points of great- est reduction in sounding-board thickness. Thus, in the case of fiber-activity beneath D, one-half of the length, or 4 inches, is above, and one-half be- low the bridge position; ditto A, and E. But, be- cause of placing the bridge 8 inches from upper end of plate, therefore bridge-position is not at the half-way point in the length of fiber-activity be- neath G, 7 inches of such length being above, and 5 inches being below bridge-position. Condensed thus: Above bridge, fiber-activity beneath G, 7 in. Below bridge, fiber-activity beneath G, 5 in. Above bridge, fiber-activity beneath D, 4 in. Below bridge, fiber-activity beneath D, 4 in. Above bridge, fiber-activity beneath A, 2 and 66-100 in. Below bridge, fiber-activity beneath A, 2 and 66-100 in. Above bridge, fiber-activity beneath E, 1 and 38-100 in. Below bridge, fiber-activity beneath E, 1 and 38-100 in. (Obviously, those fibers beneath G and D, must receive identical treatment above and below the bridge; but, those fibers beneath A and E, do not necessarily need identical treatment above and be- low the bridge, because the post prevents action of those fibers from passing its position as previously demonstrated by splitting the lower right-quarter of the sounding-board. However, 'tis but natural

231

\section*{VIOLIN TONE-PECULIARITIES.}
to give identical care to thickness throughout the

plate.) Thus, the details for lengths of fiber-activity be- neath the strings are presented. The difficulty in making these details stands out in a clear light, without aid of diagrams, is apparent. In this mat- ter, failure will cause me no surprise; but, you may understand my willingness to be "interviewed." Tis the best I can offer. Next come details for thickness. First, thick- nesses are established at bridge-position. Before giving figures for thickness, it is necessary to ex- plain that the grain of this sounding-board possess- es slightly too great density to yield the "rich" tone; that they are faultless in every other feat- ure; that this sounding-board has been in the ser- vice of two generations; that its shrinkage is com- pleted; and, that its spring-action is superlative. Without these facts, the following figures for thick- ness might appear to pass beyond the limit of safe- ty. By no means do I advise adopting these fig- ures as thickness for all samples of sounding-board wood. With wood of softer grain, and with less time from the builder's hands, I would hesitate be- fore reducing thickness to 4-64 at the ends of fib- er-activity beneath the A and E, as per details. In this case, desiring the baritone character for the G, thickness, at bridge-position, is reduced to 9-64; next, thickness, at bridge-position, beneath E, is reduced to 7-64; and, thickness is gradually diminished from G to E. Thus, immediately at bridge-position, some degree of provision is made

232

\section*{VIOLIN TONE-PECULIARITIES.}
for diminished force in the strings. Although the following table gives precise figures for thickness beneath each string at bridge-position, yet, change in thickness beneath the strings is not abrupt, but, is shaded down gradually. Condensed thus: Plate thickness beneath G, 9-64. Plate thickness beneath D, 9-64. Plate thickness beneath A, 8-64. Plate thickness beneath E, 7-64. (In sounding-board wood of soft grain, I have observed the thickness of 1-8, at bridge-posi- tion, to pass beyond the limit of safety, as shown by weakened tone-power in the D-string, often in both D and A. This phenomenon is worthy of at- tention. As a remedy for such weakened tone- power, I have met with considerable success from gluing a block of pine across the inner surface of the plate at bridge-position. The dimensions of the block are: Thickness, 1-8; width, 3-8; length, 3-4. Whenever weak tone in D and A; is caused by light sounding-board, this block operates to raise tone-pitch, and increase tone-power of the D and A, but, in nowise affecting tone of G and E. It is my conclusion that lightness of wood at bridge-position operates to lengthen fiber-activity beneath D and A, thus permitting action to pass between bridge-pedestals. Evidently, the block arrests action at bridge-position.) In this experiment, thickness is equal at the ends of fiber-activity throughout; that is, 4-64. Refer- ring to the table for lengths of fiber-activity, it is

233

\section*{VIOLIN TONE-PECULIARITIES.}
observed that, beneath E, the distance each way from bridge-position to the point where thickness reduced to 4-64, is but 1 and 38-100 inches; beneath A, but 2 and 66-100 inches. I confess that such reduction at these points re- quires "nerve;" at least in my own case. But, I am careful that the thickness of 4-64 does not ex- tend along the fibers to exceed 1-2 inch. Because these lengths of fiber-activity are determined by a constant ratio, therefore a line, drawn across their extremities, runs obliquely across the plate; and, along this line, thickness is reduced to 4-64; and, from bridge-position, thickness gradually diminish- es down to such extremities. From 4-64 to ends of the plate, thickness gradually increases to 9-64. Thus is the description of details which vary from customary details in sounding-board graduation. Treatment of the bar, being the same as hereto- fore described, is not here in point. Doubtless, further experiment will discover a better ratio for lengths of fiber- activity. Had I been permitted more time for work, trial would have been given to 3-4 as such ratio. However, 'tis but wisdom to trust that bridge which permits us to pass safely over it. The ratio 2-3 did serve me well, for equal evenness in violin tone-power mine ear hath not heard. Good luck to you, little ratio. May you meet your 200-year relatives in Cremona heaven.

234

\section*{VIOLIN TONE-PECULIARITIES.}
LECTURE XVI. GENTLEMEN: At this hour the subject of maxi- mum violin tone-power is presented. Today the greatest power in violin tone attracts great- er attention from violin users, violin build- ers, and violin students than at any other period in violin history. The cause for this fact is found in the increased seating capacity of the mod- ern auditoria. Man is a gregarious animal, and a lover of music; and, for some occult reason, his en- joyment of music is as the square of his numbers. Without doubt, this fact is due to the geometric progression of sympathetic mental action. Thus, when but a single listener is present at rehearsal, sympathetic action remains at zero; a fact demon- strated by conspicuous absence of demonstration. Upon addition of another listener, sympathetic ac- tion develops as the square of two; whereat, the single listener becomes conscious of added enjoy- ment. To the musician, such increase in enjoy- ment by the listener is a greater stimulus than the stimulus of fine gold. 'Tis ever an object of human ambition to accom- plish great feats. 'Tis the violin soloist's ambition to please his patrons in the larger auditoria equally with the patrons in the smaller auditoria. Such ambition is laudable, but, the difficulties are rapid- ly approaching the unsurmountable. Hence the "big" price for violins of big tone. But, the pres- ent demand for violins of "big" tone is an unnec- essary weakness in violin users. From the thous- ands devoting their utmost energy to the violin,

235

\section*{VIOLIN TONE-PECULIARITIES.}
but a limited number ever will be called upon for an appearance in the larger auditoria. Only genius combined with almost superhuman drudgery can command widest patronage for the violin soloist of today. Thus, there is but a limited employment for violins of "big" tone; but, 'tis our weakness to bestow an admiring glance upon that party who says, "Your violin has a big tone." "Pis said ' 'we should take humanity as we find it;" but because of implied permission not to leave humanity as we find it, therefore I say that it is a compliment of vastly Tour greater value when a compe- tent party says, ' violin has a beautiful tone. ' ' Such statement is based upon the fact that music is nothing if not beautiful. I do not mean that ' 'big" violin tone may not be beautiful; but, do mean that distance of the listener should be pro- portionate to "bigness" of tone. In the smaller auditoria, in studio, or in apartment, "big" tone causes distress to the cultivated ear, no matter what the skill in bowing. These conclusions are bas- ed upon the general understanding of the meaning in the word "big" as applied to violin tone. Thus employed, "big" tone means marked volume of tone combined with marked intensity of tone. This combination is rare; and, to the best of my observation, is rare because of the rarity of sound- ing-board wood possessing superlative spring-ac- tion. This statement is based upon repeated fail- ures to produce "big" tone at will. Thus, with every other factor at its best for augmenting vol- ume and intensity of tone, yet, lack of superlative

236

\section*{VIOLIN TONE-PECULIARITIES.}
spring-action in sounding-board wood may defeat "big" tone. In situations where space is proportionate, "big" violin tone may become agreeable. It is a fact that noise is disagreeable as its proximity. It is also a fact that the noise-wave, from all violins, travels only to comparatively short distances, ex- cept in cases of woody tone-quality. Because vio- lins of marked woody tone-quality have no 'value, therefore such violins are not included in my mean- ing when saying that the violin soloist need not fear to appear in the larger auditoria with a modern violin possessing "big" tone, even when such v tone is somewhat noisy in close proximity. There are some modern violins combining "rich" tone with marked volume and intensity of tone. When heard at a distance, it is my observation that but few ears are sufficiently acute to distinguish any difference between the tone of such modern violins and the tone of old violins yet possessing power. Violin tone has ever and, and undoubtedly will ever remain separated into two classes by reason of two irreconcilable factors, as: The business factor. The aesthetic factor. For its existence, the business factor has no oth- er object than profit; and its income is reckoned in hundredths of a dollar. The sethetic factor exists alone for human pleas- ure in beautiful sound; its income is a donation, and never reckoned at all. The difference in these two factors is as wide as the East is from the West.

237

\section*{VIOLIN TONE-PECULIARITIES.}
The business factor dates from the mists prior to the time of Da Salo. The aesthetic factor dates from Cremona. The business factor always has been, and yet is numerically the stronger. The business factor demands the greatest possible vol- ume, and the greatest possible intensity of violin- tone. The aesthetic factor demands greatest pos- sible beauty in violin-tone; yet, to satisfy the aes- thetic factor is a difficult matter. Violin aesthetic- ism reckons not the cost. To ordinary mortals, violin aestheticism seems sometimes to run into madness. Potency in violin aestheticism is second only to potency in religion. The potency in violin aestheticism created the Thomas Hall. 'Tis not at all necessary to state that said Hall is located at Chicago. The world learned of this fact by electric current. Nor in New York, nor in Boston, nor in London, nor in Paris, nor in Berlin, nor in Vienna, nor in Milan, nor in the wide world is there a parallel object lesson for the potency in old violin aesthe- ticism. Only the potency in religion can pick up a simil- ar million partly from beneath the feet of scrub women; and, such women exist only in Chicago. Quartered many years in the Auditorium, the Thomas Orchestra found nothing for complaint ex- cept the fact that a large part of its patrons could not hear first- violin tone. Such complaint was well founded. At distant seats, conspicuous ab- sence of first-violin tone was ample for disappoint- ment; and, my heart cried out against the heart-

238

\section*{VIOLIN TONE-PECULIARITIES.}
lessness displayed in forcing these old violins into a situation exciting contempt for their weakened tone-power. There are three remedies for avoiding such dis- appointment, thus: 1. Replacing the deadening carpet of wood-fiber, wood-dust, or dirt within these old violins with a permanent, and perfectly smooth surface. 2. By substitution of modern violins having ad- equate tone-power. 3. By diminishing seating capacity of the audi- torium. With but an increase of 35 per cent in intensity of first-violin tone, the Thomas orchestra might have remained at home in the Auditorium for an indefinite period, thus avoiding acceptances from scrub wt>men. Naturally, substitution of modern violins having . adequate tone-power suggests itself as the first remedy for the disappointment in not hearing first- violin tone. Herein lies the. whole difficulty in this case. This suggestion points to the employment of the modern violin for interpretation of Haydn, Mozart and Beethoven scores, a possibility not ad- mitted by Dr. Thomas. Over his signature, Dr. Thomas states that ' 'the best of Cremona violins, together with the Tourte bow, inspired the master works of Haydn and Mozart." To this statement there is no dissent; 'but, to his statement that the Cremona violin is yet a necessary vehicle for the intepretatibn of master scores, there is dissent; and, such dissent comes from every disappointed

239

\section*{VIOLIN TONE-PECULIARITIES.}
patron. Were those older composers present at rehearsal in the Auditorium doubtless they also would join in such dissent. Perhaps there is nothing concerning the "best'* of the Cremona violins, (the Stradivari, to which Dr. Thomas particularly alludes, ) more firmly es- tablished than the fact that the "best" are worn out by those disintegrating forces, heat and moist- ure; and, that "next best" are so nearly worn out as to possess but little value as vehicles for inter- pretation of either great or small compositions. It is safe to assume that none of these old violins pos- sesses the same tone-vigor as on the day when in- spiring the Haydns, Mozarts, and Beethovens of an hundred years ago. The third remedy for removal of disappointment was chosen. This example of devotion to a de- lusion cost a round million; but, that million cost nothing more than stooping to pick it up. By all the logic in history and philosophy, another decade- and-a-half will bring yet another necessity for di- minishing seating capacity. In view of the abundance of modern violins pos- sessing "rich" tone combined with ample tone- power, the statement of Dr. Thomas concerning 200-year violins as "necessary vehicles" seems to me like violin aestheticism run mad. This digression is made for a purpose. First: To make clear the error in claiming the 200-year violin to be a "necessary" vehicle for interpreta- tion of 100-year compositions. Second: To make the fact clear that the listener is in the better posi-

240

\section*{VIOLIN TONE-PECULIARITIES.}
tion to judge first-violin tone-values. Basing con- clusions upon Dr. Thomas' statement, we are com- pelled to suppose that further interpretation of Haydn, Mozart, and Beethoven scores must cease with "the best of the Cremona violins." It is my conviction that this position by Dr. Thomas is the result of old violin aestheticism'run mad; and furth- er, that his position will be proven untenable in the near future. That Dr. Thomas' statement op- erates to excite contempt for the modern violin is clearly apparent; also, that his insistence in em- ploying none but old violins by the Thomas orches- tra operated to cause contempt for "the best Cre- monas" as is in evidence by the disappointing ex- periences of an army of patrons. There is no ques- tion about interpretations by the large orchestra being intended for the pleasure of the entire audi- ence, and not intended for the especial pleasure of the conductor and occupants of the front rows. There is also no question about disappointment coming upon the listener at not hearing first violin tone in the ensemble. I only speak for myself when saying that non-appearance of first-violin tone in orchestra ensemble is a non-artistic interpretation, whether the score be 100 years old, or but one hour old. In speaking thus, I do not mean that an of- fensive over-balance may be permitted to the strings. Mere loudness of violin tone without sweetness is offensive in both solo and orchestra violins. What I consider as the most valuable fea- tures in violin tone is moderate volume combined with sweetness, evenness of power, and marked

241

\section*{VIOLIN TONE-PECULIARITIES.}
intensity. Such combination of tone-qualities can- not become offensive to the listening ear; and, there is an ample supply of modern violins possessing such combination of tone- values for the needs of large orchestras called upon for interpretation of Haydn, Mozart, and Beethoven scores, or the scores of any master composer whatever. This point brings us to details for the production of maximum violin tone-power. From my observation, each of the following fac- tors must be at its best for the production of max- imum violin tone-power: 1. Sounding-board wood: Superlative spring- quality. 2. Graduation: Form securing greatest force of blow upon contained air. 3. Back: Density, fine texture, sufficient rig-

idity. 4. Arching: Form producing maximum concen- tration of sound-wave movement at the exits. 5. Bar: Medium density, modeled for greatest spring-action, position. 6. Exits: Area, position. 7. Interior surfaces: Permanent smoothness. 8. Ribs: Depth, rigidity. 9. Blocks and linings: Solidity. 10. Post: Mass, density, length, fitting, position. 11. Cubic capacity of body. 12. Varnish: Quantity, quality. 13. Bridge: Mass, density, height, span, scroll- work, position. 14. Finger-board: Length, width, height, obli-

242

\section*{VIOLIN TONE-PECULIARITIES.}
quity of under surface to exterior surface of sound- ing-board. 15. Strings: Diameter, proportion, number of strands, twist, material. 16. Bow stick: Length, weight, balance, cam- bre, spring-action. 17. Bow hair: Male equine, evenness, method of assembling. 18. Rosin: Quality, quantity. 19. Meteoric conditions: Moderate temperature, normal humidity. 20. Elevation above sea level. 21. Bow arm: Length, condition, training, en- thusiasm. From this list of factors, I place sounding-board wood at the head; and this conspicuous position is wholly due to the fact that I have not been able to produce maximum violin tone-power without sounding-board wood possessing superlative spring- action, and also, because I have not been able to find such wood every month, nor every year, nor every decade. In fact, but a few times in fifty years search have I found such wood; whereas, the finding of hard wood possessing density, and fine texture has been comparatively a very easy matter. Only limited observation is necessary to deter- mine that neither softest wood nor hardest wood yield maximum resonance; but, widest observation cannot pre-determine that any given sample of wood will yield maximum resonance in the com- pleted violin. Only after completion and trial can the degree of resonance be known for any violin

243

t

\section*{VIOLIN TONE-PECULIARITIES.}
from the hands of any maker whatever. Two samples of sounding-board wood may appear pre- cisely similar; yet, in the completed violin, there may be a wide difference in resonance, even when details of construction are similar. Thus, the pro- duction of maximum violin tone-power becomes a matter of difficulty; and, such difficulty is charge- able to inherent variations in the spring-action of wood. It is of interest to note the fact that violins of maximum resonance combined with "rich" quality of tone are rarities, even after 400 years of effort to produce them. Manifestly, there is some obstacle preventing production of such violins; otherwise, such violins would be in abundant supply. The cause for such rarity may not appear to others as it appears to me; but, all can agree upon the fact that such violins command a high price. In view of the fact that such violins are scarce today, we may safely presume that such violins will be scarce tomorrow, and remain scarce to an indefinite period. It is self-evident that the cause for such scarcity is something not subject to the command of man. It is a fact of every-day demonstration that spring-temper, either in wood or metal, is some- thing known only as a result of action following trial; and, such result cannot be pre-determined. Although the temperer of steel springs can pro- duce either "high," or "low" degrees of temper at will, yet, to produce any given number of steel springs possessing precisely similiar action is a

244

\section*{VIOLIN TONE-PECULIARITIES.}
matter of difficulty; and, such difficulty arises from two causes: (1) Variations in molecular arrange- ment in different samples of steel: (2) Inability to pre-determine such variations. In practice, it is clearly demonstrated that imperceptible variat- ions in density, or hardness of the steel spring, pro- duces perceptible variation in spring-action. Al- so, that value of the steel spring depends upon its temper, or hardness, or density, whichever term may be chosen. Naturally, we wonder what would be the result could violin sounding-board wood be given "tem- per" to the limit of desire. Shade of Colossus! Why then, the present-day violin of "big" tone would 'nt do for the baby; and, the largest auditor- ium would' nt do for a stage; and, the gender of Music would needs be changed to "he;" 'twould kill "her;" and, "he" must needs be deaf . Find users? Believe it. As I view this matter, Nature wisely limits man's possibilities in the violin. What I shall say about appearances of sounding- board wood promising maximum resonance of tone, combined with "rich" quality of tone, is presented only as the observation and conclusion of one indi- vidual; and, such conclusion is in nowise presented as a hard-and-fast rule. Upon the matter of re- sonance in sounding-board wood, I am bound to confess that infinite surprises have awaited my bow. It is observed that maximum resonance of

245

\section*{VIOLIN TONE-PECULIARITIES.}
tone is combined with richness of tone. I desire that this combination be kept in the foreground; because, loudness without sweetness reduces vio- lin tone- value to insignificance. No fact concern- ing violin tone stands more clearly in my view than the fact that sounding-board grain may be so dense as to make "rich" tone an impossibility; whereas, mere loudness of tone follows with ease. From the softer grades of pine, "rich" tone may be secured with considerable certainty; but, from those grades, marked power of tone is impossible. Again the fact is noted that great volume of tone does not mean great intensity of tone. Great vol- ume of tone may be obtained from the softer wood; but never great intensity of tone. The vast difference between volume of tone and intensity of tone is brought out into the lime light by the long-distance test in open air. This test de- monstrates conclusively that great volume of tone is not neccessary for great power of violin-tone. Indeed, I have not observed great volume of tone and great intensity of tone existing together in the same violin. In my observation, the physical appearances of sounding-board wood yielding the maximum re- sonance combined with "rich" tone-quality are thus: Genus pine. Maturity of tree before being felled. Sample free from heart-wood and sap-wood, knots, curls, cracks, and discolored spots. Grain perfectly straight, uniform in width, neith-

246

\section*{VIOLIN TONE-PECULIARITIES.}
er wide nor extremely narrow. Color, yellowish red. Connective tissue, between dense parts of grain, rather soft and flexible. Hard part of grain distinctly marked. Wood easily subject to splitting. Splits follow a straight line. Shavings rather brittle than tough. Shavings easily split up into threads. Difficulty in securing a smooth surface. Medium transmission of artificial light. Plate rather pliable than rigid. Bent plate returns to its point of rest with celer-

ity. Unprotected, heat and moisture rapidly develop a carpet of wood fiber. Disintegration comparatively rapid. Medium density of grain. Thus are the physical qualities of my ideal sound- ing-board wood. Such wood, with correct models of arching, correct graduation, correct area and po- sition of exits, and with permanent and perfectly smooth interior surfaces yields ample intensity of tone for all practical uses of the violin, besides yielding "rich" tone. As you remember, "rich" violin tone depends upon the presence of harmonic overtones, and harmonics a bassa, or resultant tones', and, that such harmonic tones cannot be coaxed in- to an appearance by the sounding-board of dense grain. There are degrees of density between the medi- um and the extremely dense which possess consid-

247

\section*{VIOLIN TONE-PECULIARITIES.}
erable value; and, in situations where noise abounds, the sounding-board of more than medium density possesses the greater value. In such situ- ations, those ethereal tones called harmonics are annihilated; even the principal tone may have to struggle for existance; therefore, the sounding- board possessing greatest power in spring-action gives the greater satisfaction in the presence of noise. But, I do not find increasing density of sounding- board wood to be continually followed by increas- ing power of tone. On the contrary, I find weak- ened tone-power following extreme density. This result might be pre- supposed; and, because increas- ing density is followed by increased rigidity; and, increased rigidity is followed by diminished ampli- tude of oscillation; and diminished amplitude of os- cillation delivers a blow of diminished force upon contained air; hence weakened tone follows. Method of graduation has much to do with vio- lin resonance. At this moment, graduation of the sounding board is in point. Thickness of the sounding-board may be so great or so reduced as to cause weakness of tone. As precise figures for violin-plate thickness are unreliable guides, therefore such figures are not presented. In my experience, it is clearly demonstrated that sound- ing-board thickness, producing greatest force of blows upon contained air, vary with each sample of sounding-board wocd; and, for this rea- son, hard-and-fast figures for thicknesses are un- reliable guides more frequently leading to disaster

248

\section*{VIOLIN TONE-PECULIARITIES.}
than to success. It is apparent that rigidity of the sounding-board should be determined by the force in the strings, since blows of the strings must be the only reli- ance for arousing sounding-board action. There- fore, the gauge of strings must be first determin- ed, and thereafter, rigidity of sounding-board, in each violin, must be reduced to that degree deliv- ering the maximum blow upon contained air, and, such degree of rigidity can be determined only by trial. Manifestly, for certainty in this matter, there is no alternative. The point is to know when each sounding-board is delivering its maximum blow. For this matter, I know of no hard-and-fast rule, unless it be the rule that experience must govern. I can truly say that much experience in reducing sounding-board rigidity leads one to err on the safe side rather than incur disaster by too great reduction of rigid- ity. It has been my experience that, when the ex- act degree of rigidity is secured, the further remov- al of thinnest shaving weakens the force of blow upon contained air. For securing the maximum blow upon contained air it is necessary that the extent of both normal and transverse vibratory action in the sounding- board be given most careful consideration. (See Appendix for explanation of normal and transverse vibration. ) As these actions are aroused by blows of the strings, and, as the strings vary in diameter and weight, therefore, rigidity of the sounding-board

249

\section*{VIOLIN TONE-PECULIARITIES.}
must vary with the force of the strings for produc- tion of maximum tone-power. As force in the fi- shing is less than that of the A, A less than D, and D less than G, therefore, it is not reasonable to expect maximum vibratory action to follow that method of sounding-board graduation giving equal rigidity of wood to all strings alike. In discussing that vibratory action which travels across sounding-board grain, the terms "trans- verse" and "tangential" may be interchangeably employed, since practically, each term has the same meaning. In practice, increasing the dis- tance traveled by transverse vibration in the violin sounding-board operates to increase volume of tone; and, such increase in volume of tone may be car- ried to a degree causing loss to intensity of tone. In the production of maximum tone-power this point is of great importance; because great volume of tone does not mean great power of tone. Great power of tone means the distance traveled by tone. In the violin sounding-board, the distance travel- ed by, transverse vibration depends upon three factors: 1. Power of strings. 2. Rigidity of the plate along lines at a right angle with the center-join. 3. Density of grain. Manifestly, it js necessary to determine first up- on the size, or ; gauge of strings. Second, to dimin- ish rigidity along lines at a right angle with the cen- ter-join to that, degree permitting the strings to .propel transverse vibration as far from the center-

250

\section*{VIOLIN TONE-PECULIARITIES.}
join as desired. Some violin makers reduce rigidi- ty along these lines to a degree permitting force of the strings to propel transverse vibration to the edge of the plate. This plan invariably operates to increase volume of tone; but, it also operates to diminish intensity of tone; and, the reason for such loss to intensity is clearly due to dispersion of string-force. It is my observation that, to produce maximum violin tone-power, string-force must be concentrated rather than dispersed; and therefore I limit the distance traveled by transverse vibration. [It is necessary to bear in mind the fact that I employ the gauge-2 string, because, were larger strings employed, greater force would be present; hence greater distance of travel might be safely permitted to transverse vibration.] There are two methods for determining the dis- tance to be traveled by transverse vibration in the violin sounding-board after its thickness has been nearly reduced to the final degree, thus: 1. By removal of wood from the interior surface before assembling. 2. By removal of wood from the exterior sur- face after assembling and stringing up. With the first method, results depend upon guess-work when working upon new and untried wood. With the second method, results are cer- tain; because with the strings in tune, their force may be tested until transverse vibration reaches any desired point upon lines at a right angle to the center- join. For concentration of string-force upon the sound-

251

\section*{VIOLIN TONE-PECULIARITIES.}
ing-board, the following 1 method for diminishing thickness has been the most successful in my ex- perience; because, with this method, I could limit the distance traveled by transverse vibration at the point yielding what I call sufficient volume of tone without the greater dispersion of string-force; and, the result has been moderate volume of tone with marked intensity of tone, a combination in- variably reaching the greater distance in the out- of-doors long-distance test. Thus: Greatest thickness at bridge-position; thence, thickness diminishes gradually to a point half-way from bridge-position to ends of plate; from such point, thickness increases gradually to ends of the plate; and, from such half-way point, on lines at a right angle with the center-join, the same thickness at such point is secured along such lateral lines to a point one and one-fourth inches to the right, and to a point one and one-fourth inch- es to the left of the center-join; and, from these points also thicknesses gradually increases in eith- er direction from these lateral lines. Next, the violin being assembled, and the strings in tune, the force of the strings is tested. Inten- tionally, the sounding-board is too rigid as assemb- led, as a matter of precaution. Now comes work in which science cannot become a substitute for musical sense; neither can the world's wealth purchase one grain of musical sense. Musical sense is in nowise a commercial commodity. Musical sense cannot be borrowed. There's the trouble.

252

\section*{VIOLIN TONE-PECULIARITIES.}
I can borrow my friend's books, and promptly forget to return them; but, I shall never have the chance to forget the returning of his musical sense. Upon testing the tone from the newly assembled violin, and finding duration of tone to be too short to suit my taste, then I increase the distance trav- eled by normal vibration; and, accomplish such re- sult by diminishing the thickness of the sounding- board from the half-way lateral line noted to ends of the plate; and, if I find volume of tone yet too small, then I increase distance traveled by trans- verse vibration; and accomplish this result by di- minishing thickness at, above, below, and outside of the lateral points noted; and, to give an equal chance for the force in the A and E-strings, thick- ness immediately beneath those strings is dimin- ished. In working upon new and untried wood, I repeat notice of the wisdom in erring upon the safe side rather than diminishing thickness until weakness of tone follows. This precaution is made necessary by increased flexibility of connective tissue inevit- ably following use. I have known such increase of flexibility to ruin tone-values in two years use. In the work of regulating sounding-board thick- nesses with the object of maximum tone- power in view, it is of importance to bear in mind the fact that sympathetic action between vibrating bodies diminishes as the square of the distance; therefore, thicknesses in the lower half of the sounding- board should be less than in the upper half. So far as work upon the sounding-board itself is

253

\section*{VIOLIN TONE-PECULIARITIES,}
concerned, these details include the main features; but, as the previous list shows, there are many other factors necessary in the production of maxi- mum violin tone-power. Because the more import- ant factors in the list have been presented upon preceeding pages, therefore extended notice of them is unnecessary at this moment further than to say that, in my experience, no factor in the list can be neglected without inviting disaster. There are numerous examples, other than the violin, showing the disastrous effects due to dispersion of force. The horn, too light in metal, yielding, trembling, "buckling" before the force from the player's lungs, loses tone-power. Instead of con- trolling the elastic energy of air molecules within, it is itself controlled. Again, the steam whistle yields but a harmonic tone when the pressure of steam greatly exceeds rigidity of the whistle. Again, the piano, with sounding-board too light, delivers but a weak tone; a result clearly due to dispersion of force. Again, the ball, thrown against the hanging can- vass, falls straight to the ground; a result due to dispersion of force. Again, the loud voice, directed upon the tele- phone diaphragm, causes indistinctness at the re- ceiver; a result due to dispersion of force. Again, the orator, in open air, can be heard only at limited distances owing to dispersion of force. Again, the fowling-piece, recoiling, loses range. Of these examples, it may be said, force is too

254

VIOLIN TONE-PECULIARITIES. great. True, from one standpoint, not true from anoth- er standpoint. The violinist, never wanting great- er tone-power than that produced by the weight of his bow, is not pleased with the violin requiring vigorous bow-pressure. To please such violinist, the violin maker must make a violin of light weight; and such violin is not pleasing to the vio- linist whose bow-arm possesses vigor and enthus- iasm, because the light-weight violin, subjected to vigorous bow-pressure, trembles to that degree permitting too great dispersion of string-force; a degree of dispersion resulting in weak tone. Loss of power is not the only affliction brought upon the violin by too light wood. A more disas- trous affliction is woody tone quality; the kind com- ing from hard wood when its rigidity is too greatly reduced. I have seen violins, built with the best of inten- tions, yet ruined by too great reduction of rigidity in the hard-wood plate. Such violins, under light bow-pressure, and with light strings, may yield a quality of tone to be tolerated; but, under vigorous bow-pressure, the tone is intolerably woody. As a musical instrument, such violin sinks below insig- nificance.

255

\section*{VIOLIN TONE-PECULIARITIES.}
LECTURE XVII. GENTLEMEN: Philosophic reasons for the perm- anent and perfectly smooth interior surface of the violin are clearly defined and easily comprehended. After long- years of observation upon this subject, it seems to me unnecessary to present the reasons for protecting interior surfaces of the violin from such disintegrating forces as heat and moisture. It also seems unnecessary to present reasons for aug- menting tone-power by the perfectly smooth inter- ior surface; but, remembering my own doubts upon these points previous to observations, therfore I assume that others, who are considering these points for the first time, may also entertain similar doubts. At first, my own doubts would not permit trial of this experiment upon violins of value; nor did I try it upon violins of value until after several years of observation. At this moment, I find myself wondering why the violin world should fear protec- tion for interior violin surfaces more than for ex- terior surfaces. That such fear exists to a wide extent is well known; but the reasons for such fear are not altogether clear; some holding to one ob- jection, some to another. Perhaps the most uni- versal objection is due to fear that interior protec- tion may be disastrous to tone-quality. The owner of a violin possessing superior tone-quality cannot be blamed for entertaining such fear, and, because violins of superior tone-value are comparatively rare, Were such violins as common as violins of inferior tone-quality, then the case would be differ-

256

\section*{VIOLIN TONE-PECULIARITIES.}
ent. Again, the case would be different could vi- olin users live long enough to see new violins wear out under their own bows. Thus, at the moment of knowing a valuable violin to be on the down- grade, the owner would willingly adopt any reason- able method for prolonging its usefullness. But, the wearing out of a violin within a single life-time is a rare occurrence. Basing calculations upon equal use, no argument is needed to determine that of two violins, built from the same sections of hard- wood and soft-wood logs, the one lighter in wood is destined to -the briefer period of usefullness. It is my own vivid recollection that when sounding-board rigidity is reduced to that point producing greatest augmen- tation of tone, thereafter but slight loss of wood by disintegration operates to diminish tone-power; and such loss, steadily continuing, because disinte- gration steadily continues, may bring ruin to superior tone- value within a single lifetime. Thus, we may expect immediate diminution of tone- value from such new violins as possess superior tone-power at the moment of leaving the builder's hands. Again, and for the very reason of steadily continuing disintegration upon unprotected interior surfaces, we may expect increasing tone-power from such new violins as are slightly heavy in wood at the moment of leaving the builder's hands, 'Tis but human to desire immediate superiority in violin tone-values. Life is too short to wait for slow disintegration. Once possessing superior violin tone-values, 'tis

257

\section*{VIOLIN TONE-PECULIARITIES.}
but human to mourn for the loss of them. To prevent such loss is but simplicity itself. I feel a pang of regret at every thought of those wrecked gems strewn along the violin's path; and, because nearly everyone of these wrecks, due to disintegration on interior surfaces, might have been preserved to an indefinite period. The details for violin interior surface protection, as best known to me, have been presented upon previous occasions; but, in such presentation, I hes- itated to discuss the principles of philosophy in- volved therein. My hesitation was due to the fact that but comparatively few violin users care for the philosophy involved in either the production or preservation of violin tone. It is quite safe to assume that all violin users are keenly alive to find brief and positive assertions that this thing or that thing, this method or that method infallibly pro- duces certain results in violin tone. It is my obser- vation, that even today, after the violin has been an object of attention during 400 years, none can formulate infallible rules for the production of greatest tone-values. When the day of such infal- lible rules arrives, violins of greatest tone-values will be as common as are violins of inferior tone- values today. The cold fact remains that, to make violins of faulty tone is of the utmost ease and cer- tainty; whereas, to make violins of faultless tone is a matter of the utmost difficulty and uncertain- ty. Such uncertainty is due to the impossibility of formulating rules governing all phenomena in vio- lin tone. Concerning the cause for such impossi-

258

\section*{VIOLIN TONE-PECULIARITIES.}
bility, there is disagreement between philosophers of 50 years ago and philosophers of today. Today, the majority hold the opinion that such impossibil- ity is due to the capricious action of sounding-board wood. My own experience decidedly lends corrob- oration to this opinion. It is my belief that today, given sounding-board wood of best quality, violins of best tone- values might be produced indefinitely. It is self-evident that the builder, even when given the best of wood, must necessarily be master of all other factors operating to modify violin tone. In my work upon the violin, it has been an object to separate and classify all factors concerned in the production and modification of violin tone thus: Class 1. Factors constant in tone-results. Class 2. Factors not constant in tone-results. The following tone-modifying principles, being constant in results are placed in Class I, thus: 1. Lengthening a tone-producing agent, other dimensions remaining equal, lowers tone-pitch. 2. Shortening a tone-producing agent, other di- mensions remaining equal, raises tone-pitch. 3. Increasing thickness of a tone-producing agent, other dimensions remaining equal, raises tone-pitch. 4. Diminishing thickness of a tone-producing agent, other dimensions remaining equal, lowers tone-pitch. 5. Lengthening perpendicular, confined, air columns lowers tone-pitch. 6. Shortening perpendicular, confined, air col- umns raises tone-pitch.

259

\section*{VIOLIN TONE-PECULIARITIES.}
7. Position of violin exits, distant from points of sound-wave concentration, diminishes volume of tone, and intensity of tone, and noisy tone-quality. 8. Increasing area of violin exits raises tone- pitch, increases volume of tone, diminishes inten- sity of tone, and accentuates noisy tone-quality. 9. Diminishing area of violin exits lowers tone- pitch, diminishes volume of tone, increases intensi- ty of tone, and diminishes noisy tone-quality. 10. Roughened and carpeted interior surfaces of the violin diminish volume of tone, intensity of tone, and noisy tone-quality. 11. Perfectly smooth interior surfaces of the vio- lin increase power of tone, and accentuate noisy tone-quality. The following tone-modifying factors, not being; constant in action, are placed in Class II, thus: 1. The finger-board. 2. The bridge. 3. The sounding-board. 4. The post. 5. The back. 6. The.bar. 7. The'blocks. 8. The ribs. 9. The linings. Some of the modifiers in Class II nearly approach constant action, or lack of action rather, as the finger-board, blocks, ribs and linings; but, because of lacking constant effects upon tone in any de- gree, therefore they are placed in this class. [At the moment of making up the list in Class II, I

260

\section*{VIOLIN TONE-PECULIARITIES.}
discovered that the effect upon tone, due to the ribs, linings and blocks, has not heretofore received at- tention. Because of this omission, I am now think- ing there yet may be other omissions. The difficul- ties under which I work are offered as an excuse. It is my observation that the benefit to tone from the ribs, linings and blocks consists in preventing violent trembling of the violin. Because such trembling weakens tone, therefore weight and rigidity in these modifiers is their measure of value to tone; and, this measure is therefore a measure for their lack of action.] It is observed that all tone-modifiers in class I, depend wholly upon the action of air; therefore, without the shadow of doubt, their constant effect upon violin tone is due to the fact that the action of air is a constant quantity. The first six tone- modifiers in class I, affect only the single tone- quality of pitch; the seventh affects three qualities of tone; the eighth and ninth, four tone-qualities each; the tenth and eleventh, three tone-qualities each. At this point, it is interesting to note the num- ber of tone-qualities at the absolute command of the builder. I find this number not great enough to be flattering. Enumerating tone-qualities en- tering into "rich" violin tone, I find their number to equal 12, thus: Pitch of tone, volume of .tone, intensity of tone, evenness of tone, freedom from dissonant overtones, or noise, sympathy in concert, responsiveness to bow-pressure, agreeable double- stop tones, harmonic tones, resultant tones, or

261

\section*{VIOLIN TONE-PECULIARITIES.}
harmonics a bassa, brilliance in velocity, human quality of tone. [Nobility of tone, and liquidity of tone strongly appeal for place in this list, but, as nobility of tone depends upon tone-pitch and volume of tone; and, liquidity of tone depends partly upon tone-pitch and partly upon brilliance of tone, therefore they are not given place as individual tone-qualities.] Out of the 12 individual tone-qualities given, I have found but two which are absolutely at com- mand of the builder at all times. Those two are the qualities of pitch and volume. Thus, no mat- ter what the tone-value of wood may be, violins equal in tone-pitch, and, equal in volume of tone may be reproduced indefinitely. To reproduce the remaining 10 individual tone-qualities at will, and at all times is wherein lies a great difficulty. Of the remaining number, I find three which are ap- proximately at command of the builder, namely, intensity of tone, evenness of tone-power, and freedom from dissonant overtones. Intensity of violin tone is found to be a product of four factors, thus: 1. Arching of plates. 2. Position and area of exits. 3. Condition of interior surfaces. 4. Inherent spring-action of sounding-board wood. It is apparent that three of these factors are at command of the builder at all times; but, the fourth factor is not at command at all times. Inherent spring-action of sounding, board wood, being a ca-

262

\section*{VIOLIN TONE-PECULIARITIES.}
pricious quantity, operates to defeat equality in the carrying power of different violins. Only in an approximate degree can the experienced worker in sounding-board wood predetermine the tonal- quality of any given sample. Tonal quality of wood can be determined only by trial; but, the build- er can secure considerabe intensity to violin tone by arching, by area and position of exits, and by con- dition of interior surfaces. Even with sounding- board wood lacking much of possessing superlative spring-action, it is my observation that greater tone-intensity can be secured by the perfectly smooth interior surface than with sounding-board wood possessing superlative spring-action while interior surfaces remain rough and carpeted with wood-fiber, wood-dust and dirt. But, the 80 per cent, increase in intensity of tone, as recorded, should not be taken as an amount of increase whol- ly due to the perfectly smooth interior surface. This record was secured to six common violins se- lected from a general stock, and their average carrying-power was first determined without change of their conditions; and, the 80 per cent increase in carrying-power was secured after such tone-modifiers as graduation, bar, post, depth of ribs, linings, blocks, condition of interior surfaces area of exits, finger-board, bridge, varnish, and strings had received my utmost attention. Thus, the per cent of increase in intensity of tone, due to the perfectly smooth interior surface alone, does not appear in the record, nor have I made experi- ment with this factor alone. The precise value of

263

\section*{VIOLIN TONE-PECULIARITIES.}
this important factor for augmenting intensity of violin tone ought to be determined, because of its interest to owners of old violins having tone-quali- ties right in every way except in the quality of in- tensity. Although removal of the dust-carpet from inter- ior surfaces of the violin, by such agents as shelled corn, wheat, or oats, is better than no removal at all, yet, there are violin owners who have not the courage even to permit entry of these harmless agents into that sacred interior. So far as my ob- servation extends, such fear is caused by the ex- pectation that noisy tone-quality may follow. It is a fact that there are cases wherein such fear has good ground for existence. Thus, when a violin is afflicted with noisy tone-quality, placing a carpet upon its interior surfaces does operate to "im- prove" the tone; and, because whatever operates to diminish tone-power may operate to annihilate noise. This phenomen in tone is due to the fact that the noise-wave is inherently weaker than the music- wave; hence, the noise- wave is first to disap- pear. The permanent, and perfectly smooth interior surface appeals only to such violin users as desire greater intensity of tone. To them, the philoso- phy involved in this question possess interest; and, to them I now address myself. In discussing the principles of philosophy per- taining to violin interior surfaces, it is interesting to note certain facts concerning air itself, thus; In all countries, philosophers agree that air, sur-

264

\section*{VIOLIN TONE-PECULIARITIES.}
rounding earth, exists as gases; that such gases are in the form of spheres; that each sphere is a molecule; that such molecules are too small for measurement; that they touch each other as shot; that they are compressible; that when released from compression, they resume the spherical form with energy; that force is communicated from mol- ecule to molecule by their expansion and contract- ion; that, in unconfined air, energy is dispersed equally in all directions; that such lines of disper- sion may be concentrated by confining walls; that concentration of such lines increases distance trav- eled by the energy in the original blow; that such energy may be reflected from solid, smooth sur- faces without loss of force; that such energy is not only arrested, but, may be totally annihilated by striking upon soft bodies; that force of the origin- al blow travels at a right angle to the striking sur- .f ace; that such line of travel, striking upon a solid, smooth surface, is reflected at an angle equal to the angle of incidence; that volume of sound is proportionate to the number of molecules affected by the striking agent together with the amount of force in its blows; that intensity of tone, (carry- ing power of tone, ) is proportionate to the force of blow and the degree of concentration given to lines of sound-wave travel; that the distance traveled by sound-waves may be diminished by conditions of reflecting media, and by meteoric conditions. Other qualities of air, not being in point are omitted. Each of these facts concerning air is di- rectly connected with the production of violin tone.

265

\section*{VIOLIN TONE-PECULIARITIES.}
The infinitely small size of the air molecule possess- es vast interest to the violin student. This inter- est centers in the fact that unprotected, interior surfaces of the violin cannot be made perfect re- flecting media for lines of sound-wave movement. Beneath the magnifying glass, and at the limit of polishing, these surfaces yet present ridges, valleys, and open mouths of caverns. It is self-evident that each of those ridges and valleys operate to abruptly deflect lines of sound-wave movement; in other words, to defeat equality in the angles of incidence and reflection; that those open caverns operate to annihiliate sound-wave movement; that those disintegrating forces, heat and moisture, op- erate to sharpen those ridges, to deepen those val- leys, and to increase the capacity of those caverns; that all of such agents acting upon interior surfaces of the violin, combine to diminish concentration of sound-wave movement at the exits; that such di- minution defeats intensity of tone; that all of this effect upon tone is due to the infinitely small size of the air molecule, and to the infinitely rough, un- protected surfaces within the violin. From the permanent and perfectly smooth inter- ior surface, the value to intensity of violin tone- power is plainly apparent. It is also apparent that to secure such surface, dependence must be placed upon a protecting agent susceptible of taking a high polish. Aside from intensity of tone, there are, in this connection, other effects possessing interest to such violin users as desire maximum tone-power. I al-

266

\section*{VIOLIN TONE-PECULIARITIES.}
lude to such effects as are produced by: 1. Times of sound-wave reflection before reach- ing the exits. 2. The solid, or yielding back as a reflecting medium. 3. Condition of the striking surface. It is my observation that the number of times sound-waves are reflected before reaching the ex- its exerts very perceptible influence upon both brilliance of tone, (meaning distinctness of tone in rapidity of succession,) and power of tone. Obvi- ously, times of reflection may be so great as to di- minish force in any projectile; also, it is obvious that times of reflection operate to delay any projec- tile in reaching a given point. In the violin, these principles are demonstrated by the method of sounding-board graduation. Thus, when reduction in thickness places the point of widest amplitude of oscillation near to ends of the plate, then both brilliance of tone, and power of tone are diminish- ed. It is self-evident that the greatest force in- blows from the plate is found at the point of widest amplitude of oscillation; also, that times of reflection of sound-waves increases as the distance from point of origin to exits. Oft repeated exper- iment in locating the point of widest oscillation has led me to the conclusion that to this method of sounding-board graduation, maximum power of tone is an impossibility. In such experiments, it was conclusively shown that placing the point of widest oscillation half way from position of bridge to ends of plate produced maximum tone-

267

\section*{VIOLIN TONE-PECULIARITIES.}
power. From my view point, the increase in tone- power following the latter method of graduation is due to two facts, thus: The widest oscillation, possible to sounding-board fibers, can be secured only at a point half-way from the bridge-position to ends of the plate; and, shortening the distance from the point of origin to the exits, by diminish- ing times of reflection, diminishes both delay and and loss of force to the sound-waves receiving the greater blow; hence, greater brilliance of tone, and greater tone-power. It is an easy matter to demonstrate on paper that, as arching increases, times of sound-wave re- flections diminish. Indeed, this proposition is self- evident, thus: In the box fiddle, the plates are parallel, hence, sound-waves, originating at the top plate, and traveling at a right angle to the .striking agent, according to the law, must touch the back at a point perpendicular to the point of origin; therefore, the reflected wave must travel .back to the top plate directly upon the line of incidence; hence, there is no progression of sound-wave movement toward the exits. If the top plate alone be given but the slight arching of i then sound-waves, originating therefrom, will not touch the back at a point perpendicular to the point of origin, but, will strike at a ' point nearer the exits; and, will be.reflected at an angle equal to the angle of incidence according to law. Thus, sound-wave movement, at each reflection, approach- es the exits; but, in this case reaches the exits on- ly after many times reflection. If now the back

268

\section*{VIOLIN TONE-PECULIARITIES.}
be given an arching of 1-8, times reflection are di- minished by 1-2. It is thus apparent that times reflection diminish as the height of arching; but, it does not follow that tone-power follows an in- definite degree of arching because increasing the heighth of arching increases resistance; therefore the arching of violin plates may be so great that force in the strings cannot overcome such resist- ance. Experience, and but experience, determines that when widest oscillation of sounding-board fibers is placed near the ends of the plate, the arching, at bridge-position, should not be less than t to secure satisfactory brilliance of tone; and, when the point of widest oscillation is placed half-way from bridge- position to ends of the plate, then arching should not be less than i inch. It was previously shown that the less the heighth of arching, the greater the susceptibility to force, therefore heighth 7 of arching becomes a factor in the production of max- imum tone-power. As the sounding-board strikes the blow upon contained air to produce sound, therefore its strik- ing surface possesses interest in this connection. It is evident that force in any blow may be diminish- ed by a soft carpet upon the face of the striking medium. Thus, when the interior surface of the sounding-board is covered with a carpet of finely slivered wood-fiber, force in its blows upon con- tained air is diminished in proportion to the thick- ness of such carpet; or, stated in the reverse way, force in its blows is augmented by the permanent,

269

\section*{VIOLIN TONE-PECULIARITIES.}
and perfectly smooth interior surface. In this connection, the back plate possesses interest aside from the condition of its interior surface. As pre- viously stated, I find myself able to secure greater intensity of tone by treating 1 the back solely as a reflecting medium. Thus, rigidity in the back, sufficient to stand firmly against the force in charg- ing sound-waves, becomes a feature of value. But, whether or not the strong back stands still in the presence of such charge, I do not know. I only know that violent trembling of the back operates to diminish intensity of tone, as is easily demon- strated by the long distance test in open air. Right here is a fact possessing double interest, thus: In making the long-distance test for carrying power, powerful bow -pressure is employed; therefore, the back is called upon to stand before the limit of force within that particular sounding-board; and, yielding of the back to such amount of force oper- ates to diminish the recoil in air molecules. It is apparent that such recoil diminishes as the yielding of the back, and as yielding of the back diminish- es with the diminished force in charging sound- waves, therefore, with diminished bow-pressure, such violin displays greater carrying power than with greatest bow-pressure. This phenomenon may be observed frequently. So far as assistance in orchestra ensemble is concerned, such violin is worthless; and worthless because its tone is smoth- ered by sound-waves from harmony instruments. Many times have I demonstrated the fact that re- placing such weakened backs by others possessing

270

\section*{VIOLIN TONE-PECULIARITIES.}
rigidity operates to increase intensity of tone. It seems clear to me that rigidity in the back should equal the limit of force in the sounding-board, even when the sounding-board is aroused to its widest oscillation by action of the open strings under greatest bow-pressure; and, because only thus may the limit of bow-pressure be employed without dis- aster to carrying-power. I will demonstrate the reason for such conclusion by the action of this elastic ball as it rebounds from yonder brick wall, and, from yonder partition wall of light wood, As the ball is thrown against the brick wall, the distance to which it recoils is proportionate to its force at the instant of impact; thus, the greater its force, the greater its recoil. But, from the partion wall of light wood the result is widely different. Beginning with moderate force to the ball, and gradually increasing until its force equals rigidity in the wood, the distance of recoil in- creases as the amount of force; but, upon giving to the ball an amount of force exceeding rigidity in the light wood, the distance of recoil is greatly diminished. It is apparent that yielding of the wall robs the ball of its elastic energy. Air molecules are elastic balls; and, the yielding violin back robs them of a chance to display their elastic energy; hence, diminished carrying power. But, did rigidity in the wall exactly equal the greatest force given to the ball, then the distance of recoil would have exceeded former distances, because elastic energy in the wood, added to elastic energy in the ball operates to increase the distance

271

\section*{VIOLIN TONE-PECULIARITIES.}
of recoil. Form the action of the partition wall, it is apparent that such great increase in distance of recoil is possible only to one certain amount of force in the ball; a less amount of force not being able to arouse the same degree of energy in the wood; while a greater amount of force in the ball causes yielding of the wood with its disastrous results. Right here is a point of vast interest to the violin student. 'Tis the old question of regulating violin tone- power by work upon the back plate. The follow- ing conclusions, although based upon repeated practical demonstrations, are presented only as the conclusions of one individual; and, such conclus- ions may have to stand without support from oth- ers. It is my observation that rigidity of the back may be so determined as to augment the funda- mental, or open tones of each string; but, in no in- stance have I observed equal augmentation follow- ing for any other tones. The reason for this phe- nomenon seem clearly indicated by action of the ball and the light partition wall. Thus: When force in the ball exactly equals spring action in the wall, then recoil of the ball reaches the maximum, but, with less force in the ball, recoil diminishes. Applied to the violin thus: Development of great- est force in the strings demands employment of their entire length; therefore, shortening the strings operates to diminish force; therefore, when rigidity of the back exactly equals force in the en- tire length of strings, it follows that tones produc- ed by shortening the strings cannot be equally

272

VIOLIN TONE-PECULIARITIES. augmented. It goes without saying that such vio- lin as possesses powerful open tones, while marked weakness of tone follows successive shortening of the strings, is a violin of but insignificant tone- val- ue. Thus maximum violin tone-power, and perfect smoothness of violin interior surfaces are insepar- ably linked together; and perfect smoothness of violin wood is an impossibility.

273

\section*{VIOLIN TONE-PECULIARITIES.}
LECTURE XVIII. GENTLEMEN OF THE VIOLIN STUDENT CLUB The : present hour marks the close of our course. Here are two pieces of wood possessing interest; this one is soft Michigan pine; the other is Michigan white cedar; and both are presented for the purpose of demonstrating the fact that oil may be applied up- on even soft wood without being followed by pene- tration. In works upon the violin, we may often read that employment of oil as a protecting agent is hazardous because oil penetrates the wood. Again, we may read that permeation of the wood by Cremona varnish is the reason for tone values in the Cremona violin; also, that Cremona varnish is excuse me Mr. Promoter was an oil varnish. Because of the conflicting evidence in works on the violin, I could not reach a conclusion upon this im- portant point without a test being made under my own observation. It is apparent, without any test whatever, that penetration of oil must operate to diminish the ra- pidity with which the wooden spring returns to its point of rest; hence, penetration of oil into the sounding-board must operate to diminish violin tone-power. I reasoned thus: If oil can be ap- plied to these two samples of soft wood without penetration, then oil can be applied in safety to any violin. To make such test reliable, I employed raw oil, because it dries more slowly than boiled oil; also, with raw oil, I mixed gum mastic, the slowest drying gum known to me. Nothing else was put in this mixture. The surface of the wood

274

\section*{VIOLIN TONE-PECULIARITIES.}
was made smooth, but, no filling was employed. Application of this mixture was effected with the rubbing pad, and in attenuated layers. Each lay- er was given time to dry before application of succeeding layers. It is now more than ten years since these two pieces of soft wood were thus finished. With a knife, I scrape off the finish down to the wood for the purpose of showing that enough of this mixture was applied for protection. With this sharp-cutt- ing blade, I now remove a thin shaving immediate- beneath the finish. As you observe, the wood is bright; and without any discoloration whatever. I attribute this fact to the manner of application. It is my observation that a layer of oil dries with a degree of rapidity proportionate to its thickness. Thus, the attenuated layer, possible to the rubbing pad, dries in much less time than the lightest lay- er possible to the brush. Applied in attenuated layers, it is my belief that less penetration of wood follows application of oil varnish by the rubbing process than the penetration of alcohol following application of spirit varnish by the brush. My reasons for preferring oil varnish upon the violin are thus: 1. Less penetration of wood. 2. Greater attenuation. 3. Greater elasticity. Sadly I now take a parting shot at the post. This innocent appearing thing is an ever ready butt for both the pen and the tongue. Although the post is at once an object of jeers, scoffs, rail-

275

\section*{VIOLIN TONE-PECULIARITIES.}
ings and maledictions, yet its worth defies compu- tation. Its adjustment and readjustment occupy idle moments of every violin user on earth. Al- though kept; "on the move," the patience of the post is endless. Having such anatomical parts as "latete," (the head) and "la pied" (the foot) yet, the post stands equally well upon either end. In superlative fancy, the French surpass all other peoples by naming the post "1'aime" (the soul.) Oh thou Post! Thy very name Doth make us think thou art but tame; Yet, thou art the very thing Conferring title to the king.

Small thou art, of little space, Nor worth appears upon thy face, Yet, thou art a precious thought By India's wealth 'twere never bo't. To name thee "soul" doth harm to none, 'Tis A thine by right of giving tone To and E which we do hold Above the worth of Croesus' gold. We are now well across the "dry district." As you remember, we were promised something for parched lips at this moment; and, if your minds are as my mind, we will celebrate this occasion by introducing our feast of the pass-over. Have we not passed over impregnable difficult- ies? Are we not also "dry?" Verily, this dust along the violin's path is chok- ing. 'Tis not best to give our entire time to dryness;

276

\section*{VIOLIN TONE-PECULIARITIES.}
for without an occasional smile, life is scarcely worth the living. If real work is enjoyment of play, then, in all conscience, we are prepared for a moment of recreation. In amount of work, I k,now of no occupation surpassing the mastery of violin possibilities. Even the reading of it is tire- some. Although 400 years have been given to con- sideration of those possibilities, yet, this question remains new to each succeeding generation. Will it ever remain new? There's no doubt. There's a bewitching power in the enchantment of violin tone surpassing the figures in arithmetic. The daily offering of wealth to the "king" surpass- es all offerings to all other sorcerers combined. Although the germ theory of today covers the earth; although the germ hunter invades the very springs of life for his "cultures;" although he has segregated and named countless pathogenic and saprohytic micrococci, yet, in an unaccountable way, he has neglected to focus on the violinicus universalis. Considering the epidemic proclivities of this germ, such neglect is astounding. It is found all around the earth, and in all latitudes. It knows no prejudice for creeds, nor races, nor col- ors. Possibly, its neglect by the microscopist may be due to the fact that this germ is visible by sun- light, by moonlight, by starlight, self-luminous in darkness and never suffers eclipse. The violini- cus universalis is ineradicable. In all climates, its presence is manifested by uniform symptoms, to- wit; intoxication and indifference to wealth. Nor

277

\section*{VIOLIN TONE-PECULIARITIES.}
wealth, nor rank, nor position whatever secures immunity from this irrepressible germ. The mil- lionaire and the beggar jostle each other for seats beneath the artist's bow. Watch them at the mo- ment when the first sweetly tender, soulfully in- tense tone reaches their hearing. From that mo- ment, neither moves, neither breathes. Were eyes closed they remain closed. Were lips parted they remain parted. Both are drinking drinking in the sweetest sound on earth. Intoxicated? Ask yourself. Moreover, the more they drink, the more they want. Music! Beautiful Music! All around the world, from hut to palace, thou art welcome. Thy lang- uage is understood by the whole world. As the sunshine, thy presence brings warmth. Thy pow- er to touch human heart hath no parallel. Thy eloquence commands silence, and silence is willing. Thy subjects pass beyond mankind even unto the animal kingdom. Thy worship is boundless, and endless, even as the throbbing of hearts. What is music? For one thing, it is a demonstration for J\textasciicircum\{\} the ease with which the grinning idiot may question the sage. Thinking perchance you may be a sage, and not seeing danger ahead, you begin with: "Music is something built from nothing tangible; a product of genius working in ether for aBstheticism; an en- tity unknown but to hearing; matter without pond- erosity; inconceivable material from intangible

278

\section*{VIOLIN TONE-PECULIARITIES.}
realms of thought and and ." As you glance at the broadening grin of yonder idiot, you avert your face in discomforture. By the record, antiquity of music is great. The record carries the student back almost to the day when Adam, in Eden, became manager of opera. Genesis iv:21. Speaking of the descendants of Cain: And his brother's name was Jubal. He was the father of all such as handle the harp and organ. Thus we claim descent from Jubal. We are proud to know that our root handled the harp, for the harp is of strings. [We do but glance sidewise upon the organ.] We do especially delight in such of Jubal 's de- scendants as handle the strings. We are pleased with their soft voices and quiet ways. They live in a world by themselves and converse with the eyes and the strings. 'Tis true, they are distant- ly related to Cain; 'tis also true that "blood will tell;" yet, they never "murder music," neither do they "raise Cain." [Can I ever be forgiven?] Than the growth of music, there is nothing in history showing equal deliberation. Tis true that things of slower growth last the longer. Among other assets, Solomon possessed a band of four thousand trumpeters. From the best light obtain- able, the scores interpreted by this magnificent band were limited to four tones. Developement of our major and minor scales required the time from Ju- bal to Palestrina in the 16th century of the Christ-

279

\section*{VIOLIN TONE-PECULIARITIES.}
ian era. We owe much to Palestrina. Prior to Palestrina, that wonderful thing, called "har- mony" did not exist, nor could it have been em- ployed had it existed. In scale building, the gen- ius of Palestrina made the modern orchestra a pos- sibility. Today, loss of the symphony orchestra would put the world in mourning. Six thousand years from Jubal! Verily, as an example of slow growth, music permits no rival. But, development of music is not equally distributed. There are yet localities where the method of six-thousand-year Jubal re- mains in pristine purity. I know of nothing poss- essing greater interest than the moment when Ju- bal is brought face to face with an interpretation of music by the symphony orchestra of today. You remember the recent presentation of the South Sea Islander's theater its one-piece orches- trathe one hollow log drum its one drum head the one drum stick the one son of Jubal who handled the stick his inimitable look of ecstacy? After mastering his da capos, and dal segnos with my very dolcissimo I asked him, "Would you please change the tonic?" He answered never a word. There was no need. His soul-lighted eyes looked into my eyes, and from his eyes came signals older by thousands of years than the signals of Solomon. Within those illumined orbs I easily read, "'Tis all the key I know." Our musical sense is the product of six-thousand-

280

\section*{VIOLIN TONE-PECULIARITIES.}
year culture. The musical sense of this son of Ju- bal is yet at the zero point. To him, the dreary monotone from his primitive instrument of percus- sion, (probably the organ, Gen. iv:21) is sufficient to translate him into the regions of sweet dreams; and, from his dreamland point of observa- tion, the change to our major and minor scales, the change to the wide range of melody tones, the change to 24 keys, the change to harmony parts, the change to interpretation by a score of varying devices is to him but a change to incomprehens- ible noise. Of such in heaven, to him 'twere hell. To us, his ''organ" is hell. Yet, heaven cannot be denied him. There is but one alternative petition petition that distances in heaven be great enough to accom- modate all such as handle the harp and organ. Thus all may remain in the union. But yet, there's another matter for serious thought. Those very distances and planes for which we petition may cause trouble to ourselves. Thus: While here below, one trait of humanity is egotism; and, when seized of egotism, we are sure to attempt seizure of the highest seat. As egotism is blinding to introspection, therefore disappoint- ment awaits some of us who think that we can handle the harp. The "maestro" may feel certain of entering at the highest gate; but, do earthly standards compare with heavenly standards? That's the question. Even the bow of the "wa2sro" may be too rough for angel ear. That

281

\section*{VIOLIN TONE-PECULIARITIES.}
we shall try the highest gate is possible. That some of us will be turned down is probable. The "maestro" here may only enter a lower gate up there. That the violin will enter at the highest gate, we may feel assured, because, only the strings can accompany angel voice. But, we have assurances for hope. These assurances point to cer- tain guide-boards; and, these guide-boards point the way. Thus, the "msestro" may at least arrive at the highest gate; that he will enter therein de- pends upon precautions. If the world owes much to Palestrina, it owes more to the church which, never for a moment, relaxes her fostering care for Beautiful Music. The brave deserve to live. The modern violin builder's 200-year fight for life has no parallel in history. Considering that only two human traits have been the cause for this prolonged struggle, we are astounded at the vast height, and the vast depth of those two traits. Specifically, those traits are: Greed for profits. Credulity of consumer. Upon this occasion, greed for profits is much less in point than credulity of consumer; and, be- cause of the fact that equal greed for profits is equally displayed in other lines of commerce; but, equal credulity is not found in any other consumer. No fact is more patent in violin history than the fact that human credulity is colossal. Since Stradi- varius and Joseph Guarnerius pointed to the way, countless numbers of violins, equal in tone-value,

282

\section*{VIOLIN TONE-PECULIARITIES.}
have been built; but the builders have received no credit whatever. The cause for such builders not receiving credit for violins of tone-value equalling the best Strad, or the best "Joseph" is a matter worthy of thought. For this cause, there are but three classes upon which responsibility may be laid, thus: 1. The violin promoter. 2. The violin consumer. 3. The violin builder. The violin promoter interests himself only in the product of such builders as have acquired fame, and because otherwise, profits are insignificant. For large profits, the promoter depends upon con- sumers' sentiment alone; and, consumers' senti- ment is as consumers' credulity; and, consumers' credulity is colossal. Next come the builders them- selvesa procession a funeral procession appar- ently yes, in reality. Why for such mourning? Mr. Violin Builder, because of friendliness for you, my scalpel is bent upon reaching the seat of your trouble. Don't wince. 'Tis for your good. Therefore, I diagnose the cause of your mourning to be that fatal malady called ' 'imitation. " 'Tis imitation that buried your hopes. J. B. Vuillaume leads your procession. When the Strad violin came into demand, Vuillaume built Strad violins; sold them; sold them to experts in tone- values; proof enough that the Vuillaume pro- duct equalled the Strad product. Fatal mistake!

283

\section*{VIOLIN TONE-PECULIARITIES.}
Had those violins been honestly labeled, "J. B. Vuillaume, a Paris," then Paris today would be equally famous with Cremona. From Vuillaume, the procession lengthens rapid- ly; consequently Strad violins are found in every village throughout Europe and North America, and, lately they are heard from in China. When the credulous find themselves duped, credulity is followed by wrath; and, when wrathy, the dupe will not accept your imitation as a gift, even did its tone-value excel the best ever. The brave deserve to live; and the day is rapidly breaking for those courageous violin builders dar- ing to place their names upon labels. Even in high places, merit for the modern violin is now admitted, and the hand of the old- violin promoter is losing its nerve. From such accredited authority as Chas. Reade, London, comes the statement that the best Strad, stripped of its varnish, is worth today but \$125. This is an acknowledgement that violins of equal, or even greater tone-value than the Strad are existing in abundance. Indeed, new violins are commanding prices undreamed of 200 years ago. The fate of the Cremona ' 'gusher" is a sad feat- ure accompanying this change in sentiment. As the word "gusher" is applied to different objects, 'tis therefore necessary to state specifically that its employment in this connection does not indicate those mammoth geysers in Yellowstone Park, but does indicate the "Cremona gusher." There's a difference. The Yellowstone gusher displays ac- tion of deep water. The "Cremona gusher" dis-

284

\section*{VIOLIN TONE-PECULIARITIES.}
plays action of shallow water. If anyone ever had an occasion to pray for deliverance from friends, 'tis Antonius Stradivarius. A man is a man nev- er more sometimes less. Soon the Cremona "gusher" will lose his occupa- tion; and, for his fate, my pen sheds tears of inky blackness. With nothing upon which to expend his imagination, the Cremona "gusher" must die of appolexy. Sadly we turn away from the ' 'gush- er's" sickening exhibition of shallowness; and, our degree of nausea is exactly proportionate to our degree of esteem for the object of his "gush." An account in detail of Cremona "gush" is too much 'twould be an overdose inadmissable. Doubt- ing the ; ability of my pen to give the Cremona "gusher" new pointers, yet, because of its desire, 'twill be permitted trial. In such trial, 'tis but necessary to call up a single specimen of "gush." This specimen reads, "He, (Stradivarius) married the wealthy widow, Signora Capra, and thereafter pursued his chosen vocation from a purely artistic standpoint. ' ' Certainly! Ditto every poor devil of a fiddle maker when lifted above the pinching demands of quick-return fever. Although not intended by the "gusher," yet this statement virtually declares that Stradi- varius knew the pinching demands of quick-return fever prior to his adoption by that benevolent widow. My pen points to the fact that the "gusher" has entirely ignored the Signora Capra's equity in

285

\section*{VIOLIN TONE-PECULIARITIES.}
Strad's glory. A Joseph Guarnerius in jail! And in jail for a paltry debt! Tis sad reading today in view of the fact that the modern violin builder is taking but little notice of \$500 for a single violin; but, the neglect of the "gusher" for the jailer's daughter surpasses mere sadness. It does seem that even the Cremona "gusher" might understand that the "divinity which doth shape our ends," in lack of a benevo- lent widow, placed Joseph in jail for the purpose of giving him the services of the jailer's daughter in procuring the finest material. As a "regular," I am prepared to prescribe for you poor fellows, thus; As a reliable, anti-febrile treatment for "quick return" fever, take a rich widow a la Stradivarius. I sincerely hope that the Cremona "gusher" may yet retrieve lost opportunity to grant justice to the silent partners of Cremona's genii; but, for the sake of the living I sincerely pray that the ' 'gush- er" may grant rest to Stradivarius the honest, con- scientious, ambitious, tireless, violin-loving, violin- playing, violin-making man nothing more 1644 1737 equaled every day in the year One Thous- and Nine hundred and Five. 'Tis said the superior violin is a product of gen- ius. The Hon. W. E. Gladstone said the evolution of the violin cost more thought than evolution of the steam engine. In all human activities, I know of none possess- ing equal difficulty in commanding superiority of

286

\section*{VIOLIN TONE-PECULIARITIES.}
product with the art of violin making. I know of no route to superiority of violin production other than experience combined with intense application. To me, the superior violin appears as a product of su- perior mechanical skill combined with superior musical sense and intense application, all directed upon superior material nothing more nothing less. Than the art of violin making, I know of no hu- man activity possessing equal temptation for fraud. Within the violin lies much of the work upon which value depends; and, because much of its interior is hidden from view, fraud steps in. None can look only upon the exterior of a violin and determine that fraud does not sit within; yet, every purchaser believes that the violin always im- proves with age and use. 'Tis the universal be- lief that the violin always improves with age and use. The universal belief in this delusion makes my heart ache. Such belief is an ignis fatuous daily leading its victims into the Slough of De- spond; and, the procession thereto is a multitude. All around the world, in daily increasing volume, the tenderest feelings of human hearts are poured out at the feet of Beautiful Music. That fraud is permitted to thrive upon such tender feelings makes the devil wild with joy. You who have carefully read these pages are prepared to comprehend the difficulties surround- ing the path of the conscientious violin maker. You are prepared for the fact that the most skill- ful workman, try as he may, occasionally must

287

VIOLIN TONE-PECULIARITIES. meet more or less of defeat. From history, you learn that, even to the superior workman, apprec- iation is slow in coming. In the presence of this fact, and, in the presence of pinching want, and, from my own view-point, that violin maker who resolutely turns his back upon temptation to de- fraud is a person cast in heroic mold. Of such stuff are heroes made. As you remember the past, as you honor justice, and, as you love music, wherever and whenever you find such hero, you will hasten to place the crown of merit upon his brow while he yet lives; and, in the bestowal of honors, as you honor fair dealing, you will not forget the equity due his si- lent partner. Throughout the gloom of waiting days her voice has been his cheer. Grant her hon- ors. Beautiful Music! Thou canst save When others fail. To me thou gave The pow'r, th' will, th' tho't to move From put th' stubborn sinner's groove. Beautiful Music! Thou art th' leav'n Op'ning wide the door to heav'n. Love I thee? E'n to th' end, With thee to heav'n, my way I wend. Thou, Music! At thy feet, my tribute lies.

288

\section*{VIOLIN TONE-PECULIARITIES.}
APPENDIX. In every life, explanations for technical terms are desirable at one period or another; and, such period may come during youth, or it may be delay- ed until advanced age. I find it impossible to write at length upon violin tone-peculiarities without employment of some technical terms; and, remem- bering the period when technical language was confusing to myself, and remembering my pleasure accompaning elucidation of such language, and, remembering my desire that all readers receive re- ward for the punishment incidental to wading through the dry details on preceeding pages, there- fore do I willingly append an explanation for some of the more confusing terms necessarily employed herein. To those readers not interested in such explanations, 'tis unnecessary to suggest the waste of time in reading these closing pages. ALT: All tones in the first octave above the

staff. ALTISSIMO: All tones above alt. METEORIC CONDITIONS: Are such condi- tions of the air as barometic pressure, temperature, humidity, (water vapor,) winds, and clouds. As all of these items, singly, or combined, operate to di- minish or augment the distance traveled by sound- waves, therefor, such items become matters of in- terest to the student of violin tone. BAROMETRIC PRESSURE: Refers to weight of the air as indicated by action of mercury in the tube of the barometer. As the mercury falls, weight of air is lighter. As the mercury rises,

289

\section*{VIOLIN TONE-PECULIARITIES.}
weight of air is heavier. Such variations in the weight of air are frequent, especially in our sum- mer months, and, they may be so great as to cause great difference in the distance traveled by any sound from any source whatever. Strangely, the violin, of all musical devices, is the most sensitive to atmospheric conditions as is manifested by feebleness of tone-inten sity upon dates when heat and water vapor are in excess. Under such con- ditions, string-tone values suffer unavoidable de- preciation, no matter who the maker, nor what the quality of material. THE NIMBUS CLOUD: Overspreading all of the sky, operates to augment distance traveled by sound-waves. TEMPERATURE: May be either so high, or low as to greatly diminish the distance traveled by sound-waves. WINDS: Operate to diminish the distance trav- eled by sound-waves in opposition, to the current, and, to augment the distance traveled with the current. HOUR OF DAY: Is of importance in the record for the distance traveled by sound-waves, because during the midday hours, sound is propagated with greatest difficulty upon any given day. SOUND-REFLECTING AGENTS: Are such objects as buildings, hills, timber, operating to aug- ment distance traveled by violin tone in the long- distance out-of-doors test for intensity. Thus, the record of any violin for "carrying power "possesses but little interest unless accompa-

290

\section*{VIOLIN TONE-PECULIARITIES.}
nied by the readings of the barometer, hygrometer, thermometer, together with the direction and ve- locity of winds, degree of cloudiness, hour of day, and, the presence or absence of sound-reflecting agents. With such record thus complete, a violin may be safely .guaranteed to repeat its perform- ance under similar meteoric conditions. It is im- portant to the violin player to know the distance to which the tone of his violin travels with ease. Thus, the player may avoid overexertion with the bow, as such overexertion always produces a dis- agreeable effect upon the musically cultivated ear. OSCILLATION, AMPLITUDE of OSCILLA- TION, VIBRATION, both NORMAL and TAN- GENTIAL, NODES, VENTRAL SEGMENTS, HARMONIC OVERTONES, and DISSONANT OVERTONES: Are explained thus: Confining this matter to facts of practical value to violin tone, I employ these two strings as a means for assistance in making definitions clear to the understanding of every reader. These strings are of equal length, equal diameter, and, of identical material, but, are A unequal in structural perfection. String is ap- parently perfect, String B is apparently imper- fect. The action of these strings will explain, not only some technical terms applied to vibrating bod- ies, but, will explain the necessity for employing wood of structural perfection to produce the violin of best, or highest tone- values; also explain why the violin maker, scientific, or otherwise, must meet defeat when working with wood of structur- al imperfections-

291

\section*{VIOLIN TONE-PECULIARITIES.}
A motionless body is said to be at the point of rest. When a body moves to a certain distance from the point of rest, thence returns to the point of rest, thence passes to an equal distance in the op- posite direction, thence returns a second time to the point of rest, such body is said to be in vibrat- ion, or, in oscillation, whichever term is peferred. The distance traveled by such body from its point of rest is the amplitude of its oscillation. [It is of interest to note the fact that English and German philosophers differ with: French phil- osophers as to what movement of a vibrating body completes one vibration. The former hold that one complete vibration consists in one complete movement each way from the point of rest; where- as, the latter hold that one movement from the point of rest with one return completes one vibra- tion. Thus, from the French method of express- ion, concert pitch is given as A equals 900 vibrat- ions per second; whereas, from the English and German method, concert pitch is given as A equals 450 vibrations per second.] Attaching these strings separately to immovable blocks, and to separate pegs, equal tension is ap- plied. Application of a violin bow causes string A to wind rapidly around its long axis; and,, such winding continues in one direction until elastic en- ergy in string-fiber exceeds bow-friction; where- upon, the string unwinds itself, only to be instant- ly wound up again. Such rapid winding and un- winding delivers forceful blows upon contiguous air molecules, and, as such molecules are elastic

292

\section*{VIOLIN TONE-PECULIARITIES.}
spheres, and, as such spheres touch each other, therefore, the blows from the string arouse sound- wave movement; and; as such movement reaches our tympani, (ear drums, ) we become conscious of a tone proceeding from the string; and, as such tone is produced by action of the entire length of the string, therefore this tone is of the lowest pos- sible pitch for a string of this length, and, with this tension; therefore this tone is the fundament- al tone of string A. [Although this tone is fundamental to string A, yet, we may not suppose that no other fundament- al tone can be produced upon string A, because, when shortening the string by "stopping" it, as with pressure of the finger, a new tone of higher pitch is produced; and such new tone may be rightfully called the fundamental tone of a new key; and, the same holds true of all other tones of higher pitch.] The action of string A, at once causes the ap- pearance of absorbing phenomena; one of which consists in the string describing a circle around its long axis. We observe that such circle is smallest at the ends of the string, and, greatest at the half- way point. Without any reasoning whatever, it is clear that the greatest force in the blows of this string are delivered at the point of its greatest amplitude of oscillation; and, that the least force in its blows are delivered at the points of least amplitude of oscillation. It is also clear that the point of greatest amplitude of oscillation cannot occur at any other than the half-way point.

293

\section*{VIOLIN TONE-PECULIARITIES.}
[Acting upon this fact, when graduating the sounding-board for maximum tone-power, thick- ness, from bridge-position and from the ends of the plate, slightly diminishes to the half-way point; and, for the purpose of increasing the am- plitude of oscillation at such point. From the action of the string, it is clear that equal amplitude of oscillation cannot be secured at any other point between bridge-position and the ends of the plate. Another phenomena, observed upon string A, di- rects attention to nodes, ventral segments and harmonic overtones. Lightly stopping the string at the half-way point, in either direction therefrom are seen several points where the string is at rest. These points of rest are called nodes. They are equally distant from each other, and the space between two nodes is called a ventral seg- ment; and, we observe that each ventral segment acts precisely as the whole length of the string acts; that is, the widest amplitude of ventral seg- ment oscillation is at its half-way point; and, that the action of these ventral segments produces a musical tone. As the ventral segments are of sim- ilar length, therefore the tone of each segment is of the same pitch as the tone of its neighbor: and, as the pitch of such tone is in harmony with the fundamental tone, therefore such tone is called a harmonic overtone, and joining with the fundament- al tone, produces what is called "rich" tone a rare quality of tone a quality of tone highly valued by the violin soloist a quality of tone in which no

294

\section*{VIOLIN TONE-PECULIARITIES.}
other musical device approaches the "best" violin a quality of tone impossible to the sounding- board of structural imperfections; and, the reason will be demonstrated by the action of structurally imperfect string B. [Applied to music, the terms "consonant," and, "dissonant" vary in meaning as widely as "saint" and "devil" applied to humanity. Consonant har- monic overtones, and, harmonics a basso, are the causes for enchanting beauty of violin tone. They lie between such terms as "sweetness," and, "richness." To them we are indebted for the crown which we so enthusiastically offer to "The King." But, we are reminded, very of ten remind- ed, that there are other crowns. Satan wears one. Strangely, Satan's crown may come from the vio- lin; and yet more strange, Satan's activity in cornering material for his crown is supplemented by such SCIENTIFIC violin makers who ignore structural perfection in sounding-board wood.] Applying the bow to string B, the dissonant overtone, (chief gem in Satan's crown,) stands out clearly for perception by the sense of hearing; and, the cause stands out clearly for perception by the sense of sight. We observe that the nodes upon this structurally imperfect tone-producing agent exist at varying intervals; therefore the ven- tral segments are of varying lengths; therefore the tones from these ventral segments are of random pitch; therefore many of such tones are not in harmony with the fundamental tone of the string; therefore the sound from this structurally imper-

295

\section*{VIOLIN TONE-PECULIARITIES.}
feet tone-producing agent is NOISE; and, for the same reason, the tone from the structurally imper- fect sounding-board is NOISE; therefore the letter- ing on Satan's crown is N-0-I-S-E. In text books devoted to the philosophy involved in musical sound, the violin sounding-board is com- pared with a "bundle of strings;" and, determin- ing by the "rich" tone from the perfect string, and, by the "rich" tone from the perfect sounding- board, it is made clear that part of the vibratory action of the perfect string is duplicated by the perfect sounding-board fiber; also, that part of the action of the imperfect string is duplicated by the action of the imperfect sounding-board fiber. Thus, in the sounding-board of perfect fiber, nodes will occur at regular intervals; therefore ventral seg- ments will be of uniform lengths; therefore, tones from ventral segments will be of uniform pitch; therefore, such tones are in harmony with the fundamental tone; hence, the "rich" violin tone. Thus, in the sounding-board of imperfect fiber, as knots, curls, unequal density, departure of fiber from straight lines, great density and inelasticity of connective tissue, sap wood and black spots, the nodes must occur at irregular intervals; therefore, ventral segments must be of varying lengths; therefore, the tones from ventral segments must be of random pitch therefore, many of such tones must be out of harmony ; with the fundamental tone; hence, the noisy violin which defeats sci- ence; the violin which defeats the most skillful vi- olin maker; and, when noise is therein suppressed,

296

\section*{VIOLIN TONE-PECULIARITIES.}
'tis the violin of "cold" tone; and the violin of "cold" tone is no longer "The King." VIBRATORY MOVEMENT: In the violin sounding-board command deep interest from the student; and, such interest 'is due to the fact that, with other factors at the best, violin tone-values depend upon both normal vibratory movement, which travels along the fibers, and, tangential vi- bratory movement, which travels across the fibers. These two vibratory movements are the very foundation of violin tone-values. They are the factors which augment the tone of the strings. Without these factors,, violin tone- values depreciate to the level of the "broom-stick fiddle." By the presence of dry sand upon the flat sounding-board, both normal and tangential vibratory movements become defined. Uniformly distributed upon such sounding-board, the sand is thrown upward by the force of plate-oscillation; and, is thrown up to a height proportionate to the amplitude of such os- cillation. Right here is where we observe astound- ing variations in the spring-action of different samples of wood. Right here also is shown the vast difference between the action of the structur- ally perfect and imperfect fiber. Thus: Upon the perfect sounding-board, oscillation thereof, and continued during a period of time varying as the spring-action, forces the sand to leave ventral seg- ments, and, to collect upon the nodes. Thus, reg- ularity in the distance between nodes is determ- ined. Thus, irregularity in the distance between nodes is shown .on the structurally imperfect

297

\section*{VIOLIN TONE-PECULIARITIES.}
sounding-board. 'Tis true, this demonstration cannot be made successfully upon the concavo-convex violin sound- ing-board. 'Tis also true, such demonstration is unnecessary. It is self-evident that these facts are similar in these two forms of sounding-board. In either case, the presence of the exits operates to limit transverse vibratory travel, thus: As the blows of the strings are delivered upon the sound- ing-board at bridge-position, therefore, vibratory movement in the sounding-board begins at such position; and, because of proximity of the exits, it is evident that vibratory movement above the bridge is confined to the normal until after passing beyond the upper extremity of the exits. Below the bridge, on the bass side, the distance to the exit is sufficient for permitting simultaneous action of both normal and transverse vibration. Below the bridge, on the treble side, the post effectively arrests the travel of both vibratory movements, as is satisfactorily demonstrated by splitting up the lower right-quarter of the sounding-board. Bas- ing conclusions upon the evidence afforded by hundreds of used sounding-boards, it is evident that the "rich" tone depends upon the unimpeded travel of normal vibration; also, that volume of tone largely depends upon transverse vibration. An interesting fact is shown in the record for comparative velocities attaching to these vibratory movements. Thus: In pine, normal vibrations, 3322 metres per second. In pine, transverse vibrations, 1405 metres per second. Therefore,

298

\section*{VIOLIN TONE-PECULIARITIES.}
as transverse vibration travels one unit, nor- mal vibration travels 2 36-100 units; but, the evidence from used sounding-boards indicates that this ratio is not a constant quantity, that it is a quantity greatly modified by structural peculiar- ities of grain in different samples of wood. Thus, transverse vibration is dimished by both extremely soft grain, and extremely dense grain. United, these vibratory movements become the foundation of both volume of tone, and quality of tone. In the production of volume of tone from the sounding-board, transverse vibration exercises the more commanding influence. This fact is eas- ily demonstrated, thus: With that sounding-board too rigid for the force in strings of certain size, as gauge-2, gradual diminution of thickness from center join to edges operates to increase the distance traveled by transverse vibration; therefore, the area of the striking surface is increased; therefore, an increased number of air molecules within the violin body receive an identical blow; hence greater volume of tone. [It is an easy matter to increase volume of violin tone to the degree ruining intensity of tone.] The loss of power following use is a matter pos- sessing deep interest to the student of violin tone- phenomena. It is my observation that such loss is due to the following factors: 1. Increasing roughness of unprotected interior surfaces. 2. Disintegration upon unprotected surface. 3. Increasing sounding-board flexibility follow-

299

\section*{VIOLIN TONE-PECULIARITIES.}
ing use. 4. Degree of fiber-density and natural tough- ness of connective tissue. 5. Amount of use and vigor of bow. Increasing flexibility is here in point. By reason of inherent rigidity, the hard parts of sounding- board fiber are prevented from transverse bending; but, connective tissue, being elastic in both length and breadth, permits transverse bending; there- fore transverse vibration depends upon connective- tissue elasticity; and such elasticity diminishes with use; and the rate of diminution is modified by the degree of bending, frequency of bending, and the natural toughness of connective tissue peculiar to different samples of wood. The toughness of connective tissue in violin sounding-board wood is a widely varying quantity; and for this reason alone, longevity of violin tone-power presents widely varying periods. These known facts may' be applied to the violin sounding-board, thus: That sounding-board called upon for only part of its spring-force will last longer than that sounding- board called upon for the limit of its spring-force. Both violin plates may become springs. This matter depends upon thickness and inherent rigid- ity. Manifestly, marked volume of violin tone de- pends upon such reduction in plate rigidity as per- mits increased distance traveled by, and increased amplitude to transverse vibration; also manifestly, such increased amplitude is accomplished by in- creased bending of the plate; therefore, loss of spring-force is accelerated.

300

\section*{VIOLIN TONE-PECULIARITIES.}
These facts unmistakably point to the risk in giving marked volume of tone to the new violin, because, such new violin is liable to suffer damag- ing loss of tone-power at what may be rightfully called premature age. From my viewpoint, it is much the wiser plan to slightly limit both distance and amplitude of transverse vibration and trust something to increasing flexibility of the wood in- evitably following use. [The work of Prof. Pietro Blaserna, Royal Uni- versity, Rome, is an admirable treatise for persons desiring study of sound in its relation to music.] POSITION OF THE POST: As the problem of the post has received various interpretations by different philosophers, and as no authoritative ground to stand upon appears in sight, therefore all persons are at liberty to hold and express opin- ions upon this question as suits themselves. Upon this basis, the following calculations are presented: As violin string-force depends upon augmenta- tion by normal and transverse vibration in the sounding-board, plus sympathetic action due to proximity of the strings, and, as placing the post below the bridge operates to confine said activities to the upper, right quarter, as is demonstrated by splitting up the lower, right quarter, therefore, position of the post, in its relation to the bridge, governs the location of sounding-board activities augmenting force of the A and E strings; and, as placing the post above the bridge operates to trans- fer normal and transverse vibration to the lower, right quarter of the sounding-board, and, as dist-

301

\section*{VIOLIN TONE-PECULIARITIES.}
ance annihilates sympathetic action, therefore placing the post above the bridge operates to de- liver blows of diminished force upon contained air; hence diminished power of A and E-tone follows. As the influence of the post is chiefly exerted upon the A and E-strings, therefore, in selecting a posi- tion for the post, certain effects should be kept in view, as, greatest power of tone, and, best quality of tone. It happens quite often that greatest power of tone destroys best quality of tone; there- fore, in such case, a choice between greatest power and best quality of tone must be determined; and, in determining upon such choice, it is necessary to consider the post itself regardless of its position. This consideration is made necessary because such physical qualities in the post as length, dens- ity, and mass exert marked influence upon A and E-tone. Either of these attributes may defeat in- tention. Manifestly, preparatory training here be- comes of benefit, in fact, a necessity for securing best results in the work of post adjustment. As in all other matters pertaining to the modification of violin tone, the musical sense of the post-setter must be thrown into the balance; and, when such sense is a minus quantity, the failures following at- tempts at post-adjusting should not be charged up to Stradivarius. In addition to the physical qualities of the post, its position is modified by each of the following factors: 1. Height of bridge. 2. Height of arching.

302

\section*{VIOLIN TONE-PECULIARITIES.}
3. Thickness of plates at bridge-position. 4. Method of graduation. 5. Spring-quality of the plates. 6. Diameter and quality of string. Obviously, that position for the post securing greatest augmentation of A and E-tone can be de- termined only by trial. From differing methods for setting the post, the following is presented: Premising that length of the plates equals 14 inches, the bridge is first placed 8 inches from the upper end of the sounding-board for reasons here- tofore given. [This position for the bridge is modified by the method of graduation.] Applying sufficient tension to the strings to hold the bridge in position, set the post directly beneath the right foot of the bridge, and draw the strings up to either diapason normal, (international pitch, ) or up to concert pitch, thereafter maintaining the pitch and the gauge of strings decided upon. With the post in this position, observe the power of A and E-tone; then move the post downwards by 1-16 inch, and observe the increased power of A and E, and, continue thus until reaching the point where power of A and E. appears diminished. In returning the post to the point most augmenting power of A and E, move it by 1-32 inch. Upon reaching such point, observe the balance of power between A and E. , Should the E possess the greater power, move the post to the left by 1-32 inch, and observe the result; continuing thus

303

\section*{VIOLIN TONE-PECULIARITIES.}
until equal power is established. Should the A possess the greater power, move the post to the right. In this work, defeat may follow from the follow- ing fact: With the diameter of the post equalling 3-16 inches, but slight deviation from the perpen- dicular causes the sounding-board to rest upon one or the other edge of the post. Thus, when the sounding-board rests upon the lower edge of the post, the distance from the bridge is practically 3-16 greater than appearances indicate; and as 1-32, more or less, causes perceptible effect upon A and E-string tone, therefore 6-32 becomes ample reason for defeat. ANGLES of INCIDENCE and REFLECTION, within the violin body, are important factors in the production of tone-intensity. The influence of these angles is easily comprehended; yet, in no form of arching whatever have these lines been definitely traced. These angles depend for existence upon a moving body and a solid reflecting surface, thus: The ball, thrown against the wall, travels along the line of incidence; and, as the ball rebounds, it travels along the line of reflection. If the line of inci- dence is perpendicular to the wall, then the ball re- bounds directly to its starting point; thus the lines of incidence and reflection are one; and thus are these lines in the violin of perfectly flat plates. If now, the thrower stands to one side or the other of the line which is perpendicular with the wall, and plants the ball at a point where the perpendic-

304

\section*{VIOLIN TONE-PECULIARITIES.}
ular line touches the wall, then the ball will not rebound on the line of incidence nor upon the per- pendicular line, but will travel along a new line having precisely the same angle with the wall as the line of incidence; and thus are lines of sound- wave travel within the violin of arched plates. It is self-evident that such interior walls as turn sound-wave lines of travel away from the exits cause loss to intensity of tone; and per contra, such interior walls as direct sound-wave lines of travel toward the exits cause increased intensity of tone proportionate to the degree of concentration. Verbum sap First, establish the interior walls; later, the exterior walls. LABEL, VARNISH and PRICE vs SWEET TONE: Johann Holtzhammer is a herder of sheep. His summer home is on the mountain. From early morn till dewy eve, Johann listens to sweet sounds. The song of birds is sweet. The song of the rippling brook is sweet. The rustle of swaying boughs is sweet. The odor of pine is sweet. Even the noisy whirr of industry's wheels is bereft of the noise- wave ere it reaches up to Jo- hann's ears. At the base of detached rock, wild berries offer their sweet appearances and sweeter juices. The mountain air is sweet. The water in the rivulet is sweet. Indeed, life to Jo- hann is one round of sweetness; yet, Johann longs for something to keep his hands employed. With NATURE throbbing all around him, idleness be- comes oppressive. What would he do? What could he do? Make a violin? Why not? Here is

305

\section*{VIOLIN TONE-PECULIARITIES.}
a fallen, giant pine. Yonder is a dead maple. There's the ax. Here's a jackknife. Where there's a will, there's a way. Those giant logs have been seasoning since Johann's grandsire herded the sheep; but, the ax is keen, and Jo- hann's jackknife is always ready for business. Time is nothing. Result is everything-. As Johann works, a nearby lark sings with undoubted approv- al. 'Tis true, Johann knows nothing of the Py- thagorean scale; nothing of the Palestrina scale; nothing of the piano, or temperate scale; nothing- of consonant overtones; nothing of harmonics a bassa; but, he does know sweet sound. 'Tis all he needs. Johann stains his violin with the juice of moun- tain berries.

Enough! Imitation is defied. The secret dies with Johann. The thrilling, enchanting, soulful, sweetly tend- er tone of Johann's violin had long been the cause for bird wonderment and worship prior to his passing; and, after his passing, each bird on that mountain vied with its neighbor in singing sweet and low upon the recurring date when Johann de- parted with his violin. Not one bird on that moun- tain ever mentioned the jackknifed angles upon Johann's violin; neither did Charon as Johann ap- plied for ferriage across the Styx; neither did the keeper of the great gate, as appears in the sequel. As Johann entered the long avenue, a multitude of shades with downcast mein, were standing at a

306

VIOLIN TONE-PECULIARITIES. halt upon either hand. Some of them bore ' 'mast- erpieces" upon which were conspicuous tickets lettered with the words, LABEL, and, VARNISH, and, a limited number bore conspicuous figures up- on the ticket, as, \$10,000, and, \$12,000. Although no affair of Johann's, yet, 'twas a wonderment whyfor such halting. As this gate keeper requires only the billionth part of a moment for computing the ad valorem in any invoice, there- fore, upon the instant Johann's immortal consign- ment arrived, the great gate swung open; and, in that brief interval, a glimpse of glory flashed down the long avenue; and, as the great gate closed be- hind Johann, there came up from the long avenue a sound as the gnashing of teeth.

307

\$

\section*{Univ Calif - Digitized by Microsoft ®}
\section*{Univ Calif - Digitized by Microsoft ®}
\section*{Univ Calif - Digitized by Microsoft ®}
\section*{UNIVERSITY OF CALIFORNIA LIBRARY Los Angeles This book is DUE on the last date stamped below. REC'D MUS-UB I}
OCT 21 REC'D MUS'LIB 1 NOV 1 J977 OCT 25 1977 MAR 31 1975 IWR

\section*{NOV 11982 SEP 29 1975 REC'D MUS-L1B AUG 3 1 1982 4 1975}
REC'D MUS-LW FEB 1 FEB 1 9 Form L9-Series 4939

\section*{UC SOUTHERN REGIONAL LIBRARY FACILITY i mil mil linn 111 mil ill! 111 MUSIC LI3RARY A 000754194 9 ML}
C2?9v

\end{document}